% Lesson 8 — Grammar: Revise question tags — Anglais Tle A — Cours signé OnBuch+
% Programme officiel MINESEC. Compiler avec : tectonic EN08-grammar-revise-question-tags.tex
\newcommand{\DOCMATIERE}{Anglais}
\newcommand{\DOCNIVEAU}{Tle A}
\newcommand{\DOCMODULE}{Chapter 1 : Family and Social Life}
\newcommand{\DOCLECON}{EN08}
\newcommand{\DOCTITRE}{Lesson 8 — Grammar: Revise question tags}
% OnBuch+ — préambule « cours signés » ENRICHI (compilé avec tectonic / XeLaTeX)
% Système visuel « Pop » d'OnBuch+ : fond crème (jamais blanc), bordures encre
% épaisses, ombres « dures » décalées, titres Archivo Black en capitales.
% Version MATHÉMATIQUES : graphes (pgfplots/tikz), boîtes pédagogiques
% (définition, propriété, méthode, attention, le savais-tu, exercice résolu,
% exercice, TP, bilan), page de garde et signature OnBuch+ ancrée en bas.
%
% Macros attendues dans main.tex AVANT \input{preamble} :
%   \DOCMATIERE  \DOCNIVEAU  \DOCMODULE  \DOCLECON (numéro)  \DOCTITRE
\documentclass[11pt]{article}

\usepackage{fontspec}
\usepackage[english,french]{babel} % français = langue principale ; l'anglais via \eng{..} / textebox
\usepackage[a4paper,margin=2cm,top=3.4cm,bottom=2.8cm,headheight=1.4cm,footskip=1.3cm]{geometry}
\usepackage{amsmath,amssymb,amsfonts}
\usepackage{mathtools} % \xmapsto, \coloneqq... souvent utilisés par les modèles
\usepackage{cancel} % simplifications barrées (bilans de forces)
% \abs{z} (module, valeur absolue) et \norm{u} (norme), utilisés par les modèles
\providecommand{\abs}{}\let\abs\relax\DeclarePairedDelimiter{\abs}{\lvert}{\rvert}
\providecommand{\norm}{}\let\norm\relax\DeclarePairedDelimiter{\norm}{\lVert}{\rVert}
\usepackage[table]{xcolor} % [table] : fournit aussi \rowcolors (alternance de lignes)
\usepackage{pagecolor}
\usepackage{tikz}
\usetikzlibrary{arrows.meta,positioning,calc,shapes.geometric,shapes.misc,shapes.arrows,decorations.pathmorphing,decorations.pathreplacing,decorations.markings,patterns,shadows,mindmap,trees,fit,backgrounds,matrix,intersections,angles,quotes}
\usepackage{pgfplots}
\usepackage[version=4]{mhchem} % équations nucléaires \ce{^{235}_{92}U}
\usepackage[european,siunitx]{circuitikz} % schémas électriques
\pgfplotsset{compat=1.18}
\usepgfplotslibrary{fillbetween,groupplots} % aires entre courbes (calcul intégral)
\usepackage{enumitem}
\usepackage{fancyhdr}
% [most] charge toutes les bibliothèques tcolorbox (listings, minted, raster,
% documentation...) et gonfle inutilement la mémoire TeX sur un document long
% (~300 boîtes) : on ne charge que ce qui est réellement utilisé.
\usepackage[skins,breakable]{tcolorbox}
\usepackage{graphicx}
\usepackage{colortbl}
\usepackage{booktabs}
\usepackage{tabularx}
\usepackage{diagbox} % case d'en-tête diagonale des tableaux à double entrée
% Colonnes extensibles centrée / gauche / droite (C, L, R), souvent utilisées par les modèles
\newcolumntype{C}{>{\centering\arraybackslash}X}
\newcolumntype{L}{>{\raggedright\arraybackslash}X}
\newcolumntype{R}{>{\raggedleft\arraybackslash}X}
\usepackage{array}
\usepackage{multicol}
\usepackage{multirow}
\usepackage{siunitx}
% parse-numbers=false : accepte aussi \SI{2,5 \times 10^{-3}}{...}
\sisetup{locale=FR,output-decimal-marker={,},group-minimum-digits=4,parse-numbers=false}
\DeclareSIUnit\ohm{\ensuremath{\symup{\Omega}}} % U+2126 absent de la police
\DeclareSIUnit\division{div} % réglages d'oscilloscope (ms/div, V/div)
\DeclareSIUnit\tr{tr} % tours (vitesse de rotation, ex. \SI{1500}{\tr\per\minute})
\DeclareSIUnit\molar{M} % concentration molaire (ex. \SI{3}{\molar}, \SI{1}{\micro\molar})
\DeclareSIUnit\an{an} % année (le modèle confond parfois avec \year, un compteur TeX) — ex. \SI{5}{\kilo\meter\per\an}
\DeclareSIUnit\FCFA{FCFA} % franc CFA (ex. \SI{500}{\FCFA\per\kilo\gram})
\DeclareSIUnit\degreeBrix{\textdegree Brix} % teneur en sucre d'un sirop/jus (réfractométrie)
\DeclareSIUnit\month{mois} % le modèle utilise parfois \month, un compteur TeX, comme unité de durée
\usepackage{qrcode}
% \degree : commande fréquemment utilisée par les modèles pour un angle ou une
% température en dehors de \SI{}{} (non fournie par les paquets chargés).
\providecommand{\degree}{^{\circ}}
% \qed (fin de démonstration), utilisé par les modèles sans amsthm
\providecommand{\qed}{\hfill$\square$}
% \barre : double barre des valeurs interdites dans un tableau de variations (tkz-tab)
\providecommand{\barre}{\ensuremath{\|}}
% environnement proof (amsthm non chargé) utilisé par les modèles
\ifdefined\proof\else\newenvironment{proof}[1][Démonstration]{\par\noindent\textit{#1.}\ }{\qed\par}\fi
\usepackage{lastpage}
\usepackage{hyperref}
\hypersetup{hidelinks}

% ============================================================
% Palette « Pop » OnBuch+ — mêmes hex que la classe P (plus_ui.dart)
% ============================================================
\definecolor{poporange}{HTML}{F59321}
\definecolor{popdark}{HTML}{E07A0C}
\definecolor{popcream}{HTML}{FFF6E8}
\definecolor{popink}{HTML}{1C1714}
\definecolor{poppurple}{HTML}{7A5AE0}
\definecolor{popgreen}{HTML}{2E9E6A}
\definecolor{poppink}{HTML}{E0457B}
\definecolor{popblue}{HTML}{2F7DE1}
\definecolor{popgold}{HTML}{C98A12}
\definecolor{popmuted}{HTML}{8A7F76}
\definecolor{popline}{HTML}{EADFD2}
\definecolor{popgray}{HTML}{8A7F76}
\colorlet{popgrey}{popgray}
% Alias tolérants (le modèle nomme parfois les couleurs différemment)
\colorlet{popwhite}{white}
\colorlet{popblack}{popink}
\colorlet{popred}{poppink}
\colorlet{popyellow}{popgold}
% Teintes claires (fonds de tableaux/encadrés) — le modèle nomme parfois
% les couleurs avec un suffixe L (ex. popblueL) sans les avoir définies.
\colorlet{poporangeL}{poporange!20}
\colorlet{popdarkL}{popdark!20}
\colorlet{poppurpleL}{poppurple!20}
\colorlet{popgreenL}{popgreen!20}
\colorlet{poppinkL}{poppink!20}
\colorlet{popblueL}{popblue!20}
\colorlet{popgoldL}{popgold!20}
\colorlet{popredL}{popred!20}
\colorlet{popgrayL}{popgray!20}
% Teintes claires pour fonds de tableaux / zones
\colorlet{popblueL}{popblue!10!white}
\colorlet{poporangeL}{poporange!14!white}
\colorlet{popgreenL}{popgreen!12!white}
\colorlet{poppurpleL}{poppurple!12!white}
\colorlet{poppinkL}{poppink!12!white}
\colorlet{popgoldL}{popgold!14!white}

\pagecolor{popcream}

% ============================================================
% Typographie
% ============================================================
\defaultfontfeatures{Ligatures=TeX}
\setmainfont{PlusJakartaSans-Medium.ttf}[Path=../../../fonts/,BoldFont=PlusJakartaSans-Bold.ttf]
% Symboles, API et emojis absents de la police principale : repli sur DejaVu Sans (caractères actifs)
\newfontface\symfont{DejaVuSans.ttf}[Path=../../../fonts/]
\makeatletter
\newcommand\symactive[1]{\catcode#1=13 \begingroup\lccode`\~=#1 \lowercase{\endgroup\def~}{{\symfont\char#1}}}
\newcommand\symblank[1]{\catcode#1=13 \begingroup\lccode`\~=#1 \lowercase{\endgroup\def~}{}}
\newcommand\symhyph[1]{\catcode#1=13 \begingroup\lccode`\~=#1 \lowercase{\endgroup\def~}{\mbox{-}}}
\newcount\symc
\newcommand\symrange[2]{\symc=#1 \@whilenum\symc<#2 \do{\expandafter\symactive\expandafter{\the\symc}\advance\symc by 1 }}
\newcommand\symblankrange[2]{\symc=#1 \@whilenum\symc<#2 \do{\expandafter\symblank\expandafter{\the\symc}\advance\symc by 1 }}
\makeatother
\symrange{"0250}{"02FF}\symactive{"2713}\symactive{"2714}\symactive{"2717}\symactive{"2718}\symactive{"2610}\symactive{"2611}\symactive{"2612}
\symhyph{"2011}\symhyph{"2E1A}\symblankrange{"1F300}{"1FAFF}\symblank{"FE0F}
% Maths en sans-serif (Fira Math) avec chiffres et lettres droites de Plus
% Jakarta, pour que les formules se fondent dans le texte.
\usepackage[math-style=ISO,bold-style=ISO]{unicode-math}
\setmathfont{FiraMath-Regular.otf}[Path=../../../fonts/]
\setmathfont{PlusJakartaSans-Medium.ttf}[Path=../../../fonts/,range={up/num,up/{Latin,latin},\mathup}]
\usepackage{icomma} % virgule décimale sans espace en mode math
% Caractères mathématiques Unicode absents des polices de texte (Plus Jakarta,
% Archivo Black) : rendus par leur équivalent mathématique, en texte comme en
% formule — sinon ils s'affichent en carré vide (ex. « Arithmétique dans ℤ »).
\usepackage{newunicodechar}
\newunicodechar{ℝ}{\ensuremath{\mathbb{R}}}
\newunicodechar{ℤ}{\ensuremath{\mathbb{Z}}}
\newunicodechar{ℂ}{\ensuremath{\mathbb{C}}}
\newunicodechar{ℕ}{\ensuremath{\mathbb{N}}}
\newunicodechar{ℚ}{\ensuremath{\mathbb{Q}}}
\newunicodechar{𝔻}{\ensuremath{\mathbb{D}}}
\newunicodechar{α}{\ensuremath{\alpha}}
\newunicodechar{β}{\ensuremath{\beta}}
\newunicodechar{γ}{\ensuremath{\gamma}}
\newunicodechar{δ}{\ensuremath{\delta}}
\newunicodechar{ε}{\ensuremath{\varepsilon}}
\newunicodechar{θ}{\ensuremath{\theta}}
\newunicodechar{λ}{\ensuremath{\lambda}}
\newunicodechar{μ}{\ensuremath{\mu}}
\newunicodechar{σ}{\ensuremath{\sigma}}
\newunicodechar{τ}{\ensuremath{\tau}}
\newunicodechar{φ}{\ensuremath{\varphi}}
\newunicodechar{ψ}{\ensuremath{\psi}}
\newunicodechar{ω}{\ensuremath{\omega}}
\newunicodechar{Σ}{\ensuremath{\Sigma}}
% majuscules produites par \MakeUppercase dans les titres (α-aminés -> Α)
\newunicodechar{Α}{\ensuremath{\alpha}}
\newunicodechar{Β}{\ensuremath{\beta}}
\newunicodechar{Γ}{\ensuremath{\Gamma}}
\newunicodechar{Δ}{\ensuremath{\Delta}}
\newunicodechar{Ε}{\ensuremath{\varepsilon}}
\newunicodechar{Θ}{\ensuremath{\Theta}}
\newunicodechar{Λ}{\ensuremath{\Lambda}}
\newunicodechar{Μ}{\ensuremath{\mu}}
\newunicodechar{Τ}{\ensuremath{\tau}}
\newunicodechar{Φ}{\ensuremath{\Phi}}
\newunicodechar{Ψ}{\ensuremath{\Psi}}
\newunicodechar{Ω}{\ensuremath{\Omega}}
\newunicodechar{↦}{\ensuremath{\mapsto}}
\newunicodechar{∈}{\ensuremath{\in}}
\newunicodechar{∉}{\ensuremath{\notin}}
\newunicodechar{∧}{\ensuremath{\wedge}}
\newunicodechar{⇔}{\ensuremath{\Leftrightarrow}}
\newunicodechar{⟺}{\ensuremath{\Longleftrightarrow}}
\newunicodechar{⇒}{\ensuremath{\Rightarrow}}
\newunicodechar{⟹}{\ensuremath{\Longrightarrow}}
\newunicodechar{∘}{\ensuremath{\circ}}
\newunicodechar{∪}{\ensuremath{\cup}}
\newunicodechar{∩}{\ensuremath{\cap}}
\newunicodechar{≡}{\ensuremath{\equiv}}
\newunicodechar{⊂}{\ensuremath{\subset}}
\newunicodechar{⊥}{\ensuremath{\perp}}
\newunicodechar{∃}{\ensuremath{\exists}}
\newunicodechar{∀}{\ensuremath{\forall}}
\newunicodechar{∛}{\ensuremath{\sqrt[3]{\,}}}
\newunicodechar{∜}{\ensuremath{\sqrt[4]{\,}}}
\newunicodechar{⃗}{\ensuremath{{}^{\to}}}
\newunicodechar{ⁿ}{\ifmmode^{n}\else\textsuperscript{\ensuremath{n}}\fi}
\newunicodechar{ˣ}{\ifmmode^{x}\else\textsuperscript{\ensuremath{x}}\fi}
\newunicodechar{ᵇ}{\ifmmode^{b}\else\textsuperscript{\ensuremath{b}}\fi}
\newunicodechar{ʳ}{\ifmmode^{r}\else\textsuperscript{\ensuremath{r}}\fi}
\newunicodechar{ᵘ}{\ifmmode^{u}\else\textsuperscript{\ensuremath{u}}\fi}
\newunicodechar{ʸ}{\ifmmode^{y}\else\textsuperscript{\ensuremath{y}}\fi}
\newunicodechar{ᵃ}{\ifmmode^{a}\else\textsuperscript{\ensuremath{a}}\fi}
\newunicodechar{ᵉ}{\ifmmode^{e}\else\textsuperscript{\ensuremath{e}}\fi}
\newunicodechar{ᵏ}{\ifmmode^{k}\else\textsuperscript{\ensuremath{k}}\fi}
\newunicodechar{ᵖ}{\ifmmode^{p}\else\textsuperscript{\ensuremath{p}}\fi}
\newunicodechar{ᴺ}{\ifmmode^{N}\else\textsuperscript{\ensuremath{N}}\fi}
\newunicodechar{ᶜ}{\ifmmode^{c}\else\textsuperscript{\ensuremath{c}}\fi}
\newunicodechar{ʰ}{\ifmmode^{h}\else\textsuperscript{\ensuremath{h}}\fi}
\newunicodechar{ᵠ}{\ifmmode^{q}\else\textsuperscript{\ensuremath{q}}\fi}
\newunicodechar{ₙ}{\ifmmode_{n}\else\textsubscript{\ensuremath{n}}\fi}
\newunicodechar{ₐ}{\ifmmode_{a}\else\textsubscript{\ensuremath{a}}\fi}
\newunicodechar{ᵢ}{\ifmmode_{i}\else\textsubscript{\ensuremath{i}}\fi}
\newunicodechar{ᵣ}{\ifmmode_{r}\else\textsubscript{\ensuremath{r}}\fi}
\newunicodechar{ₖ}{\ifmmode_{k}\else\textsubscript{\ensuremath{k}}\fi}
\newunicodechar{ᵦ}{\ifmmode_{\beta}\else\textsubscript{\ensuremath{\beta}}\fi}
\newunicodechar{ₓ}{\ifmmode_{x}\else\textsubscript{\ensuremath{x}}\fi}
\newfontfamily\popdisplay{ArchivoBlack-Regular.ttf}[Path=../../../fonts/]
\newfontfamily\poplogo{SpaceGrotesk-Bold.ttf}[Path=../../../fonts/]
\newfontfamily\popsemi{PlusJakartaSans-SemiBold.ttf}[Path=../../../fonts/]
\newfontfamily\popxbold{PlusJakartaSans-ExtraBold.ttf}[Path=../../../fonts/]
\newfontfamily\popxboldls{PlusJakartaSans-ExtraBold.ttf}[Path=../../../fonts/,LetterSpace=3.0]

\setlist[enumerate]{leftmargin=1.7em,itemsep=5pt,topsep=4pt}
\setlist[itemize]{leftmargin=1.7em,itemsep=4pt,topsep=4pt,label={\color{poporange}\textbullet}}
\setlength{\parindent}{0pt}
\setlength{\parskip}{5pt}
\renewcommand{\arraystretch}{1.25}

% pgfplots : style Pop par défaut
\pgfplotsset{
  every axis/.append style={axis line style={popink,line width=1pt},tick style={popink},
    grid style={popline,line width=0.5pt},label style={font=\small},tick label style={font=\footnotesize}},
  popcurve/.style={poporange,line width=1.8pt,no markers,smooth},
}
% Même style disponible dans un \draw TikZ simple (hors pgfplots)
\tikzset{popcurve/.style={poporange,line width=1.8pt}}
\tikzset{popgrid/.style={popline,very thin}}
\colorlet{popcurve}{poporange} % le nom du style est parfois utilisé comme couleur

% ============================================================
% Pill / squiggle
% ============================================================
\newcommand{\poppill}[3]{%
  \tikz[baseline=-0.55ex]\node[draw=popink,line width=1.2pt,rounded corners=2.6pt,%
    fill=#2,inner xsep=6.5pt,inner ysep=3pt]{%
    \popxboldls\fontsize{7}{7}\selectfont\color{#3}\MakeUppercase{#1}};%
}
\newcommand{\popsquiggle}[1]{%
  \tikz[baseline=-0.4ex]{
    \draw[#1,line width=2.6pt,line cap=round]
      plot [smooth] coordinates {(0,0.09) (0.34,0.01) (0.68,0.21) (1.02,0.01) (1.36,0.10)};
  }%
}
% Mot-clé surligné (marqueur orange sous le texte)
\newcommand{\cle}[1]{{\popxbold\color{popdark}#1}}

% ============================================================
% En-tête / pied
% ============================================================
\pagestyle{fancy}
\fancyhf{}
\renewcommand{\headrulewidth}{0pt}
\renewcommand{\footrulewidth}{0pt}
\lhead{\raisebox{-0.15ex}{\poplogo\fontsize{13}{13}\selectfont\color{popink}OnBuch\color{poporange}+}}
\rhead{\poppill{\DOCMATIERE\ \textbullet\ \DOCNIVEAU\ \textbullet\ Leçon \DOCLECON}{white}{popink}}
\lfoot{\popxbold\fontsize{7.6}{7.6}\selectfont\color{popmuted}\MakeUppercase{Cours signé OnBuch+}\ \textbullet\ {\popsemi\DOCTITRE}}
\rfoot{\tikz[baseline=-0.6ex]\node[rounded corners=8.5pt,fill=popink,minimum height=17pt,minimum width=17pt,inner xsep=5pt,inner sep=0]{\popxbold\fontsize{8}{8}\selectfont\color{popcream}\thepage\,/\,\pageref*{LastPage}};}

% ============================================================
% Titres
% ============================================================
\newcommand{\chaptitre}[2]{%
  \begin{tcolorbox}[enhanced,colback=poporange,colframe=popink,boxrule=2.6pt,arc=11pt,
      drop shadow=popink,left=15pt,right=15pt,top=12pt,bottom=11pt,before skip=3pt,after skip=18pt]
    \poppill{#2}{popink}{popcream}\par\vspace{9pt}
    {\raggedright\hyphenpenalty=10000\exhyphenpenalty=10000\popdisplay\fontsize{22}{24}\selectfont\color{popcream}\MakeUppercase{#1}\par}
  \end{tcolorbox}
}
% Section numérotée : pastille orange avec le numéro + titre Archivo
\newcounter{popsec}
\newcommand{\coursec}[1]{%
  \stepcounter{popsec}\setcounter{popsub}{0}%
  \par\vspace{16pt}\noindent
  \begin{minipage}{\linewidth}
  \tikz[baseline=-0.7ex]\node[draw=popink,line width=1.6pt,fill=poporange,rounded corners=4pt,minimum size=22pt,inner sep=0]{\popdisplay\fontsize{12}{12}\selectfont\color{popcream}\thepopsec};%
  \hspace{8pt}{\popdisplay\fontsize{15}{17}\selectfont\color{popink}\MakeUppercase{#1}}\par
  \vspace{4pt}\noindent{\color{poporange}\rule{36pt}{3.6pt}}
  \end{minipage}\par\nopagebreak\vspace{8pt}
}
\newcounter{popsub}
\newcommand{\courssub}[1]{%
  \stepcounter{popsub}%
  \par\vspace{9pt}\noindent{\popxbold\fontsize{11.5}{13}\selectfont\color{popink}{\color{poporange}\thepopsec.\thepopsub}\hspace{6pt}#1}\par\nopagebreak\vspace{4pt}
}

% ============================================================
% Boîtes pédagogiques — même recette : fond blanc, bordure épaisse colorée,
% ombre dure encre, badge plein. Titre optionnel [#1].
% ============================================================
\tcbset{popbase/.style={breakable,enhanced,colback=white,boxrule=2.3pt,arc=9pt,
  shadow={3pt}{-3pt}{0pt}{popink},left=14pt,right=14pt,top=11pt,bottom=11pt,before skip=11pt,after skip=15pt,
  fontupper=\popsemi\fontsize{10.6}{14.5}\selectfont\color{popink},
  fontlower=\popsemi\fontsize{10.6}{14.5}\selectfont\color{popink}}}
% Coche dessinée (le glyphe ✓ n'existe pas dans les polices du document)
\newcommand{\popcheck}[1]{\tikz[baseline=-0.5ex]\draw[#1,line width=1.6pt,line cap=round,line join=round](-0.13,0.0)--(-0.04,-0.09)--(0.13,0.1);}
\providecommand{\coche}{\popcheck{popgreen}}
\providecommand{\resultat}[1]{\par\smallskip\noindent\fcolorbox{popgreen}{popgreenL}{\parbox{\dimexpr\linewidth-2\fboxsep-2\fboxrule\relax}{\textbf{Résultat.} #1}}\par\smallskip}
\newcommand{\popboxhead}[3]{\poppill{#1}{#2}{white}\ifx\relax#3\relax\else\hspace{6pt}{\popxbold\fontsize{10.6}{12}\selectfont\color{popink}#3}\fi\par\vspace{7pt}}

\newtcolorbox{aretenir}[1][]{popbase,colframe=popgreen,before upper={\popboxhead{À retenir}{popgreen}{#1}}}
% Texte en anglais (document de lecture, dialogue, extrait) : cadre
% d'étude. Un titre optionnel (source, auteur) s'affiche dans l'en-tête.
\newtcolorbox{textebox}[1][]{popbase,colframe=popblue,before upper={\popboxhead{Text}{popblue}{#1}\begin{otherlanguage}{english}},after upper={\end{otherlanguage}}}
% Marques ✓ / ✗ (forme correcte / fautive) souvent utilisées par les modèles
\providecommand{\cross}{\textcolor{poppink}{\ensuremath{\times}}}
\providecommand{\xmark}{\cross}\providecommand{\crossmark}{\cross}\providecommand{\cmark}{\ensuremath{\checkmark}}
% Ligne de réponse / justification à compléter (inventée par les modèles)
\providecommand{\justification}[1]{\par\noindent\textit{Justification~:} \dotfill\par}
% \textipa (package tipa non chargé) : la police ne couvre pas l'API -> simple passage
\providecommand{\textipa}[1]{#1}
% Soulignement (ulem non chargé) : \uline, \ul
\providecommand{\uline}[1]{\underline{#1}}\providecommand{\ul}[1]{\underline{#1}}
% \glossaire{\item ...} inventée par les modèles : liste de glossaire
\providecommand{\glossaire}[1]{\begin{itemize}#1\end{itemize}}
% \alert (beamer) inventée par les modèles : texte mis en évidence
\providecommand{\alert}[1]{\textbf{\textcolor{poppink}{#1}}}
% variantes ASCII d'ordinaux écrites en macro
\providecommand{\iere}{\textsuperscript{re}}\providecommand{\ieme}{\textsuperscript{e}}\providecommand{\ere}{\textsuperscript{re}}\providecommand{\eme}{\textsuperscript{e}}\providecommand{\er}{\textsuperscript{er}}\providecommand{\ie}{\textsuperscript{e}}
% Ordinaux français écrits en macro par les modèles : 1\ière, 3\ième
\newcommand{\ière}{\textsuperscript{re}}\newcommand{\ième}{\textsuperscript{e}}
% \exemple (inventée par les modèles) : étiquette « Exemple : »
\providecommand{\exemple}{\textbf{Exemple~:}\ }
% \square utilisable aussi hors mode math (cases à cocher)
\AtBeginDocument{\let\squareMath\square\renewcommand{\square}{\ensuremath{\squareMath}}}
% Mot ou phrase en anglais à l'intérieur d'un paragraphe français
\newcommand{\eng}[1]{\ifmmode\text{\textit{#1}}\else\begingroup\selectlanguage{english}\itshape #1\endgroup\fi}
\let\esp\eng % alias toléré
% flèches d'intonation inventées par les modèles
\providecommand{\textsearrow}{}\renewcommand{\textsearrow}{\ensuremath{\searrow}}\providecommand{\textnearrow}{}\renewcommand{\textnearrow}{\ensuremath{\nearrow}}
% \fr{..} : mot français (inventé par les modèles)
\providecommand{\fr}[1]{\textit{#1}}
\newtcolorbox{exemplebox}[1][]{popbase,colframe=poporange,before upper={\popboxhead{Exemple}{poporange}{#1}}}
\newtcolorbox{definition}[1][]{popbase,colframe=popblue,before upper={\popboxhead{Définition}{popblue}{#1}}}
\newtcolorbox{propriete}[1][]{popbase,colframe=poppurple,before upper={\popboxhead{Propriété}{poppurple}{#1}}}
\newtcolorbox{methode}[1][]{popbase,colframe=popgold,before upper={\popboxhead{Méthode}{popgold}{#1}}}
\newtcolorbox{attention}[1][]{popbase,colframe=poppink,before upper={\popboxhead{Attention}{poppink}{#1}}}
\newtcolorbox{experience}[1][]{popbase,colframe=popblue,colback=popblueL,before upper={\popboxhead{Expérience}{popblue}{#1}}}
% « Le savais-tu ? » : carte violette pleine (contexte, culture, Cameroun)
\newtcolorbox{savaistu}[1][]{popbase,colback=poppurple,colframe=popink,
  fontupper=\popsemi\fontsize{10.4}{14.2}\selectfont\color{white},
  before upper={\poppill{Le savais-tu\,?}{white}{poppurple}\ifx\relax#1\relax\else\hspace{6pt}{\popxbold\color{white}#1}\fi\par\vspace{7pt}}}
% Exercice résolu : énoncé en haut, \tcblower puis la solution sur fond crème
\newcounter{popexo}\newcounter{popexr}
\newtcolorbox{exoresolu}[1][]{popbase,colframe=poporange,colbacklower=popcream,
  segmentation style={popink,line width=1.2pt,dashed},
  before upper={\stepcounter{popexr}\popboxhead{Exercice résolu \thepopexr}{poporange}{#1}},
  before lower={\poppill{Solution}{popink}{popcream}\par\vspace{6pt}}}
% Exercice (non résolu en place) — la correction est dans la partie Corrigés
\newtcolorbox{exercice}[1][]{popbase,colframe=popink,
  before upper={\stepcounter{popexo}\popboxhead{Exercice \thepopexo}{popink}{#1}}}
% Corrigé
\newtcolorbox{corrige}[1][]{popbase,colframe=popgreen,colback=popgreenL,
  before upper={\popboxhead{Corrigé}{popgreen}{#1}}}
% Objectifs / prérequis / bilan
\newtcolorbox{objectifs}{popbase,colframe=popink,colback=poporangeL,
  before upper={\popboxhead{Objectifs de la leçon}{popink}{}}}
\newtcolorbox{prerequis}{popbase,colframe=popink,colback=popblueL,
  before upper={\popboxhead{Prérequis}{popblue}{}}}
\newtcolorbox{bilan}{popbase,colframe=popink,colback=white,boxrule=2.8pt,
  before upper={\popboxhead{Fiche bilan}{poporange}{L'essentiel à connaître pour l'examen}}}
% Figure encadrée avec légende
\newtcolorbox{popfigure}[1][]{enhanced,breakable=false,colback=white,colframe=popline,boxrule=1.4pt,arc=8pt,
  left=10pt,right=10pt,top=10pt,bottom=8pt,before skip=10pt,after skip=12pt,halign=center,
  lower separated=false,halign lower=center,
  fontlower=\popsemi\fontsize{9}{11.5}\selectfont\color{popmuted},
  #1}
\newcommand{\legende}[1]{\par\vspace{6pt}{\popsemi\fontsize{9}{11.5}\selectfont\color{popmuted}#1}}

% ============================================================
% Signature OnBuch+
% ============================================================
\newcommand{\poplogoinline}[1]{{\poplogo\fontsize{#1}{#1}\selectfont\color{popink}OnBuch\color{poporange}+}}
\newcommand{\signatureonbuch}{%
  \begin{tcolorbox}[enhanced,colback=popink,colframe=popink,boxrule=2.2pt,arc=10pt,
      drop shadow=poporange,left=14pt,right=14pt,top=10pt,bottom=10pt,before skip=0pt,after skip=0pt]
    \tikz[baseline=-0.7ex]\node[circle,fill=popgreen,draw=popcream,line width=1.3pt,minimum size=22pt,inner sep=0]{\popcheck{white}};%
    \hspace{9pt}\begin{minipage}[c]{0.62\linewidth}
      {\popxbold\fontsize{12}{13}\selectfont\color{popcream}Signé\ \textbullet\ {\poplogo OnBuch\color{poporange}+}}\par\vspace{2pt}
      {\raggedright\popsemi\fontsize{8}{10.5}\selectfont\color{popline}\MakeUppercase{Équipe pédagogique OnBuch+}\\\MakeUppercase{Conforme au programme officiel MINESEC}\par}
    \end{minipage}\hfill
    {\poplogo\fontsize{22}{22}\selectfont\color{popcream}OnBuch\color{poporange}+}
  \end{tcolorbox}%
}

% ============================================================
% Page de garde
% #1 titre · #2 matière · #3 niveau · #4 module · #5 n° leçon · #6 durée ·
% #7 description / objectifs (texte ou itemize)
% ============================================================
\newcommand{\pagegarde}[7]{%
  \thispagestyle{empty}
  \newgeometry{a4paper,margin=1.7cm,top=1.6cm,bottom=1.4cm}%
  \begin{tikzpicture}[remember picture,overlay]
    \fill[poporange!18] ([xshift=-3.2cm,yshift=-3.6cm]current page.north east) circle (4.6cm);
    \fill[poppurple!14] ([xshift=2.4cm,yshift=7.6cm]current page.south west) circle (3.8cm);
  \end{tikzpicture}%
  {\poplogo\fontsize{30}{30}\selectfont\color{popink}OnBuch\color{poporange}+}\hfill
  \raisebox{0.6ex}{\poppill{Cours signé \textbullet\ Premium}{popink}{popcream}}\par
  \vspace{1pt}\popsquiggle{poppurple}\par
  \vspace{8pt}
  \begin{tcolorbox}[enhanced,colback=poporange,colframe=popink,boxrule=2.8pt,arc=13pt,
      drop shadow=popink,left=18pt,right=18pt,top=12pt,bottom=13pt,after skip=10pt]
    \poppill{#2 \textbullet\ #3}{popink}{popcream}\hspace{5pt}\poppill{#4}{white}{popink}\par\vspace{8pt}
    {\popdisplay\fontsize{14}{14}\selectfont\color{popink}LEÇON #5}\par\vspace{4pt}
    {\raggedright\hyphenpenalty=10000\exhyphenpenalty=10000\popdisplay\fontsize{25}{27}\selectfont\color{popcream}\MakeUppercase{#1}\par}
    \vspace{8pt}
    {\popxbold\fontsize{9.5}{9.5}\selectfont\color{popink}\MakeUppercase{Durée conseillée : #6}}
  \end{tcolorbox}
  \begin{tcolorbox}[enhanced,colback=white,colframe=popink,boxrule=2.2pt,arc=10pt,
      drop shadow=popink,left=16pt,right=16pt,top=10pt,bottom=10pt,after skip=8pt]
    \poppill{Dans cette leçon}{poppurple}{white}\par\vspace{5pt}
    \popsemi\fontsize{9.3}{12.6}\selectfont\color{popink}#7
  \end{tcolorbox}
  \noindent\begin{minipage}[t]{0.487\linewidth}
    \begin{tcolorbox}[enhanced,colback=popgreen,colframe=popink,boxrule=2.2pt,arc=10pt,
        drop shadow=popink,left=11pt,right=11pt,top=10pt,bottom=10pt,height=3.1cm,valign=center]
      \tcbox[colback=white,colframe=popink,boxrule=1.3pt,arc=5pt,boxsep=0pt,nobeforeafter,
        left=3pt,right=3pt,top=3pt,bottom=3pt,valign=center]{\qrcode[height=1.9cm]{https://play.google.com/store/apps/details?id=com.luvvix.onbuch&hl=fr}}%
      \hspace{8pt}\begin{minipage}[c]{0.5\linewidth}
        {\popdisplay\fontsize{11}{12.5}\selectfont\color{white}\MakeUppercase{Télécharge OnBuch}}\par\vspace{4pt}
        {\raggedright\popsemi\fontsize{8.4}{11}\selectfont\color{white}Gratuit \textbullet\ Google Play\\Tuteur IA, annales, concours, résultats\par}
      \end{minipage}
    \end{tcolorbox}
  \end{minipage}\hfill
  \begin{minipage}[t]{0.487\linewidth}
    \begin{tcolorbox}[enhanced,colback=poppurple,colframe=popink,boxrule=2.2pt,arc=10pt,
        drop shadow=popink,left=11pt,right=11pt,top=10pt,bottom=10pt,height=3.1cm,valign=center]
      \tikz[baseline=-0.7ex]\node[circle,fill=white,draw=popink,line width=1.3pt,minimum size=1.6cm,inner sep=0]{\popdisplay\fontsize{24}{24}\selectfont\color{poppurple}+};%
      \hspace{8pt}\begin{minipage}[c]{0.6\linewidth}
        {\popdisplay\fontsize{11}{12.5}\selectfont\color{white}\MakeUppercase{Passe à OnBuch+}}\par\vspace{4pt}
        {\raggedright\popsemi\fontsize{8.4}{11}\selectfont\color{white}Tous les cours signés, Tuteur IA illimité, fiches PDF — onglet OnBuch+ de l'app\par}
      \end{minipage}
    \end{tcolorbox}
  \end{minipage}
  \vspace{8pt}
  \popcontenu
  \vfill
  \signatureonbuch
  \restoregeometry
  \newpage
}

% Rangée « Ce cours contient » (page de garde)
\newcommand{\poptile}[3]{% #1 couleur #2 gros texte #3 légende
  \begin{tcolorbox}[enhanced,colback=#1,colframe=popink,boxrule=2pt,arc=9pt,shadow={2.5pt}{-2.5pt}{0pt}{popink},
      left=6pt,right=6pt,top=9pt,bottom=9pt,halign=center,valign=center,height=1.75cm,width=\linewidth,nobeforeafter]
    {\popdisplay\fontsize{12.5}{13}\selectfont\color{white}\MakeUppercase{#2}}\par\vspace{3pt}
    {\centering\popsemi\fontsize{7.8}{9.5}\selectfont\color{white}#3\par}
  \end{tcolorbox}}
\newcommand{\popcontenu}{%
  \par\noindent{\popxboldls\fontsize{7.5}{7.5}\selectfont\color{popmuted}CE COURS SIGNÉ CONTIENT}\par\vspace{7pt}
  \noindent\begin{minipage}[t]{0.235\linewidth}\poptile{popblue}{Cours}{complet, illustré, pas à pas}\end{minipage}\hfill
  \begin{minipage}[t]{0.235\linewidth}\poptile{poporange}{Méthodes}{et exercices résolus}\end{minipage}\hfill
  \begin{minipage}[t]{0.235\linewidth}\poptile{poppink}{Exercices}{type examen + corrigés}\end{minipage}\hfill
  \begin{minipage}[t]{0.235\linewidth}\poptile{popgreen}{Fiche bilan}{l'essentiel à retenir}\end{minipage}\par}

% Sommaire
\newtcolorbox{sommaire}{enhanced,colback=white,colframe=popink,boxrule=2.2pt,arc=10pt,
  drop shadow=popink,left=18pt,right=18pt,top=13pt,bottom=13pt,before skip=6pt,after skip=20pt,
  before upper={\poppill{Sommaire}{poppurple}{white}\par\vspace{9pt}\popsemi\fontsize{10.6}{14.5}\selectfont\color{popink}}}

% Signature de fin de cours — poussée tout en bas de la dernière page
\newcommand{\findecours}{%
  \par\vfill\nopagebreak
  \begin{center}\popsquiggle{poporange}\end{center}\vspace{4pt}
  \signatureonbuch
}
\AtBeginDocument{\def\llbracket{\mathopen{[\mkern-3.5mu[}}\def\rrbracket{\mathclose{]\mkern-3.5mu]}}} % intervalles d'entiers (glyphe ⟦ absent de la police)
% Glyphes absents de Fira Math : redéfinis à partir de glyphes disponibles
\AtBeginDocument{%
  \def\setminus{\mathbin{\backslash}}\let\smallsetminus\setminus
  \def\vdots{\mathord{\vbox{\baselineskip3.2pt\lineskiplimit0pt\kern4pt\hbox{.}\hbox{.}\hbox{.}}}}%
  \def\ddots{\mathinner{\mkern1mu\raise6.5pt\hbox{.}\mkern2mu\raise3.6pt\hbox{.}\mkern2mu\raise0.7pt\hbox{.}\mkern1mu}}%
  \def\triangle{\mathord{\scalebox{0.9}{$\symup{\Delta}$}}}%
  \def\nabla{\mathord{\raisebox{\depth}{\scalebox{1}[-1]{$\symup{\Delta}$}}}}%
  \def\checkmark{\popcheck{popdark}}%
  \def\star{\mathbin{\ast}}\def\bigstar{\mathord{\ast}}%
  \def\diamond{\mathbin{\scalebox{0.6}{$\Diamond$}}}%
  \def\bigcup{\mathop{\mathchoice{\raisebox{-0.3em}{\scalebox{1.7}{$\cup$}}}{\scalebox{1.25}{$\cup$}}{\cup}{\cup}}\displaylimits}%
  \def\bigcap{\mathop{\mathchoice{\raisebox{-0.3em}{\scalebox{1.7}{$\cap$}}}{\scalebox{1.25}{$\cap$}}{\cap}{\cap}}\displaylimits}%
  \def\lesseqqgtr{\mathrel{\lessgtr}}\def\bigoplus{\mathop{\oplus}\displaylimits}%
}
\providecommand{\ssi}{si et seulement si} % abréviation fréquente
\AtBeginDocument{\providecommand{\wideparen}{\overparen}} % arc de cercle

%% FIGFIX-BEGIN v3 — correctifs globaux des figures (docs/onbuch-generation/tools/figfix)
\usepackage{adjustbox}\usepackage{etoolbox}
% toute figure TikZ/pgfplots plus large que la ligne est réduite à la largeur disponible
\BeforeBeginEnvironment{tikzpicture}{\begin{adjustbox}{max width=\linewidth}}
\AfterEndEnvironment{tikzpicture}{\end{adjustbox}}
% légendes pgfplots lisibles par-dessus les courbes
\pgfplotsset{every axis/.append style={legend style={fill=white,fill opacity=0.92,text opacity=1,draw=black!15,inner sep=3pt}}}
% étiquettes d'arêtes (syntaxe edge["..."]) avec fond blanc
\tikzset{every edge quotes/.append style={fill=white,inner sep=1.2pt}}
% bulles de cartes mentales : texte à la ligne dans la bulle, bulle un peu plus grande
\tikzset{every concept/.append style={text width=2.0cm,align=center,minimum size=2.3cm,inner sep=2pt}}
%% FIGFIX-END
\begin{document}

\pagegarde{Lesson 8 — Grammar: Revise question tags}{Anglais}{Tle A}{Chapter 1 : Family and Social Life}{EN08}{2 h}{Cette leçon de grammaire révise en profondeur les question tags (tags questions) : leur formation, leur accord avec la phrase principale, leurs emplois et leur intonation. Elle insiste sur les cas particuliers et sur les erreurs typiques des apprenants francophones, avec des exercices progressifs de type Bac.
\vspace{4pt}
\begin{multicols}{2}\small
\begin{itemize}
\item À la fin de cette leçon, je sais reconnaître un question tag dans une phrase ou un dialogue.
\item À la fin de cette leçon, je sais former correctement un tag à partir d'une phrase affirmative (tag négatif) et d'une phrase négative (tag affirmatif).
\item À la fin de cette leçon, je sais choisir le bon auxiliaire (do/does/did, be, have, modaux) pour construire le tag.
\item À la fin de cette leçon, je sais accorder le tag avec le sujet, le temps et l'auxiliaire de la phrase principale.
\item À la fin de cette leçon, je sais traiter les cas particuliers : I am → aren't I, impératifs → will you, Let's → shall we, there is → isn't there, everybody/somebody → they.
\item À la fin de cette leçon, je sais répondre correctement à un tag en anglais (Yes/No selon la réalité, pas selon la forme).
\end{itemize}\end{multicols}}

\begin{sommaire}
\begin{multicols}{2}
\begin{enumerate}
\item Observation et découverte
\item Règle de formation : le principe de base
\item Cas particuliers et exceptions
\item Répondre à un question tag et intonation
\item Comparaison français/anglais et pièges des francophones
\item Entraînement guidé et production
\item Activité d'intégration
\item Méthodes et exercices résolus
\item Exercices
\item Corrigés des exercices
\item Fiche bilan
\end{enumerate}
\end{multicols}
\end{sommaire}

\begin{objectifs}
\begin{itemize}
\item À la fin de cette leçon, je sais reconnaître un question tag dans une phrase ou un dialogue.
\item À la fin de cette leçon, je sais former correctement un tag à partir d'une phrase affirmative (tag négatif) et d'une phrase négative (tag affirmatif).
\item À la fin de cette leçon, je sais choisir le bon auxiliaire (do/does/did, be, have, modaux) pour construire le tag.
\item À la fin de cette leçon, je sais accorder le tag avec le sujet, le temps et l'auxiliaire de la phrase principale.
\item À la fin de cette leçon, je sais traiter les cas particuliers : I am → aren't I, impératifs → will you, Let's → shall we, there is → isn't there, everybody/somebody → they.
\item À la fin de cette leçon, je sais répondre correctement à un tag en anglais (Yes/No selon la réalité, pas selon la forme).
\item À la fin de cette leçon, je sais éviter les erreurs fréquentes des francophones (répétition du nom, mauvais auxiliaire, « isn't it » universel).
\item À la fin de cette leçon, je sais utiliser les question tags dans un dialogue oral ou écrit avec la bonne intonation.
\end{itemize}
\end{objectifs}

\begin{prerequis}
\begin{itemize}
\item Les auxiliaires anglais : be, have, do/does/did et les modaux (can, will, must, should...)
\item La formation des phrases affirmatives et négatives aux temps principaux (present simple, past simple, present perfect, future)
\item Les pronoms personnels sujets (I, you, he, she, it, we, they)
\item La distinction entre verbe lexical et auxiliaire dans une phrase
\item Les questions fermées (Yes/No questions) vues en Seconde et Première
\end{itemize}
\end{prerequis}

\newpage

\coursec{Observation et découverte}

\courssub{Une scène de la vie quotidienne : la récréation au lycée de Bafoussam}

Pendant la récréation au lycée de Bafoussam, Amina s'approche de son ami Tabi, l'air inquiet. Elle lui lance : \eng{You didn't do the English homework, did you?} Tabi répond avec un grand sourire : \eng{Yes, I did!} Amina pousse un soupir de soulagement : il \emph{a bien} fait les devoirs. Un élève francophone aurait traduit la réponse par « si ! », car en français, pour contredire une question négative, on dit « si ». En anglais, pas de « si » : le \eng{yes} ou le \eng{no} s'accorde avec la réalité des faits, pas avec la forme de la question.

Ces petites phrases qui « collent » à la fin d'une phrase, les \cle{question tags}, sont partout dans l'anglais parlé : au marché de Bamenda, dans les séries nigérianes, dans les chansons ghanéennes, dans les conversations des étudiants de Buea. Les maîtriser, c'est parler un anglais naturel et vivant, pas un anglais de manuel scolaire. C'est aussi un point régulièrement testé au Baccalauréat A, à l'épreuve écrite comme à l'oral.

\textbf{Problématique :} comment l'anglais construit-il ces petites questions finales, et pourquoi le « n'est-ce pas » français, si simple, devient-il un casse-tête en anglais ?

\courssub{Lecture et observation : un dialogue entre deux élèves camerounais}

Lisons d'abord ce dialogue entre Amina et Tabi, un vendredi après les cours. Lis-le à voix haute, puis souligne mentalement toutes les petites questions placées en fin de phrase.

\begin{textebox}[At the school gate — a dialogue between two classmates]
\textbf{Amina:} You're going to Douala for the holidays, aren't you?\\
\textbf{Tabi:} Yes, I am. My uncle lives in Bonabéri.\\
\textbf{Amina:} Lucky you! You've been there before, haven't you?\\
\textbf{Tabi:} Yes, several times. But the bus journey is long, isn't it?\\
\textbf{Amina:} It is! You didn't travel alone last time, did you?\\
\textbf{Tabi:} No, I didn't. My cousin came with me.\\
\textbf{Amina:} Your cousin can't swim, can she?\\
\textbf{Tabi:} No, she can't. Why do you ask?\\
\textbf{Amina:} Because you'll go to the beach in Limbe, won't you?\\
\textbf{Tabi:} Of course we will! You aren't jealous, are you?\\
\textbf{Amina:} Me? Never! Well... maybe a little.
\end{textebox}

\textbf{Glossaire :}

\begin{center}
\begin{tabularx}{\linewidth}{|l|l|X|}
\hline
\rowcolor{popblueL} \textbf{Mot anglais} & \textbf{Nature} & \textbf{Traduction française} \\
\hline
holidays & nom pluriel & les vacances \\
\hline
lucky you & expression & chanceux ! quelle chance ! \\
\hline
journey & nom & le trajet, le voyage \\
\hline
last time & locution & la dernière fois \\
\hline
to swim & verbe & nager \\
\hline
jealous & adjectif & jaloux, jalouse \\
\hline
\end{tabularx}
\end{center}

\courssub{Observer pour comprendre : quatre constats guidés}

\textbf{Constat n° 1 — La position du tag.} Reprends le dialogue et souligne les six question tags : \eng{aren't you?}, \eng{haven't you?}, \eng{isn't it?}, \eng{did you?}, \eng{can she?}, \eng{won't you?}, \eng{are you?}. Où se trouvent-ils ? \textbf{Toujours à la toute fin de la phrase}, jamais au début, jamais au milieu. On remarque aussi qu'ils sont \textbf{séparés de la phrase principale par une virgule} et \textbf{suivis d'un point d'interrogation}.

\textbf{Constat n° 2 — La polarité inversée.} Regardons de plus près :

\begin{center}
\begin{tabularx}{\linewidth}{|X|X|X|}
\hline
\rowcolor{popblueL} \textbf{Phrase principale} & \textbf{Polarité} & \textbf{Tag} \\
\hline
You're going to Douala... & affirmative (+) & aren't you? (–) \\
\hline
You've been there before... & affirmative (+) & haven't you? (–) \\
\hline
The bus journey is long... & affirmative (+) & isn't it? (–) \\
\hline
You didn't travel alone... & négative (–) & did you? (+) \\
\hline
Your cousin can't swim... & négative (–) & can she? (+) \\
\hline
You'll go to the beach... & affirmative (+) & won't you? (–) \\
\hline
\end{tabularx}
\end{center}

La loi est claire : \textbf{phrase affirmative → tag négatif ; phrase négative → tag affirmatif}. C'est le principe de la \cle{polarité inversée}.

\textbf{Constat n° 3 — Le tag reprend l'auxiliaire et un pronom.} Dans \eng{You've been there before, haven't you?}, le tag reprend l'auxiliaire \eng{have} de la phrase principale et remplace le nom ou le sujet par un \textbf{pronom} (\eng{you}). Jamais on ne dira : \emph{Your cousin can't swim, can't your cousin?} — le nom est toujours remplacé par un pronom personnel sujet (\eng{she}).

\textbf{Constat n° 4 — La ponctuation.} La structure complète est : phrase déclarative + \textbf{virgule} + tag + \textbf{point d'interrogation}. La virgule marque une petite pause à l'oral ; le point d'interrogation signale qu'on attend une réaction de l'interlocuteur.

\begin{definition}[Question tag]
Un \cle{question tag} (ou \emph{tag question}) est une \textbf{courte question ajoutée à la fin d'une phrase déclarative}, formée d'un \textbf{auxiliaire} et d'un \textbf{pronom sujet}. Il sert à \textbf{demander confirmation} d'une information ou à \textbf{inviter l'interlocuteur à réagir}. Équivalent français approximatif : « n'est-ce pas ? », « hein ? », « non ? ».
\end{definition}

\begin{popfigure}
\begin{tikzpicture}[font=\small]
\node[draw=popblue, fill=popblueL, rounded corners, thick, minimum width=4.6cm, minimum height=1.1cm, align=center] (main) at (0,0) {\textbf{Phrase principale}\\ You are tired};
\node[draw=popmuted, thick, circle, inner sep=1.5pt] (comma) at (3.15,0) {\textbf{,}};
\node[draw=poporange, fill=poporangeL, rounded corners, thick, minimum width=3.4cm, minimum height=1.1cm, align=center] (tag) at (5.6,0) {\textbf{Tag : auxiliaire + pronom}\\ aren't you};
\node[draw=popmuted, thick, circle, inner sep=1.5pt] (qm) at (7.9,0) {\textbf{?}};
\node[draw=popgreen, fill=popgreenL, rounded corners, thick, align=center, minimum width=2.2cm] (plus) at (0,-1.9) {polarité \textbf{+}};
\node[draw=poppink, fill=poppinkL, rounded corners, thick, align=center, minimum width=2.2cm] (minus) at (5.6,-1.9) {polarité \textbf{–}};
\draw[-{Stealth}, very thick, popgreen] (plus.east) -- (minus.west) node[midway, above, align=center, font=\footnotesize]{phrase affirmative $\rightarrow$ tag négatif};
\draw[-{Stealth}, very thick, poppink] (minus.west) .. controls (2.8,-3.1) .. (plus.east) node[midway, below, align=center, font=\footnotesize]{phrase négative $\rightarrow$ tag affirmatif};
\end{tikzpicture}
\legende{Structure du question tag : phrase principale + virgule + auxiliaire + pronom + point d'interrogation, avec inversion de polarité.}
\end{popfigure}

\courssub{Exemples traduits : observer la correspondance avec le français}

\begin{exemplebox}[Des exemples pour sentir le mécanisme]
\eng{You are tired, aren't you?} « Tu es fatigué, n'est-ce pas ? »\\
\eng{He didn't call, did he?} « Il n'a pas appelé, n'est-ce pas ? »\\
\eng{She speaks English well, doesn't she?} « Elle parle bien anglais, n'est-ce pas ? »\\
\eng{They won't be late, will they?} « Ils ne seront pas en retard, n'est-ce pas ? »\\
\eng{You haven't eaten, have you?} « Tu n'as pas mangé, si ? »
\end{exemplebox}

Remarque bien : dans chaque exemple français, on utilise le \textbf{même} mot invariable « n'est-ce pas », alors qu'en anglais le tag \textbf{change à chaque fois} : \eng{aren't}, \eng{did}, \eng{doesn't}, \eng{will}, \eng{have}... C'est précisément là que réside la difficulté pour nous, francophones.

\courssub{La fonction communicative : à quoi sert un question tag ?}

Pourquoi les Anglophones ajoutent-ils ces petites questions ? Trois grandes fonctions :

\begin{itemize}
\item \textbf{Demander confirmation} d'une information dont on est presque sûr : \eng{You come from Douala, don't you?} « Tu viens de Douala, n'est-ce pas ? » — le locuteur croit savoir et veut vérifier.
\item \textbf{Engager la conversation}, créer du lien : \eng{The match wasn't boring, was it?} « Le match n'était pas ennuyeux, hein ? » — on invite l'autre à donner son avis, comme au sujet du football après une victoire des Lions Indomptables.
\item \textbf{Adoucir une affirmation} ou une remarque, pour ne pas paraître brutal : \eng{She's been to London, hasn't she?} — on affirme, mais on laisse une porte ouverte à la contradiction.
\end{itemize}

\begin{exoresolu}[Identifier les tags dans le dialogue]
\textbf{Énoncé.} Relevez dans le dialogue d'Amina et Tabi trois question tags, indiquez pour chacun : (a) la polarité de la phrase principale ; (b) l'auxiliaire repris ; (c) le pronom utilisé.
\tcblower
\textbf{Solution détaillée.}
\begin{enumerate}
\item \eng{You're going to Douala for the holidays, aren't you?} — (a) phrase \textbf{affirmative} ; (b) auxiliaire \eng{are} (verbe \eng{be} au présent), passé à la forme négative contractée \eng{aren't} ; (c) pronom \eng{you}, qui reprend le sujet \eng{you}.
\item \eng{You didn't travel alone last time, did you?} — (a) phrase \textbf{négative} (\eng{didn't}) ; (b) auxiliaire \eng{did} (marque du prétérit), repris à la forme affirmative ; (c) pronom \eng{you}.
\item \eng{Your cousin can't swim, can she?} — (a) phrase \textbf{négative} (\eng{can't}) ; (b) auxiliaire modal \eng{can}, repris à la forme affirmative ; (c) pronom \eng{she}, qui remplace le nom \eng{your cousin} — jamais le nom lui-même dans le tag.
\end{enumerate}
\end{exoresolu}

\begin{attention}[Pièges des francophones]
\begin{itemize}
\item Ne traduis jamais « n'est-ce pas » mot à mot par \emph{isn't it} dans toutes les phrases ! \eng{You are tired, isn't it?} est \textbf{faux} : il faut \eng{aren't you?}. Le tag dépend du sujet et de l'auxiliaire de \emph{ta} phrase.
\item N'oublie pas la \textbf{virgule} avant le tag et le \textbf{point d'interrogation} après.
\item Ne répète jamais le nom dans le tag : \emph{Your cousin can't swim, can't your cousin?} est incorrect ; dis \eng{can she?}.
\item Attention à la réponse : à une phrase négative comme \eng{You didn't do the homework, did you?}, on répond \eng{Yes, I did} si on l'a fait (le français dirait « si ! »). Le \eng{yes/no} anglais suit les faits, pas la forme de la question.
\end{itemize}
\end{attention}

\begin{savaistu}[Pourquoi le question tag est-il si difficile pour les francophones ?]
En français, « n'est-ce pas » est \textbf{invariable} : on peut le coller derrière n'importe quelle phrase sans réfléchir (« Tu viens, n'est-ce pas ? », « Il est parti, n'est-ce pas ? »). En anglais, le tag est une véritable \textbf{mini-phrase} qui doit s'accorder avec le sujet, le temps et l'auxiliaire de la phrase principale : \eng{You came, didn't you?} / \eng{He has left, hasn't he?} / \eng{She will come, won't she?}. C'est ce qui le rend difficile... mais aussi si élégant quand on le maîtrise ! Petit détail culturel : les Anglais utilisent les tags sans cesse pour être polis et éviter les affirmations trop directes — un vrai marqueur de conversation naturelle.
\end{savaistu}

\begin{exercice}[À toi de jouer]
Souligne les question tags dans les phrases suivantes, puis indique si la phrase principale est affirmative ou négative :
\begin{enumerate}
\item \eng{The match wasn't boring, was it?}
\item \eng{You can't speak Hausa, can you?}
\item \eng{She's been to London, hasn't she?}
\item \eng{You come from Douala, don't you?}
\end{enumerate}
\end{exercice}

\begin{corrige}[Exercice « À toi de jouer »]
\begin{enumerate}
\item Tag : \eng{was it?} — phrase principale \textbf{négative} (\eng{wasn't}).
\item Tag : \eng{can you?} — phrase principale \textbf{négative} (\eng{can't}).
\item Tag : \eng{hasn't she?} — phrase principale \textbf{affirmative} (\eng{'s been} = \eng{has been}).
\item Tag : \eng{don't you?} — phrase principale \textbf{affirmative} (\eng{come}).
\end{enumerate}
\end{corrige}

\begin{aretenir}[L'essentiel de cette première étape]
Un question tag est une courte question placée \textbf{à la fin} d'une phrase, après une \textbf{virgule}, suivie d'un \textbf{point d'interrogation}. Il est formé d'un \textbf{auxiliaire} (repris de la phrase principale) et d'un \textbf{pronom sujet} (jamais le nom). Sa polarité est \textbf{inversée} : phrase affirmative → tag négatif ; phrase négative → tag affirmatif. Il sert à demander confirmation, engager la conversation ou adoucir une affirmation.
\end{aretenir}

\coursec{Règle de formation : le principe de base}

Dans la section précédente, nous avons observé des phrases comme \eng{Your brother plays for the school team, doesn't he?} et nous avons remarqué qu'un \cle{question tag} est une petite question accrochée à la fin d'une phrase. Nous allons maintenant apprendre à le \textbf{construire} de manière systématique. Bonne nouvelle : la méthode est mécanique et fonctionne dans presque tous les cas. Il suffit de suivre trois étapes, dans l'ordre.

\courssub{Les trois étapes de construction}

\begin{methode}[Construire un question tag en 3 questions]
Pour fabriquer un question tag correct, pose-toi toujours ces trois questions, dans cet ordre :
\begin{enumerate}
\item \textbf{Quel est l'auxiliaire de la phrase ?} Cherche \eng{be}, \eng{have} ou un modal (\eng{can}, \eng{will}, \eng{must}...). S'il n'y en a pas (present simple ou past simple), note le temps du verbe : tu utiliseras \eng{do}, \eng{does} ou \eng{did}.
\item \textbf{La phrase est-elle affirmative ou négative ?} Si elle est affirmative, le tag sera négatif ; si elle est négative, le tag sera affirmatif. On inverse toujours la polarité.
\item \textbf{Quel pronom remplace le sujet ?} Un nom est toujours remplacé par un pronom : \eng{John} $\rightarrow$ \eng{he}, \eng{the girls} $\rightarrow$ \eng{they}, \eng{the bus} $\rightarrow$ \eng{it}, \eng{Mary and I} $\rightarrow$ \eng{we}.
\end{enumerate}
\end{methode}

Appliquons la méthode sur un exemple complet. Phrase de départ : \eng{Your father works in Yaoundé.}

\begin{enumerate}
\item Auxiliaire ? Aucun auxiliaire apparent ; le verbe \eng{works} est au present simple, 3\up{e} personne $\rightarrow$ auxiliaire du tag : \eng{does}.
\item Polarité ? La phrase est affirmative $\rightarrow$ tag négatif : \eng{doesn't}.
\item Pronom ? \eng{Your father} $\rightarrow$ \eng{he}.
\end{enumerate}

Résultat : \eng{Your father works in Yaoundé, doesn't he?} « Ton père travaille à Yaoundé, n'est-ce pas ? »

\begin{propriete}[La règle d'or des question tags]
\begin{itemize}
\item \textbf{PHRASE affirmative $\rightarrow$ TAG négatif} : \eng{She is late, isn't she?}
\item \textbf{PHRASE négative $\rightarrow$ TAG affirmatif} : \eng{She isn't late, is she?}
\item Le tag reprend \textbf{le même auxiliaire} que la phrase (ou \eng{do/does/did} s'il n'y en a pas), suivi du \textbf{pronom sujet}.
\item Le tag négatif est \textbf{toujours contracté} : \eng{isn't}, \eng{aren't}, \eng{don't}, \eng{didn't}, \eng{can't}, \eng{won't}, \eng{hasn't}. Les formes non contractées (\eng{is she not?}) sont très rares et réservées à un style extrêmement soutenu.
\end{itemize}
\end{propriete}

\begin{attention}[Le piège numéro 1 du francophone : « isn't it? » partout]
En français, \eng{n'est-ce pas} est invariable : « Tu parles français, n'est-ce pas ? », « Il est parti, n'est-ce pas ? ». Beaucoup d'élèves calquent ce fonctionnement et mettent \eng{isn't it?} à la fin de toutes les phrases. C'est faux en anglais : le tag \textbf{change} selon l'auxiliaire, le temps et le sujet.

\begin{center}
\begin{tabularx}{\linewidth}{|X|X|}\hline
\rowcolor{poppinkL} \textbf{Phrase fausse} & \textbf{Phrase correcte} \\ \hline
You speak French, isn't it? & You speak French, don't you? \\ \hline
They came yesterday, isn't it? & They came yesterday, didn't they? \\ \hline
She can drive, isn't it? & She can drive, can't she? \\ \hline
We are late, isn't it? & We are late, aren't we? \\ \hline
\end{tabularx}
\end{center}

Retiens : \eng{isn't it?} ne convient que lorsque la phrase contient \eng{is} avec un sujet à la 3\up{e} personne du singulier, par exemple \eng{It is hot today, isn't it?}
\end{attention}

\courssub{Étape 1 : identifier l'auxiliaire de la phrase principale}

L'auxiliaire est la petite pièce du moteur verbal. En anglais, il y a trois familles :

\begin{itemize}
\item \textbf{\eng{be}} (am, is, are, was, were) : \eng{She \textbf{is} tired.} / \eng{They \textbf{were} playing.}
\item \textbf{\eng{have}} (have, has, had) : \eng{We \textbf{have} finished.} / \eng{He \textbf{has} left.}
\item \textbf{les modaux} : \eng{can/could}, \eng{will/would}, \eng{shall/should}, \eng{may/might}, \eng{must}, \eng{ought to} : \eng{You \textbf{can} swim.} / \eng{She \textbf{will} come.}
\end{itemize}

Si la phrase ne contient \textbf{aucun} de ces auxiliaires, c'est que le verbe est au \textbf{present simple} ou au \textbf{past simple}. Dans ce cas, on introduit l'auxiliaire invisible : \eng{do}, \eng{does} (present) ou \eng{did} (past).

\begin{exemplebox}[Repérer l'auxiliaire]
\begin{itemize}
\item \eng{The students \textbf{are} revising.} $\rightarrow$ auxiliaire : \eng{are} (be au present continu).
\item \eng{Manga \textbf{has} sold his cocoa.} $\rightarrow$ auxiliaire : \eng{has} (present perfect).
\item \eng{You \textbf{will} pass the exam.} $\rightarrow$ auxiliaire : \eng{will} (modal).
\item \eng{They \textbf{like} fufu.} $\rightarrow$ aucun auxiliaire apparent ; present simple $\rightarrow$ \eng{do}.
\item \eng{She \textbf{called} me yesterday.} $\rightarrow$ aucun auxiliaire apparent ; past simple $\rightarrow$ \eng{did}.
\end{itemize}
\end{exemplebox}

\courssub{Étape 2 : repérer la polarité et l'inverser}

La \cle{polarité} est le signe de la phrase : affirmative (+) ou négative (−). Le tag prend toujours le signe \textbf{opposé}.

\begin{itemize}
\item Phrase affirmative : \eng{You are coming,} $\rightarrow$ tag négatif : \eng{\textbf{aren't you}?} « Tu viens, n'est-ce pas ? »
\item Phrase négative : \eng{You aren't coming,} $\rightarrow$ tag affirmatif : \eng{\textbf{are you}?} « Tu ne viens pas, si ? »
\end{itemize}

\begin{attention}[Une phrase négative peut ne pas contenir « not »]
Certains mots rendent la phrase négative sans \eng{not} : \eng{never}, \eng{no}, \eng{nobody}, \eng{nothing}, \eng{hardly}, \eng{scarcely}, \eng{rarely}, \eng{seldom}. Devant eux, le tag est \textbf{affirmatif} :
\begin{itemize}
\item \eng{She \textbf{never} complains, \textbf{does she}?} « Elle ne se plaint jamais, si ? »
\item \eng{\textbf{Nobody} came, \textbf{did they}?} « Personne n'est venu, si ? »
\item \eng{He \textbf{hardly} works, \textbf{does he}?} « Il ne travaille guère, si ? »
\end{itemize}
Erreur fréquente : ✗ \eng{She never complains, doesn't she?} — il y aurait deux négations.
\end{attention}

\courssub{Étape 3 : remplacer le sujet par le pronom correspondant}

Le sujet du tag est \textbf{toujours un pronom personnel}, jamais un nom. C'est une différence majeure avec le français, où l'on répète souvent le nom ou où « n'est-ce pas » n'a pas de sujet.

\begin{center}
\begin{tabularx}{\linewidth}{|X|X|X|}\hline
\rowcolor{popblueL} \textbf{Sujet de la phrase} & \textbf{Pronom du tag} & \textbf{Exemple} \\ \hline
John / my father / the teacher & he & John is tall, isn't he? \\ \hline
Mary / your sister / the girl & she & Mary sings well, doesn't she? \\ \hline
the bus / the dog / it & it & The bus is late, isn't it? \\ \hline
the girls / my parents / they & they & The girls left, didn't they? \\ \hline
Mary and I & we & We won, didn't we? \\ \hline
you and your brother & you & You two are ready, aren't you? \\ \hline
\end{tabularx}
\end{center}

\begin{exemplebox}[Exemples traduits]
\begin{itemize}
\item \eng{Your father works in Yaoundé, doesn't he?} « Ton père travaille à Yaoundé, n'est-ce pas ? »
\item \eng{We haven't finished, have we?} « Nous n'avons pas fini, si ? »
\item \eng{The bendskin drivers are waiting, aren't they?} « Les moto-taximen attendent, n'est-ce pas ? »
\item \eng{Amina's grandmother doesn't speak English, does she?} « La grand-mère d'Amina ne parle pas anglais, si ? »
\end{itemize}
\end{exemplebox}

\courssub{Cas 1 : la phrase a un auxiliaire apparent — on le reprend}

C'est le cas le plus simple : l'auxiliaire de la phrase est simplement \textbf{recopié} dans le tag, avec la polarité inversée.

\begin{itemize}
\item \eng{She \textbf{is} late, \textbf{isn't} she?} « Elle est en retard, n'est-ce pas ? »
\item \eng{They \textbf{have} left, \textbf{haven't} they?} « Ils sont partis, n'est-ce pas ? »
\item \eng{You \textbf{were} sleeping, \textbf{weren't} you?} « Tu dormais, n'est-ce pas ? »
\item \eng{He \textbf{hasn't} paid, \textbf{has} he?} « Il n'a pas payé, si ? »
\item \eng{We \textbf{aren't} lost, \textbf{are} we?} « Nous ne sommes pas perdus, si ? »
\end{itemize}

\courssub{Cas 2 : present simple ou past simple — on introduit do, does ou did}

Quand la phrase n'a pas d'auxiliaire, on utilise l'auxiliaire du temps correspondant : \eng{do} (I, you, we, they), \eng{does} (he, she, it) au present simple ; \eng{did} (toutes les personnes) au past simple.

\begin{itemize}
\item \eng{You \textbf{like} fufu, \textbf{don't} you?} « Tu aimes le fufu, n'est-ce pas ? »
\item \eng{She \textbf{teaches} biology, \textbf{doesn't} she?} « Elle enseigne la biologie, n'est-ce pas ? »
\item \eng{He \textbf{came} to Douala, \textbf{didn't} he?} « Il est venu à Douala, n'est-ce pas ? »
\item \eng{They \textbf{didn't} win the match, \textbf{did} they?} « Ils n'ont pas gagné le match, si ? »
\end{itemize}

\begin{attention}[Le verbe principal ne change pas]
Dans le tag, on ne reprend \textbf{jamais} le verbe principal, seulement l'auxiliaire. Et le verbe de la phrase garde sa forme normale :
\begin{itemize}
\item ✗ \eng{She teaches biology, doesn't teach she?}
\item ✗ \eng{He came, didn't came he?}
\item ✓ \eng{She teaches biology, doesn't she?} / \eng{He came, didn't he?}
\end{itemize}
Autre piège : ne pas confondre \eng{does} auxiliaire et \eng{does} verbe. Dans \eng{She does her homework, doesn't she?}, le premier \eng{does} est le verbe \eng{do} (« faire »), le second est l'auxiliaire du tag.
\end{attention}

\courssub{Cas 3 : la phrase contient un modal — on reprend le modal}

Avec les verbes modaux, le tag reprend exactement le même modal, à la polarité inversée.

\begin{itemize}
\item \eng{You \textbf{can} swim, \textbf{can't} you?} « Tu sais nager, n'est-ce pas ? »
\item \eng{She \textbf{won't} come, \textbf{will} she?} « Elle ne viendra pas, si ? »
\item \eng{We \textbf{should} leave now, \textbf{shouldn't} we?} « Nous devrions partir maintenant, non ? »
\item \eng{He \textbf{must} work hard, \textbf{mustn't} he?} « Il doit travailler dur, n'est-ce pas ? »
\item \eng{They \textbf{couldn't} hear us, \textbf{could} they?} « Ils ne pouvaient pas nous entendre, si ? »
\end{itemize}

Note la contraction irrégulière : \eng{will not} $\rightarrow$ \eng{won't} (et non « willn't »), \eng{can not} $\rightarrow$ \eng{can't} (un seul \emph{n} dans la forme contractée écrite), \eng{shall not} $\rightarrow$ \eng{shan't}.

\courssub{Tableau récapitulatif des formes}

\begin{center}
\begin{tabularx}{\linewidth}{|l|X|X|}\hline
\rowcolor{popblueL} \textbf{Temps / auxiliaire} & \textbf{Phrase + / tag −} & \textbf{Phrase − / tag +} \\ \hline
be (present) & She is tired, isn't she? & She isn't tired, is she? \\ \hline
be (past) & They were late, weren't they? & He wasn't there, was he? \\ \hline
present simple & You like tea, don't you? & He doesn't smoke, does he? \\ \hline
past simple & She called, didn't she? & They didn't come, did they? \\ \hline
present perfect & We have eaten, haven't we? & He hasn't finished, has he? \\ \hline
will & You will help, won't you? & She won't forget, will she? \\ \hline
can & He can drive, can't he? & They can't swim, can they? \\ \hline
must & We must go, mustn't we? & She mustn't cheat, must she? \\ \hline
should & You should rest, shouldn't you? & He shouldn't lie, should he? \\ \hline
\end{tabularx}
\end{center}

\courssub{Organigramme de décision}

\begin{popfigure}
\begin{center}
\begin{tikzpicture}[
node distance=1.0cm,
every node/.style={font=\small},
start/.style={rectangle, rounded corners, draw=popblue, fill=popblueL, align=center, inner sep=6pt},
decision/.style={diamond, draw=poporange, fill=poporangeL, align=center, inner sep=2pt, aspect=2.0},
action/.style={rectangle, rounded corners, draw=popgreen, fill=popgreenL, align=center, inner sep=6pt},
final/.style={rectangle, rounded corners, draw=poppurple, fill=poppurpleL, align=center, inner sep=6pt},
arrow/.style={-{Stealth[length=2.5mm]}, thick, popdark}
]
\node[start] (s) {Phrase de départ};
\node[decision, below=of s, align=center] (d){La phrase a-t-elle\\un auxiliaire ?\\(be, have, modal)};
\node[action, below left=1.5cm and 1.5cm of d, align=center] (yes){OUI : reprendre\\le même auxiliaire\\(is, have, can, will...)};
\node[action, below right=1.5cm and 1.5cm of d, align=center] (no){NON : utiliser\\do / does (present)\\ou did (past)};
\node[decision, below=3.5cm of d, align=center] (d2){Phrase affirmative\\ou négative ?};
\node[action, below left=1.5cm and 1.5cm of d2, align=center] (neg){Affirmative :\\tag NÉGATIF\\contracté (isn't, don't...)};
\node[action, below right=1.5cm and 1.5cm of d2, align=center] (pos){Négative :\\tag AFFIRMATIF\\(is, do, will...)};
\node[final, below=3.5cm of d2, align=center] (f){Ajouter le pronom sujet\\(he, she, it, we, they...)};
\draw[arrow] (s) -- (d);
\draw[arrow] (d) -| node[above left, font=\footnotesize]{oui} (yes);
\draw[arrow] (d) -| node[above right, font=\footnotesize]{non} (no);
\draw[arrow] (yes) |- (d2);
\draw[arrow] (no) |- (d2);
\draw[arrow] (d2) -| node[above left, font=\footnotesize]{affirmative} (neg);
\draw[arrow] (d2) -| node[above right, font=\footnotesize]{négative} (pos);
\draw[arrow] (neg) |- (f);
\draw[arrow] (pos) |- (f);
\end{tikzpicture}
\end{center}
\legende{Organigramme de décision : construire un question tag en trois étapes — auxiliaire, polarité, pronom.}
\end{popfigure}

\courssub{Application guidée}

\begin{exoresolu}[Construire les tags étape par étape]
Complète avec le question tag correct : (a) \eng{Your mother works at the market, ...?} (b) \eng{The students haven't done their homework, ...?} (c) \eng{Paul can speak three languages, ...?} (d) \eng{It rained heavily last night, ...?}
\tcblower
\textbf{Solution détaillée, avec la méthode en 3 questions :}

(a) \eng{Your mother works at the market, \textbf{doesn't she}?}
1) Pas d'auxiliaire apparent ; \eng{works} = present simple, 3\up{e} personne $\rightarrow$ \eng{does}. 2) Phrase affirmative $\rightarrow$ tag négatif : \eng{doesn't}. 3) \eng{Your mother} $\rightarrow$ \eng{she}.
Traduction : « Ta mère travaille au marché, n'est-ce pas ? »

(b) \eng{The students haven't done their homework, \textbf{have they}?}
1) Auxiliaire apparent : \eng{have} (present perfect). 2) Phrase négative (\eng{haven't}) $\rightarrow$ tag affirmatif : \eng{have}. 3) \eng{The students} $\rightarrow$ \eng{they}.
Traduction : « Les élèves n'ont pas fait leurs devoirs, si ? »

(c) \eng{Paul can speak three languages, \textbf{can't he}?}
1) Auxiliaire : le modal \eng{can}. 2) Phrase affirmative $\rightarrow$ tag négatif contracté : \eng{can't}. 3) \eng{Paul} $\rightarrow$ \eng{he}.
Traduction : « Paul parle trois langues, n'est-ce pas ? »

(d) \eng{It rained heavily last night, \textbf{didn't it}?}
1) Pas d'auxiliaire apparent ; \eng{rained} = past simple $\rightarrow$ \eng{did}. 2) Phrase affirmative $\rightarrow$ tag négatif : \eng{didn't}. 3) \eng{It} reste \eng{it}.
Traduction : « Il a beaucoup plu la nuit dernière, n'est-ce pas ? »
\end{exoresolu}

\begin{aretenir}[L'essentiel de la règle de formation]
\begin{itemize}
\item Un question tag = \textbf{auxiliaire + pronom sujet}, à la polarité opposée de la phrase.
\item Auxiliaire apparent (be, have, modal) $\rightarrow$ on le reprend ; sinon $\rightarrow$ \eng{do/does/did}.
\item Phrase + $\rightarrow$ tag − contracté ; phrase − $\rightarrow$ tag + non contracté.
\item Le sujet du tag est toujours un pronom, jamais un nom.
\item \eng{isn't it?} n'est pas universel : c'est l'erreur de calque la plus fréquente des francophones.
\end{itemize}
\end{aretenir}

\begin{exercice}[À toi de jouer]
Ajoute le question tag correct à chaque phrase :
\begin{enumerate}
\item You live in Bamenda, ...?
\item She hasn't called her brother, ...?
\item The match started at four o'clock, ...?
\item We must respect our parents, ...?
\item Your friends are waiting outside, ...?
\item He won't be late, ...?
\end{enumerate}
\end{exercice}

\begin{corrige}[Exercice « À toi de jouer »]
\begin{enumerate}
\item \eng{You live in Bamenda, don't you?} — present simple, phrase affirmative, sujet \eng{you}.
\item \eng{She hasn't called her brother, has she?} — auxiliaire \eng{have}, phrase négative $\rightarrow$ tag affirmatif.
\item \eng{The match started at four o'clock, didn't it?} — past simple sans auxiliaire $\rightarrow$ \eng{did} ; \eng{the match} $\rightarrow$ \eng{it}.
\item \eng{We must respect our parents, mustn't we?} — modal \eng{must}, phrase affirmative $\rightarrow$ \eng{mustn't}.
\item \eng{Your friends are waiting outside, aren't they?} — auxiliaire \eng{are} ; \eng{your friends} $\rightarrow$ \eng{they}.
\item \eng{He won't be late, will he?} — modal \eng{will} à la forme négative $\rightarrow$ tag affirmatif \eng{will}.
\end{enumerate}
\end{corrige}

\coursec{Cas particuliers et exceptions}

Dans la section précédente, nous avons posé le principe de base du \cle{question tag} : une phrase affirmative appelle un tag négatif (\eng{You live in Yaoundé, don't you?}), une phrase négative appelle un tag affirmatif (\eng{She doesn't smoke, does she?}), et le tag reprend l'auxiliaire de la phrase principale. Cette règle couvre la grande majorité des cas. Mais l'anglais, comme toute langue vivante, comporte des situations où le principe de base ne suffit pas : certains sujets, certains verbes et certains mots imposent des tags spéciaux. C'est précisément sur ces cas particuliers que portent souvent les questions du baccalauréat. Observons-les un à un.

\courssub{Le cas de « I am » : aren't I?}

Commençons par une petite scène de la vie quotidienne. Un élève de terminale arrive en retard au lycée et dit à son camarade :

\begin{exemplebox}[Observation]
\eng{I'm late again, aren't I?} « Je suis encore en retard, n'est-ce pas ? »

\eng{I'm your best friend, aren't I?} « Je suis ton meilleur ami, non ? »
\end{exemplebox}

Avec \eng{I am}, on s'attendrait logiquement à un tag \eng{amn't I?}. Or cette forme \textbf{n'existe pas} en anglais standard : la contraction de \eng{am not} ne s'emploie pas dans les tags. L'anglais utilise à la place \eng{aren't I?}.

\begin{propriete}[Le tag de « I am »]
Après une phrase affirmative à la première personne du singulier avec le verbe \eng{be}, le tag est \cle{aren't I?} (et jamais \eng{amn't I?}).

\eng{I'm right, aren't I?} « J'ai raison, n'est-ce pas ? »

\eng{I'm invited to the wedding, aren't I?} « Je suis invité au mariage, n'est-ce pas ? »
\end{propriete}

\begin{attention}[Erreur classique des francophones]
✗ \eng{I'm right, amn't I?}

✓ \eng{I'm right, aren't I?}

La forme \eng{amn't} existe uniquement dans certains parlers régionaux (Irlande, Écosse) ; elle est à proscrire à l'écrit et à l'examen. Retenez : \textbf{I am → aren't I?}.
\end{attention}

\courssub{Les impératifs : will you? / would you? / can you?}

Observez ces phrases entendues dans une famille camerounaise un samedi matin :

\begin{exemplebox}[Observation]
\eng{Open the window, will you?} « Ouvre la fenêtre, veux-tu ? »

\eng{Pass me the salt, would you?} « Passe-moi le sel, s'il te plaît. »

\eng{Don't be late, will you?} « Ne sois pas en retard, d'accord ? »

\eng{Help me carry this bag, can you?} « Aide-moi à porter ce sac, tu peux ? »
\end{exemplebox}

Un impératif n'a ni sujet exprimé ni auxiliaire : le principe de base ne peut donc pas s'appliquer. L'anglais utilise alors un tag fixe, le plus souvent \eng{will you?}.

\begin{propriete}[Le tag après un impératif]
Après un impératif (affirmatif ou négatif), le tag est généralement \cle{will you?}. On peut aussi utiliser \eng{would you?} (plus poli), \eng{can you?} ou \eng{can't you?} (plus pressant, parfois agacé).

\eng{Sit down, will you?} « Assieds-toi, veux-tu ? »

\eng{Don't forget your keys, will you?} « N'oublie pas tes clés, hein ? »

\eng{Be quiet, can't you?} « Tais-toi, tu veux ? » (ton agacé)
\end{propriete}

\begin{savaistu}[Une fausse question]
\eng{Will you?} après un impératif n'est pas une vraie question : personne ne s'attend à ce que vous répondiez \eng{Yes, I will!} C'est simplement une façon polie — ou pressante selon le ton de la voix — de formuler une demande. En français, on rend cette nuance par « s'il te plaît », « veux-tu ? », « d'accord ? » ou simplement par le ton de la voix.
\end{savaistu}

\courssub{Le cas de « Let's » : shall we?}

\begin{exemplebox}[Observation]
\eng{Let's go to the stadium, shall we?} « Allons au stade, d'accord ? »

\eng{Let's visit Grandma, shall we?} « Allons voir grand-mère, d'accord ? »

\eng{Let's not argue, shall we?} « Ne nous disputons pas, veux-tu ? »
\end{exemplebox}

\eng{Let's} est la contraction de \eng{let us} : c'est une invitation à faire quelque chose ensemble. Le tag qui lui correspond est toujours \cle{shall we?}, car \eng{shall} exprime la proposition à la première personne du pluriel.

\begin{propriete}[Le tag de « Let's »]
Après \eng{Let's...} (y compris la forme négative \eng{Let's not...}), le tag est invariablement \cle{shall we?}.
\end{propriete}

\courssub{There is / There are : isn't there? / aren't there?}

\begin{exemplebox}[Observation]
\eng{There's a problem, isn't there?} « Il y a un problème, n'est-ce pas ? »

\eng{There are many pupils absent today, aren't there?} « Il y a beaucoup d'élèves absents aujourd'hui, n'est-ce pas ? »

\eng{There won't be a meeting tomorrow, will there?} « Il n'y aura pas de réunion demain, n'est-ce pas ? »
\end{exemplebox}

\begin{propriete}[Le tag avec « there »]
Quand la phrase commence par \eng{there is / there are / there will...}, le mot \eng{there} joue le rôle de sujet : on le \textbf{reprend dans le tag}, avec l'auxiliaire correspondant.

\eng{There is} → \eng{isn't there?} ; \eng{there are} → \eng{aren't there?} ; \eng{there won't be} → \eng{will there?}
\end{propriete}

\begin{attention}[Piège de traduction]
Le français « il y a » incite à traduire le sujet par \eng{it} : ✗ \eng{There's a problem, isn't it?} C'est faux. En anglais, \eng{there} est le sujet grammatical de la phrase, c'est donc \eng{there} qui revient dans le tag : ✓ \eng{There's a problem, isn't there?}
\end{attention}

\courssub{Everybody, everyone, somebody, nobody : le pronom « they »}

Observez ces phrases prononcées après une fête de famille à Douala :

\begin{exemplebox}[Observation]
\eng{Everybody enjoyed the party, didn't they?} « Tout le monde a apprécié la fête, n'est-ce pas ? »

\eng{Someone has taken my umbrella, haven't they?} « Quelqu'un a pris mon parapluie, non ? »

\eng{Nobody phoned, did they?} « Personne n'a téléphoné, n'est-ce pas ? »
\end{exemplebox}

En français, « tout le monde », « personne », « quelqu'un » sont grammaticalement singuliers. En anglais aussi, le verbe se met au singulier (\eng{everybody \textbf{is}}). Mais dans le tag, l'anglais utilise le pronom pluriel \cle{they}, car ces mots désignent en réalité un groupe de personnes dont on ignore le sexe.

\begin{propriete}[Les pronoms indéfinis en -body / -one]
Après \eng{everybody, everyone, somebody, someone, nobody, no one, anybody, anyone}, le verbe de la phrase est au \textbf{singulier}, mais le pronom du tag est \cle{they} (avec l'auxiliaire au pluriel).

\eng{Everyone is here, aren't they?} « Tout le monde est là, n'est-ce pas ? »

\eng{Somebody called you, didn't they?} « Quelqu'un t'a appelé, non ? »
\end{propriete}

\begin{attention}[Erreur très fréquente]
✗ \eng{Everyone is here, isn't it?}

✓ \eng{Everyone is here, aren't they?}

Les francophones, influencés par le singulier français (« tout le monde est là »), utilisent à tort \eng{it} ou \eng{he}. Retenez : \textbf{-body / -one → they dans le tag}.
\end{attention}

\courssub{Les mots à sens négatif : le tag devient affirmatif}

Voici maintenant l'un des pièges préférés des examinateurs. Observez :

\begin{exemplebox}[Observation]
\eng{He never calls, does he?} « Il n'appelle jamais, n'est-ce pas ? »

\eng{She hardly eats, does she?} « Elle ne mange presque rien, n'est-ce pas ? »

\eng{Nothing happened, did it?} « Rien ne s'est passé, n'est-ce pas ? »

\eng{Few students understood the lesson, did they?} « Peu d'élèves ont compris la leçon, n'est-ce pas ? »

\eng{Nobody came, did they?} « Personne n'est venu, n'est-ce pas ? »
\end{exemplebox}

Ces phrases ne contiennent ni \eng{not}, ni \eng{n't}. Pourtant, leur \textbf{sens} est négatif : \eng{never} = « ne... jamais », \eng{nobody} = « personne ne... », \eng{hardly} = « à peine », \eng{few} = « peu de ». En anglais, c'est le sens qui commande : une phrase contenant un mot négatif est traitée comme une phrase négative, et le tag est donc \textbf{affirmatif}.

\begin{propriete}[Les mots à sens négatif]
Les mots \eng{never, nobody, no one, nothing, nowhere, hardly, scarcely, barely, few, little} rendent la phrase négative : le question tag est donc \cle{affirmatif} (sans \eng{n't}).

\eng{She hardly eats, does she?} « Elle mange à peine, n'est-ce pas ? »

\eng{He scarcely ever goes out, does he?} « Il ne sort presque jamais, n'est-ce pas ? »

\eng{There is little water left, is there?} « Il reste peu d'eau, n'est-ce pas ? »
\end{propriete}

\begin{attention}[Double négation interdite]
Contrairement au français (« il ne vient jamais »), l'anglais n'accumule jamais deux négations dans la même phrase. On dira \eng{He never comes} ou \eng{He doesn't ever come}, mais ✗ \eng{He doesn't never come}. De même, le tag reste affirmatif : ✗ \eng{He never comes, doesn't he?} → ✓ \eng{He never comes, does he?}
\end{attention}

\courssub{This / that → it ; these / those → they}

\begin{exemplebox}[Observation]
\eng{That's your bag, isn't it?} « C'est ton sac, n'est-ce pas ? »

\eng{This is your pen, isn't it?} « Ceci est ton stylo, n'est-ce pas ? »

\eng{Those are John's shoes, aren't they?} « Ce sont les chaussures de John, n'est-ce pas ? »
\end{exemplebox}

\begin{propriete}[Les démonstratifs dans le tag]
Les pronoms démonstratifs sont remplacés dans le tag par un pronom personnel : \eng{this} et \eng{that} deviennent \cle{it} ; \eng{these} et \eng{those} deviennent \cle{they}.
\end{propriete}

\courssub{I don't think + subordonnée : le tag s'accorde avec la subordonnée}

Dernier cas, précieux pour le baccalauréat. Observez :

\begin{exemplebox}[Observation]
\eng{I don't think he will win, will he?} « Je ne pense pas qu'il gagnera, n'est-ce pas ? »

\eng{I don't believe she is coming, is she?} « Je ne crois pas qu'elle vienne, n'est-ce pas ? »

\eng{I don't suppose they know the truth, do they?} « Je ne suppose pas qu'ils connaissent la vérité, n'est-ce pas ? »
\end{exemplebox}

\begin{propriete}[Le tag après « I don't think / believe / suppose »]
Quand la phrase principale est \eng{I don't think}, \eng{I don't believe} ou \eng{I don't suppose}, la négation porte en réalité sur la subordonnée (le français fait de même : « je ne pense pas qu'il \emph{vienne} »). Le tag s'accorde donc avec la \textbf{subordonnée}, et il est \textbf{affirmatif} : sujet et auxiliaire de la subordonnée, sans \eng{n't}.

\eng{I don't think it will rain, will it?} « Je ne pense pas qu'il pleuvra, n'est-ce pas ? »
\end{propriete}

\begin{attention}[Ne pas accorder le tag avec « I »]
✗ \eng{I don't think he will win, don't I?}

✓ \eng{I don't think he will win, will he?}

Le tag ne porte jamais sur \eng{I don't think} lui-même : il interroge le contenu de la subordonnée.
\end{attention}

\courssub{Tableau récapitulatif des cas particuliers}

\begin{center}
\begin{tabularx}{\linewidth}{|X|X|X|}
\hline
\rowcolor{popblueL} Début de phrase & Tag & Exemple \\
\hline
I am... & aren't I? & I'm late, aren't I? \\
\hline
Impératif (Open..., Don't...) & will you? / would you? & Open the door, will you? \\
\hline
Let's... & shall we? & Let's go, shall we? \\
\hline
There is / are... & isn't / aren't there? & There's a problem, isn't there? \\
\hline
Everybody, someone, nobody... & ...they? & Everyone is here, aren't they? \\
\hline
Never, hardly, few, nothing... & tag affirmatif & He never calls, does he? \\
\hline
This / that... & ...it? & That's mine, isn't it? \\
\hline
These / those... & ...they? & Those are yours, aren't they? \\
\hline
I don't think + subordonnée & tag affirmatif de la subordonnée & I don't think he knows, does he? \\
\hline
\end{tabularx}
\end{center}

\courssub{Carte mentale : les exceptions des question tags}

\begin{popfigure}
\begin{tikzpicture}[mindmap, grow cyclic,
  every node/.style={concept, align=center, font=\small},
  root concept/.append style={concept color=popblue, text=white, font=\small\bfseries},
  level 1 concept/.append style={sibling angle=60, level distance=3.4cm, concept color=poporangeL, text=popdark}]
\node[root concept, align=center]{Question tags\\exceptions}
  child { node[align=center]{I am\\→ aren't I?} }
  child { node[align=center]{Impératifs\\→ will you?} }
  child { node[align=center]{Let's\\→ shall we?} }
  child { node[align=center]{There is/are\\→ isn't there?} }
  child { node[align=center]{Everybody\\→ they} }
  child { node[align=center]{never, hardly\\→ tag affirmatif} };
\end{tikzpicture}
\legende{Carte mentale des cas particuliers des question tags : chaque branche correspond à une exception à mémoriser.}
\end{popfigure}

\courssub{Exercice résolu et entraînement}

\begin{exoresolu}[Complétez avec le question tag correct]
Énoncé. Complétez chaque phrase avec le tag approprié.

1. \eng{I'm your captain, ...?}

2. \eng{Let's watch the football match, ...?}

3. \eng{There's no water in the tank, ...?}

4. \eng{Everybody likes fried plantains, ...?}

5. \eng{She hardly ever complains, ...?}

6. \eng{Don't touch that wire, ...?}

7. \eng{I don't think they will come, ...?}

\tcblower
Solution détaillée.

1. \eng{I'm your captain, \textbf{aren't I}?} — avec \eng{I am}, le tag est toujours \eng{aren't I?} (jamais \eng{amn't I?}).

2. \eng{Let's watch the football match, \textbf{shall we}?} — après \eng{Let's}, tag invariable : \eng{shall we?}.

3. \eng{There's no water in the tank, \textbf{is there}?} — le sujet est \eng{there} ; la phrase est négative (\eng{no water}), donc tag affirmatif.

4. \eng{Everybody likes fried plantains, \textbf{don't they}?} — \eng{everybody} appelle \eng{they} dans le tag ; phrase affirmative, tag négatif.

5. \eng{She hardly ever complains, \textbf{does she}?} — \eng{hardly} a un sens négatif : le tag est affirmatif.

6. \eng{Don't touch that wire, \textbf{will you}?} — après un impératif (même négatif), le tag est \eng{will you?}.

7. \eng{I don't think they will come, \textbf{will they}?} — le tag s'accorde avec la subordonnée (\eng{they will}) et il est affirmatif.
\end{exoresolu}

\begin{exercice}[À vous de jouer]
Complétez avec le question tag correct, puis traduisez chaque phrase en français.

1. \eng{Nobody answered the phone, ...?}

2. \eng{Pass me my bag, ...?}

3. \eng{I'm invited to the party, ...?}

4. \eng{There are many mangoes this year, ...?}

5. \eng{Let's help our parents, ...?}

6. \eng{That's your bicycle, ...?}

7. \eng{Few people came to the meeting, ...?}

8. \eng{I don't believe he is ill, ...?}
\end{exercice}

\begin{corrige}[Exercice « À vous de jouer »]
1. \eng{Nobody answered the phone, did they?} « Personne n'a répondu au téléphone, n'est-ce pas ? » (\eng{nobody} = sens négatif → tag affirmatif ; pronom \eng{they}).

2. \eng{Pass me my bag, will you?} « Passe-moi mon sac, veux-tu ? » (impératif → \eng{will you?}).

3. \eng{I'm invited to the party, aren't I?} « Je suis invité à la fête, n'est-ce pas ? » (\eng{I am} → \eng{aren't I?}).

4. \eng{There are many mangoes this year, aren't there?} « Il y a beaucoup de mangues cette année, n'est-ce pas ? » (\eng{there} repris dans le tag).

5. \eng{Let's help our parents, shall we?} « Aidons nos parents, d'accord ? » (\eng{Let's} → \eng{shall we?}).

6. \eng{That's your bicycle, isn't it?} « C'est ta bicyclette, n'est-ce pas ? » (\eng{that} → \eng{it}).

7. \eng{Few people came to the meeting, did they?} « Peu de gens sont venus à la réunion, n'est-ce pas ? » (\eng{few} = sens négatif → tag affirmatif).

8. \eng{I don't believe he is ill, is he?} « Je ne crois pas qu'il soit malade, n'est-ce pas ? » (tag accordé avec la subordonnée, affirmatif).
\end{corrige}

\begin{aretenir}[L'essentiel de la section]
\begin{itemize}
\item \eng{I am} → \cle{aren't I?} (jamais \eng{amn't I?}).
\item Impératif → \cle{will you?} ; \eng{Let's} → \cle{shall we?}.
\item \eng{There is/are} → le tag reprend \cle{there}.
\item \eng{Everybody, someone, nobody} → pronom \cle{they} dans le tag.
\item \eng{Never, hardly, few, little, nothing} = phrase négative → tag \cle{affirmatif}.
\item \eng{This/that} → \eng{it} ; \eng{these/those} → \eng{they}.
\item \eng{I don't think...} → tag affirmatif accordé avec la \cle{subordonnée}.
\end{itemize}
\end{aretenir}

Maintenant que les cas particuliers sont maîtrisés, une question reste ouverte : comment faut-il \emph{répondre} à un question tag, et quel rôle joue l'intonation ? C'est l'objet de la section suivante.

\coursec{Répondre à un question tag et intonation}

Dans la section précédente, nous avons étudié les cas particuliers et les exceptions des question tags (\eng{Let's go, shall we?}, \eng{I am late, aren't I?}…). Nous savons donc maintenant \emph{former} un tag correctement. Mais un question tag est fait pour obtenir une \emph{réponse} : comment répondre correctement ? Et comment prononcer le tag lui-même ? C'est l'objet de cette section, essentielle pour l'expression orale et pour l'épreuve de compréhension de l'oral.

\courssub{Répondre selon la réalité, pas selon la forme}

\begin{textebox}[A dialogue after a football match]
\textbf{Context:} Two students, Amina and Peter, are talking the day after a match of the Indomitable Lions.

— You didn't watch the match, did you?\\
— Yes, I did! It was fantastic! The Lions played very well.\\
— Really? I thought you had gone to your aunt's house in Mvolyé.\\
— I had gone there, but they have a television. We all watched it together. You didn't miss it, did you?\\
— No, I didn't. I watched it too. What a victory!\\
— The striker scored twice, didn't he?\\
— Yes, he did. And the goalkeeper was brilliant.\\
— You aren't going to school by bendskin tomorrow, are you?\\
— No, I'm not. My father will drive me.\\
— Good. We can talk about the match on the way, can't we?\\
— Yes, we can!
\end{textebox}

\textbf{Glossaire :}
\begin{itemize}
\item \eng{to miss} (v.) : rater, manquer
\item \eng{a victory} (n.) : une victoire
\item \eng{a striker} (n.) : un attaquant, un buteur
\item \eng{to score} (v.) : marquer (un but)
\item \eng{a goalkeeper} (n.) : un gardien de but
\item \eng{brilliant} (adj.) : remarquable, excellent
\item \eng{a bendskin} (n.) : une moto-taxi (Cameroun)
\item \eng{to drive somebody} (v.) : conduire quelqu'un (en voiture)
\item \eng{on the way} (expr.) : en chemin
\end{itemize}

\textbf{Questions de compréhension :}
\begin{enumerate}
\item Did Amina watch the match? How do you know?
\item Where was Amina during the match?
\item How will Amina go to school tomorrow?
\item Find three question tags in the dialogue and their answers.
\end{enumerate}

Observez la première réplique : la question est \emph{négative} (\eng{You didn't watch...}), pourtant Amina répond \eng{Yes, I did!}. Pourquoi \eng{Yes} ? Parce qu'elle a réellement regardé le match. C'est toute la logique des réponses anglaises.

\begin{aretenir}[Le principe fondamental]
En anglais, on répond à un question tag selon la \cle{réalité des faits}, jamais selon la forme (affirmative ou négative) de la question. Si le fait est vrai, on dit \eng{Yes} ; si le fait est faux, on dit \eng{No} — même quand la question est négative.
\end{aretenir}

\begin{exemplebox}[Observer et comparer]
\eng{You aren't Cameroonian, are you?} — \eng{Yes, I am.}\\
« Tu n'es pas camerounais, n'est-ce pas ? » — « Si, je le suis. »

\eng{You aren't Cameroonian, are you?} — \eng{No, I'm not.}\\
« Tu n'es pas camerounais, n'est-ce pas ? » — « Non, je ne le suis pas. »

\eng{She hasn't finished her homework, has she?} — \eng{Yes, she has.}\\
« Elle n'a pas fini ses devoirs, n'est-ce pas ? » — « Si, elle les a finis. »

\eng{It won't rain tomorrow, will it?} — \eng{No, it won't.}\\
« Il ne pleuvra pas demain, n'est-ce pas ? » — « Non, il ne pleuvra pas. »
\end{exemplebox}

\courssub{Le « si » français n'existe pas en anglais}

\begin{propriete}[Contredire une question négative : « Yes » + auxiliaire]
En français, pour contredire une question négative, on utilise le mot spécial \textbf{« si »} : « Tu n'aimes pas le ndolé ? — \textbf{Si}, j'aime ça ! »

En anglais, ce mot \textbf{n'existe pas}. On répond simplement \cle{Yes} suivi de la réponse courte avec l'auxiliaire :

\eng{You don't like ndolé, do you? — Yes, I do!}\\
« Tu n'aimes pas le ndolé, n'est-ce pas ? — Si, j'adore ça ! »

La réponse courte reprend \textbf{l'auxiliaire du tag}, jamais le verbe principal seul.
\end{propriete}

\begin{attention}[Erreur fréquente des francophones : répondre « No » pour dire « si »]
Sous l'influence du français, beaucoup d'élèves répondent \eng{No} quand ils veulent dire « si ». Cela crée une contradiction incompréhensible pour un anglophone :

✗ \eng{You don't like ndolé, do you? — No, I like it.} (contradiction !)\\
✓ \eng{You don't like ndolé, do you? — Yes, I do!}

Autre exemple :\\
✗ \eng{You haven't eaten, have you? — No, I have eaten.}\\
✓ \eng{You haven't eaten, have you? — Yes, I have.}

Retenez : \eng{Yes} = le fait est vrai ; \eng{No} = le fait est faux. La négation de la question ne change rien.
\end{attention}

\begin{propriete}[La réponse courte est obligatoire]
Dans une réponse soignée (à l'oral comme à l'écrit, et surtout à l'examen), on ne dit jamais \eng{Yes} ou \eng{No} tout seul : on ajoute toujours la \cle{réponse courte} avec le sujet et l'auxiliaire.

\begin{center}
\begin{tabularx}{\linewidth}{|X|X|X|}
\hline
\rowcolor{popblueL} Question tag & Réponse affirmative & Réponse négative \\
\hline
You like fufu, don't you? & Yes, I do. & No, I don't. \\
\hline
She has arrived, hasn't she? & Yes, she has. & No, she hasn't. \\
\hline
They can swim, can't they? & Yes, they can. & No, they can't. \\
\hline
He will come, won't he? & Yes, he will. & No, he won't. \\
\hline
You didn't go, did you? & Yes, I did. & No, I didn't. \\
\hline
It is hot, isn't it? & Yes, it is. & No, it isn't. \\
\hline
\end{tabularx}
\end{center}

On peut ensuite ajouter un commentaire : \eng{Yes, I did! It was fantastic!}
\end{propriete}

\begin{center}
\begin{tabularx}{\linewidth}{|X|X|X|}
\hline
\rowcolor{popblueL} Français & English & Remarque \\
\hline
« Tu ne viens pas ? — Si, je viens. » & \eng{You aren't coming, are you? — Yes, I am.} & Pas de « si » en anglais : \eng{Yes} suffit. \\
\hline
« Tu ne viens pas ? — Non, je ne viens pas. » & \eng{You aren't coming, are you? — No, I'm not.} & \eng{No} = le fait est faux. \\
\hline
« Tu as fini, n'est-ce pas ? — Oui. » & \eng{You've finished, haven't you? — Yes, I have.} & Réponse courte obligatoire. \\
\hline
« Il ne pleut pas ? — Non. » & \eng{It isn't raining, is it? — No, it isn't.} & Auxiliaire \eng{is}, pas \eng{does}. \\
\hline
\end{tabularx}
\end{center}

\courssub{L'intonation du tag : descendante ou montante}

L'intonation du question tag change son \emph{sens}. C'est un point capital pour l'épreuve orale et pour la compréhension de l'anglais parlé.

\begin{propriete}[Les deux intonations du question tag]
\textbf{1. Intonation descendante (↘) sur le tag :} le locuteur est \cle{presque sûr} de la réponse ; il attend une simple \emph{confirmation}, un accord. C'est le cas le plus fréquent.

\eng{It's a beautiful day, isn't it? ↘} (« C'est une belle journée, n'est-ce pas ? » — je m'attends à \eng{Yes, it is.})

\textbf{2. Intonation montante (↗) sur le tag :} le locuteur \cle{n'est pas sûr} ; il pose une \emph{vraie question} et attend une information qu'il ne connaît pas.

\eng{You aren't Cameroonian, are you? ↗} (« Tu n'es pas camerounais, si ? » — je ne connais vraiment pas la réponse.)
\end{propriete}

\begin{popfigure}
\begin{tikzpicture}[scale=1]
% Tag commun
\node[align=center, font=\bfseries] (tag) at (0,0) {\eng{You passed the exam,}\\ \eng{didn't you?}};
% Flèche descendante
\draw[-{Stealth[length=3mm]}, line width=1.5pt, popgreen] (-5,0.4) -- (-2.2,0.4) -- (-2.2,-0.6);
\node[align=center, text=popgreen, font=\small] at (-4.6,-1.3) {intonation \textbf{descendante} ↘\\ confirmation attendue\\ (je suis presque sûr)};
% Flèche montante
\draw[-{Stealth[length=3mm]}, line width=1.5pt, poporange] (2.2,-0.6) -- (2.2,0.4) -- (5,0.4);
\node[align=center, text=poporange, font=\small] at (4.6,-1.3) {intonation \textbf{montante} ↗\\ vraie question\\ (je ne sais pas)};
\end{tikzpicture}
\legende{Les deux intonations du question tag : la voix descend quand on attend une confirmation ; elle monte quand on pose une vraie question.}
\end{popfigure}

\begin{exemplebox}[La même phrase, deux sens]
\eng{The concert starts at eight, doesn't it? ↘}\\
« Le concert commence à huit heures, n'est-ce pas ? » (je le sais, je veux juste confirmer)

\eng{The concert starts at eight, doesn't it? ↗}\\
« Le concert commence à huit heures, non ? » (je ne suis pas sûr de l'heure, je demande)

\eng{You've met my sister, haven't you? ↘} (je crois me souvenir que oui)

\eng{You've met my sister, haven't you? ↗} (je ne sais vraiment pas)
\end{exemplebox}

\begin{methode}[Choisir la bonne intonation]
\begin{enumerate}
\item Pose-toi la question : \textbf{« Est-ce que je connais déjà la réponse ? »}
\item \textbf{Oui}, je la connais (ou j'en suis presque sûr) : la voix \textbf{descend} sur le tag ↘. J'attends seulement que mon interlocuteur confirme.
\item \textbf{Non}, je ne la connais pas : la voix \textbf{monte} sur le tag ↗. Je pose une vraie question.
\item À l'écoute, fais l'inverse : si la voix monte, ton interlocuteur attend une vraie information ; si elle descend, il attend surtout ton accord.
\end{enumerate}
\end{methode}

\begin{exoresolu}[Réponses et intonation]
\textbf{Énoncé.} Répondez à chaque question tag selon l'indication entre parenthèses (la réalité), avec une réponse courte complète, puis indiquez l'intonation la plus probable du tag.

1. \eng{You didn't do your homework, did you?} (tu l'as fait)\\
2. \eng{Your father works in Douala, doesn't he?} (c'est vrai)\\
3. \eng{We aren't late, are we?} (vous êtes en retard)\\
4. \eng{She can't speak French, can she?} (elle le parle)\\
5. \eng{The teacher will give us a test, won't he?} (ce n'est pas sûr, tu demandes)

\tcblower
\textbf{Solution détaillée.}

1. La réalité est positive (tu as fait tes devoirs) : on contredit la question négative par \eng{Yes}. Réponse : \eng{Yes, I did.} (« Si, je les ai faits. ») Intonation probable : descendante ↘ (le locuteur suppose que tu ne les as pas faits et attend une confirmation).\\
2. La réalité est positive : \eng{Yes, he does.} Intonation descendante ↘ : simple confirmation.\\
3. La réalité est que vous \emph{êtes} en retard : le fait exprimé dans la phrase (« nous ne sommes pas en retard ») est donc \emph{faux}. On répond \eng{Yes, we are.} (« Si, nous sommes en retard. ») — et non \eng{No}, qui signifierait « nous ne sommes pas en retard ». Intonation montante ↗ possible si le locuteur doute vraiment.\\
4. La réalité est positive : \eng{Yes, she can.} (« Si, elle le parle. ») — jamais \eng{No, she can}.\\
5. Le locuteur n'est pas sûr : intonation montante ↗ sur \eng{won't he?}. Réponse possible : \eng{I don't know. He might.} (« Je ne sais pas. C'est possible. »)
\end{exoresolu}

\begin{attention}[Autres pièges à éviter]
\begin{itemize}
\item Ne répondez jamais avec le verbe principal seul : ✗ \eng{Yes, I like.} ✓ \eng{Yes, I do.}
\item Ne traduisez pas « si » par \eng{if} : ✗ \eng{If, I did!} ✓ \eng{Yes, I did!}
\item Gardez le même auxiliaire que dans le tag : \eng{You've eaten, haven't you? — Yes, I have} (pas \eng{Yes, I do}).
\item Si vous ne savez pas, dites-le honnêtement : \eng{I don't know.}, \eng{I'm not sure.}, \eng{Maybe.}
\end{itemize}
\end{attention}

\begin{experience}[Entraînement oral : les deux voix]
Par groupes de deux, chaque élève lit les paires suivantes deux fois : d'abord avec l'intonation descendante ↘ (confirmation), puis avec l'intonation montante ↗ (vraie question). Le partenaire doit dire, à chaque lecture, si le locuteur était sûr ou non, puis répondre avec une réponse courte complète.

1. \eng{You come from Buea, don't you?}\\
2. \eng{The exam is on Monday, isn't it?}\\
3. \eng{You haven't met the new principal, have you?}\\
4. \eng{We can play football after class, can't we?}\\
5. \eng{Your brother didn't pass, did he?}

Variante : le partenaire invente une réalité (vraie ou fausse) et répond en conséquence, en ajoutant un commentaire : \eng{Yes, I do! I was born there.}
\end{experience}

\begin{exercice}[À vous de jouer]
Répondez à chaque question tag selon la réalité indiquée, avec une réponse courte complète.

1. \eng{You don't speak English, do you?} (vous le parlez)\\
2. \eng{They haven't left yet, have they?} (ils sont partis)\\
3. \eng{It won't be difficult, will it?} (ce sera difficile)\\
4. \eng{You were at the market yesterday, weren't you?} (vous n'y étiez pas)\\
5. \eng{She can drive, can't she?} (elle ne sait pas conduire)
\end{exercice}

\begin{corrige}[Exercice : Réponses aux question tags]
1. \eng{Yes, I do.} (« Si, je le parle. »)\\
2. \eng{Yes, they have.} (« Si, ils sont partis. »)\\
3. \eng{Yes, it will.} (« Si, ce sera difficile. »)\\
4. \eng{No, I wasn't.} (« Non, je n'y étais pas. »)\\
5. \eng{No, she can't.} (« Non, elle ne sait pas conduire. »)

Rappel de méthode : identifiez d'abord la réalité (vraie ou fausse), choisissez \eng{Yes} ou \eng{No} en conséquence, puis reprenez l'auxiliaire du tag.
\end{corrige}

\begin{savaistu}[Pourquoi les Anglais utilisent-ils tant les question tags ?]
Dans la culture anglophone, le question tag à intonation descendante est un outil de \emph{politesse} et de \emph{conversation} : il invite l'autre à participer sans l'obliger à donner un avis tranché. Dire \eng{Lovely weather, isn't it? ↘} à un inconnu dans un bus à Londres est une manière classique d'engager la conversation. Au Cameroun anglophone, on entend souvent des variantes locales comme \eng{isn't it?} généralisé à tous les auxiliaires dans la langue parlée familière — mais à l'examen, exigez-vous la forme standard : \eng{You came, didn't you?}
\end{savaistu}

\coursec{Comparaison français/anglais et pièges des francophones}

Dans la section précédente, nous avons appris à \emph{répondre} à un question tag et à choisir la bonne intonation. Nous allons maintenant nous placer du point de vue du francophone que vous êtes : pourquoi les question tags sont-ils si difficiles pour nous ? Parce que le français résout le même problème avec \emph{un seul outil}, là où l'anglais en exige des dizaines. Cette section confronte les deux langues, puis passe en revue les cinq erreurs qui coûtent le plus de points au baccalauréat.

\courssub{Le choc des systèmes : un outil français contre une machine anglaise}

\begin{exemplebox}[Observez d'abord]
Comparez attentivement ces paires de phrases :

\eng{You are coming, aren't you?} « Tu viens, n'est-ce pas ? »

\eng{She isn't coming, is she?} « Elle ne vient pas, n'est-ce pas ? »

\eng{They played well, didn't they?} « Ils ont bien joué, n'est-ce pas ? »

\eng{He can swim, can't he?} « Il sait nager, n'est-ce pas ? »

Que remarquez-vous ? En français, la petite phrase finale est \textbf{toujours la même} : « n'est-ce pas ». En anglais, le tag \textbf{change à chaque phrase} : \eng{aren't you}, \eng{is she}, \eng{didn't they}, \eng{can't he}. L'anglais fabrique son tag sur mesure, comme un tailleur ; le français utilise un vêtement unique, taille universelle.
\end{exemplebox}

\begin{propriete}[La règle d'or de la comparaison]
\begin{itemize}
\item \textbf{Français} : « n'est-ce pas » est \cle{invariable}. Il s'accroche à n'importe quelle phrase, quel que soit le verbe, le temps, le sujet ou la polarité.
\item \textbf{Anglais} : le tag est \cle{variable}. Il dépend de trois informations contenues dans la phrase principale : l'\textbf{auxiliaire} (ou le verbe principal), le \textbf{sujet} (repris par un pronom) et la \textbf{polarité} (affirmative ou négative).
\end{itemize}
Autrement dit : chaque fois qu'un francophone veut dire « n'est-ce pas », il doit d'abord \emph{analyser sa propre phrase} avant de choisir le bon tag. C'est ce travail d'analyse que le français nous dispense de faire — et c'est exactement ce que le baccalauréat teste.
\end{propriete}

\begin{center}
\begin{tabularx}{\linewidth}{|X|X|X|}
\hline
\rowcolor{popblueL} Français & English & Remarque \\
\hline
Tu viens, n'est-ce pas ? & You are coming, aren't you? & tag avec \eng{be} au présent \\
\hline
Elle ne vient pas, n'est-ce pas ? & She isn't coming, is she? & polarité inversée \\
\hline
Ils ont bien joué, n'est-ce pas ? & They played well, didn't they? & passé : auxiliaire \eng{did} \\
\hline
Il sait nager, n'est-ce pas ? & He can swim, can't he? & modal \eng{can} repris \\
\hline
Vous êtes fatigués, n'est-ce pas ? & You are tired, aren't you? & pronom adapté au sujet \\
\hline
\end{tabularx}
\end{center}

\begin{popfigure}
\begin{tikzpicture}[
  box/.style={draw=popblue, thick, rounded corners, fill=popblueL, align=center, text width=5.2cm, inner sep=6pt},
  ebox/.style={draw=poporange, thick, rounded corners, fill=poporangeL, align=center, text width=5.2cm, inner sep=6pt},
  lab/.style={font=\bfseries, align=center},
  arr/.style={-{Stealth[length=3mm]}, thick, popdark}
]
\node[lab, popblue] at (-3,2.6) {Français : outil unique};
\node[lab, poporange] at (3,2.6) {Anglais : tag sur mesure};
\node[box, align=center] (f1) at (-3,1.2){« n'est-ce pas »\\ invariable};
\node[box, align=center] (f2) at (-3,-0.6){« si »\\ pour contredire une négation};
\node[ebox, align=center] (e1) at (3,1.2){auxiliaire + pronom\\ adaptés à la phrase\\\eng{aren't you? / didn't they?}};
\node[ebox, align=center] (e2) at (3,-0.6){\eng{Yes} + auxiliaire\\\eng{Yes, I do!}};
\draw[arr] (f1) -- (e1);
\draw[arr] (f2) -- (e2);
\end{tikzpicture}
\legende{Deux stratégies linguistiques pour les mêmes fonctions : le français fige, l'anglais adapte.}
\end{popfigure}

\courssub{Le cas de « si » : contredire une négation}

\begin{exemplebox}[Observez]
\eng{— You don't like ndolé, do you? — Yes, I do! I love it!}

« — Tu n'aimes pas le ndolé, n'est-ce pas ? — Si ! J'adore ça ! »
\end{exemplebox}

\begin{propriete}[« Si » français = \eng{Yes} + auxiliaire en anglais]
Quand on veut \textbf{contredire une question négative}, le français utilise le mot spécial « si ». L'anglais n'a pas de mot spécial : il utilise \cle{Yes} suivi de l'auxiliaire affirmative, souvent avec insistance.

\eng{— You haven't finished, have you? — Yes, I have!} « — Tu n'as pas fini, n'est-ce pas ? — Si, j'ai fini ! »

\eng{— She can't drive, can she? — Yes, she can!} « — Elle ne sait pas conduire, n'est-ce pas ? — Si, elle sait ! »

À l'inverse, pour confirmer la négation : \eng{No, I haven't.} « Non, je n'ai pas fini. »
\end{propriete}

\begin{attention}[Ne traduisez jamais « si » par \eng{yes} seul]
Le mot « si » de contradiction n'a \textbf{pas d'équivalent mot à mot} en anglais. Répondre simplement \eng{Yes} est possible mais ambigu ; la réponse complète et naturelle est \eng{Yes, I have!} / \eng{Yes, she can!}. À l'écrit du Bac, donnez toujours la réponse avec son auxiliaire.
\end{attention}

\courssub{Les cinq pièges majeurs des francophones}

\begin{attention}[Les 5 erreurs les plus fréquentes au Bac]
Voici, par ordre de fréquence dans les copies, les fautes à éradiquer absolument.

\textbf{Piège 1 — le tag universel \eng{isn't it}.} Sous l'influence de « n'est-ce pas », beaucoup d'élèves collent \eng{isn't it} partout.

✗ \eng{You are coming, isn't it?} \quad ✓ \eng{You are coming, aren't you?} « Tu viens, n'est-ce pas ? »

\textbf{Piège 2 — le nom répété au lieu du pronom.} Le sujet du tag est toujours un \textbf{pronom}, jamais un nom.

✗ \eng{Mary is kind, isn't Mary?} \quad ✓ \eng{Mary is kind, isn't she?} « Marie est gentille, n'est-ce pas ? »

\textbf{Piège 3 — \eng{do/does/did} oublié.} Quand la phrase n'a pas d'auxiliaire (present simple ou past simple), on ne peut pas contracter le verbe lexical : on recrute \eng{do}, \eng{does} ou \eng{did}.

✗ \eng{You play football, playn't you?} \quad ✓ \eng{You play football, don't you?} « Tu joues au football, n'est-ce pas ? »

\textbf{Piège 4 — la polarité non inversée.} Phrase affirmative → tag négatif ; phrase négative → tag affirmatif. Ne jamais additionner deux négations.

✗ \eng{He came, didn't he not?} \quad ✓ \eng{He came, didn't he?}

✗ \eng{She is ill, is she?} \quad ✓ \eng{She is ill, isn't she?} « Elle est malade, n'est-ce pas ? »

\textbf{Piège 5 — le mauvais temps dans le tag.} Le tag doit reprendre le temps (et l'auxiliaire) de la phrase principale.

✗ \eng{They went home, don't they?} \quad ✓ \eng{They went home, didn't they?} « Ils sont rentrés, n'est-ce pas ? »
\end{attention}

\begin{attention}[La contraction est obligatoire à l'usage courant]
À l'oral comme à l'écrit courant, le tag négatif est \textbf{toujours contracté} : on dit et on écrit \eng{isn't she}, \eng{don't you}, \eng{can't he}, et non \eng{is not she} ou \eng{do not you}, qui sonnent artificiels ou solennels.

✗ \eng{She is your sister, is not she?} \quad ✓ \eng{She is your sister, isn't she?}

Seule exception connue : \eng{am I not} à la place de \eng{aren't I} dans un style très formel (\eng{I am late, am I not?}), mais \eng{aren't I} reste la forme attendue au lycée.
\end{attention}

\begin{savaistu}[Les « same tags » : l'exception qui confirme la règle]
Il existe des tags de \textbf{même polarité} : \eng{You like him, do you?} ou \eng{Oh, you've failed, have you?}. Ils n'expriment pas une demande de confirmation mais l'\textbf{ironie}, la \textbf{surprise} ou le \textbf{doute} : « Ah bon, tu l'aimes, vraiment ? ». Ces \emph{same tags} sont rares et hors usage courant : au baccalauréat, il faut savoir les \emph{reconnaître} à la compréhension, mais ne jamais les \emph{produire} à la place d'un tag ordinaire. Votre règle de production reste : polarité inversée, toujours.
\end{savaistu}

\courssub{Méthode de vérification à l'écrit}

\begin{methode}[Trois gestes avant de valider un tag]
Au brouillon comme en composition, prenez l'habitude de ce contrôle systématique :

\begin{enumerate}
\item \textbf{Soulignez l'auxiliaire} de la phrase principale (\eng{is}, \eng{can}, \eng{have}, \eng{will}...). S'il n'y en a pas, c'est un present simple ou un past simple : recrutez \eng{do}, \eng{does} ou \eng{did}.
\item \textbf{Entourez le sujet} et remplacez-le par le pronom correspondant : \eng{Mary} → \eng{she}, \eng{the boys} → \eng{they}, \eng{my father and I} → \eng{we}.
\item \textbf{Vérifiez la polarité} : phrase affirmative → tag négatif contracté ; phrase négative → tag affirmatif.
\end{enumerate}

Exemple appliqué : \eng{The students finished early, \ldots ?} → pas d'auxiliaire, past simple → \eng{did} ; sujet \eng{the students} → \eng{they} ; phrase affirmative → négatif. Résultat : \eng{The students finished early, didn't they?}
\end{methode}

\begin{exoresolu}[Corrigez les erreurs d'un élève de Tle A]
Énoncé — Voici cinq tags écrits par un élève francophone. Corrigez chaque erreur et justifiez.

1. \eng{Your brother works in Douala, isn't it?}

2. \eng{Ngono speaks three languages, doesn't Ngono?}

3. \eng{You enjoyed the match, enjoyn't you?}

4. \eng{She hasn't called, hasn't she?}

5. \eng{They left Yaoundé at six, don't they?}

\tcblower
Solution détaillée.

1. ✗ \eng{isn't it} → ✓ \eng{Your brother works in Douala, doesn't he?} Le verbe \eng{works} est au present simple sans auxiliaire : on recrute \eng{does} ; le sujet masculin devient \eng{he}. Le tag universel \eng{isn't it} est la traduction fautive de « n'est-ce pas ».

2. ✗ \eng{doesn't Ngono} → ✓ \eng{Ngono speaks three languages, doesn't she?} Le sujet du tag est obligatoirement un pronom personnel, jamais le nom répété.

3. ✗ \eng{enjoyn't you} → ✓ \eng{You enjoyed the match, didn't you?} On ne contracte jamais un verbe lexical ; le past simple exige l'auxiliaire \eng{did}.

4. ✗ \eng{hasn't she} → ✓ \eng{She hasn't called, has she?} La phrase est déjà négative (\eng{hasn't}) : le tag doit être affirmatif. Deux négations ne s'additionnent pas.

5. ✗ \eng{don't they} → ✓ \eng{They left Yaoundé at six, didn't they?} Le tag reprend le temps de la phrase : \eng{left} est au past simple, donc \eng{didn't}, pas \eng{don't}.
\end{exoresolu}

\begin{exercice}[À vous de jouer]
Ajoutez le bon question tag, puis traduisez chaque phrase en français.

1. \eng{Your parents live in Bamenda, \ldots ?}

2. \eng{The bendskin driver wasn't careful, \ldots ?}

3. \eng{We will pass the baccalauréat, \ldots ?}

4. \eng{She doesn't like cocoyams, \ldots ?}

5. \eng{The students have finished their composition, \ldots ?}
\end{exercice}

\begin{corrige}[Exercice « À vous de jouer »]
1. \eng{Your parents live in Bamenda, don't they?} « Tes parents habitent à Bamenda, n'est-ce pas ? » (present simple sans auxiliaire → \eng{do}).

2. \eng{The bendskin driver wasn't careful, was he?} « Le conducteur de bendskin n'était pas prudent, n'est-ce pas ? » (phrase négative → tag affirmatif).

3. \eng{We will pass the baccalauréat, won't we?} « Nous réussirons au baccalauréat, n'est-ce pas ? » (modal \eng{will} repris, contracté en \eng{won't}).

4. \eng{She doesn't like cocoyams, does she?} « Elle n'aime pas les macabos, n'est-ce pas ? » (négative → affirmative).

5. \eng{The students have finished their composition, haven't they?} « Les élèves ont fini leur rédaction, n'est-ce pas ? » (present perfect : auxiliaire \eng{have} repris).
\end{corrige}

\begin{savaistu}[Pourquoi les anglophones ne peuvent pas s'en passer]
Les anglophones utilisent les question tags \textbf{en permanence} : non pas seulement pour vérifier une information, mais pour maintenir la conversation, inviter l'autre à parler et adoucir une affirmation. \eng{Lovely weather, isn't it?} n'attend pas vraiment une réponse météorologique : c'est une porte ouverte au dialogue. À l'inverse, un anglais parlé sans aucun tag sonne souvent \textbf{froid, sec ou trop direct}, presque impoli. Maîtriser les tags, c'est donc maîtriser aussi la politesse conversationnelle anglaise — une compétence que les correcteurs valorisent dans les productions écrites dialoguées.
\end{savaistu}

\begin{aretenir}[L'essentiel de la section]
\begin{itemize}
\item « N'est-ce pas » est invariable ; le tag anglais est toujours adapté : auxiliaire + pronom + polarité inversée.
\item « Si » de contradiction se rend par \eng{Yes} + auxiliaire : \eng{Yes, I do!}
\item Jamais de nom dans le tag, jamais de verbe lexical contracté, jamais de double négation, jamais de temps différent.
\item Avant de valider : souligner l'auxiliaire, entourer le sujet, vérifier la polarité.
\end{itemize}
\end{aretenir}

\coursec{Entraînement guidé et production}

Nous savons désormais former les \cle{question tags}, traiter les cas particuliers et éviter les pièges classiques des francophones (le fameux \eng{isn’t it?} universel, calque de « n’est-ce pas »). Il reste l’étape décisive : passer de la règle à l’automatisme. Cette dernière partie propose un entraînement progressif, de la reconnaissance à la production libre, dans l’esprit des épreuves du Baccalauréat de la série A.

\begin{popfigure}
\begin{center}
\begin{tikzpicture}[font=\small]
\node[draw=popblue, fill=popblueL, rounded corners, align=center, minimum width=2.6cm, minimum height=1.1cm] (a) at (0,0) {\textbf{Observer}\\ repérer les tags};
\node[draw=poporange, fill=poporangeL, rounded corners, align=center, minimum width=2.6cm, minimum height=1.1cm, right=0.5cm of a] (b) {\textbf{Former}\\ compléter};
\node[draw=popgreen, fill=popgreenL, rounded corners, align=center, minimum width=2.6cm, minimum height=1.1cm, right=0.5cm of b] (c) {\textbf{Exceptions}\\ transformer};
\node[draw=poppurple, fill=poppurpleL, rounded corners, align=center, minimum width=2.6cm, minimum height=1.1cm, right=0.5cm of c] (d) {\textbf{Répondre}\\ Yes / No + aux.};
\node[draw=poppink, fill=poppinkL, rounded corners, align=center, minimum width=2.6cm, minimum height=1.1cm, right=0.5cm of d] (e) {\textbf{Produire}\\ dialogue libre};
\draw[-{Stealth}, thick, popmuted] (a) -- (b);
\draw[-{Stealth}, thick, popmuted] (b) -- (c);
\draw[-{Stealth}, thick, popmuted] (c) -- (d);
\draw[-{Stealth}, thick, popmuted] (d) -- (e);
\end{tikzpicture}
\end{center}
\legende{Frise de progression de l’entraînement : de l’observation à la production.}
\end{popfigure}

\courssub{Exercice 1 : reconnaissance des tags dans un dialogue}

\begin{methode}[Repérer un question tag dans un texte]
\begin{enumerate}
\item Cherche d’abord les \textbf{virgules suivies d’un point d’interrogation} : le tag se trouve presque toujours entre la virgule et le « ? » final.
\item Vérifie que le fragment trouvé contient bien \textbf{un auxiliaire + un pronom} (et non un sujet nominal).
\item Contrôle la \textbf{polarité} : phrase affirmative $\rightarrow$ tag négatif ; phrase négative $\rightarrow$ tag affirmatif.
\item Signale les tags fautifs et propose une correction.
\end{enumerate}
\end{methode}

\begin{exercice}[Souligne les tags et vérifie-les]
Dans le dialogue suivant entre deux élèves de Terminale A à Yaoundé, souligne les cinq question tags, puis dis s’ils sont corrects. Corrige ceux qui sont fautifs.\par
\eng{— You are going to the school library after class, aren’t you?\\
— Yes, I am. The new English novels have arrived, isn’t it?\\
— Really? Let’s go together, will we?\\
— Wait! You haven’t finished your homework, haven’t you?\\
— Don’t worry, I did it this morning. Mr Etoundi never checks it, doesn’t he?}
\end{exercice}

\begin{corrige}[Exercice 1]
\begin{itemize}
\item \eng{aren’t you?} : correct (auxiliaire \eng{are}, polarité inversée, pronom \eng{you}).
\item \eng{isn’t it?} : \textbf{faux}. Le sujet est \eng{the new English novels} (pluriel) et l’auxiliaire \eng{have} : \eng{The new English novels have arrived, haven’t they?} « Les nouveaux romans anglais sont arrivés, n’est-ce pas ? »
\item \eng{will we?} : \textbf{faux}. Après \eng{Let’s}, le tag est toujours \eng{shall we?} : \eng{Let’s go together, shall we?}
\item \eng{haven’t you?} : \textbf{faux}. Phrase négative $\rightarrow$ tag affirmatif : \eng{You haven’t finished your homework, have you?}
\item \eng{doesn’t he?} : correct (présent simple, 3\up{e} personne, polarité inversée).
\end{itemize}
\end{corrige}

\courssub{Exercice 2 : complétion (cas de base et exceptions mélangés)}

\begin{methode}[Réussir l’exercice de complétion au Bac]
Pour chaque phrase, applique méthodiquement les quatre étapes :
\begin{enumerate}
\item \textbf{Repère le temps et l’auxiliaire} de la phrase principale (s’il n’y a pas d’auxiliaire visible, utilise \eng{do / does / did}).
\item \textbf{Vérifie la polarité} : affirmative $\rightarrow$ tag négatif ; négative $\rightarrow$ tag affirmatif.
\item \textbf{Choisis le pronom} qui remplace exactement le sujet (attention à \eng{there}, \eng{nobody}, \eng{this}).
\item \textbf{Contracte} le tag négatif : \eng{isn’t}, \eng{don’t}, \eng{won’t} (et non \eng{willn’t} !).
\end{enumerate}
\end{methode}

\begin{exemplebox}[Deux modèles à imiter]
\eng{The students haven’t arrived, have they?} « Les élèves ne sont pas arrivés, si ? »\par
\eng{Close the window, will you?} « Ferme la fenêtre, veux-tu ? »
\end{exemplebox}

\begin{exercice}[Complète avec le tag approprié]
\begin{enumerate}
\item Your father works at the hospital in Douala, \ldots ?
\item The rain won’t stop before evening, \ldots ?
\item Let’s watch the football match tonight, \ldots ?
\item I am late for the lesson, \ldots ?
\item Nobody phoned while I was out, \ldots ?
\item She has never visited the Bamenda highlands, \ldots ?
\item Open your books, \ldots ?
\item There is a cybercafé near the market, \ldots ?
\item The students had already left when the bell rang, \ldots ?
\item You’d rather stay at home, \ldots ?
\end{enumerate}
\end{exercice}

\begin{corrige}[Exercice 2]
\begin{enumerate}
\item \eng{doesn’t he} (présent simple, 3\up{e} pers. sing., affirmatif $\rightarrow$ négatif).
\item \eng{will it} (négatif \eng{won’t} $\rightarrow$ tag affirmatif ; \eng{the rain} = \eng{it}).
\item \eng{shall we} (exception : après \eng{Let’s}, toujours \eng{shall we}).
\item \eng{aren’t I} (exception : \eng{I am} $\rightarrow$ \eng{aren’t I}, jamais \eng{amn’t I}).
\item \eng{did they} (\eng{nobody} est négatif $\rightarrow$ tag affirmatif ; pronom \eng{they}).
\item \eng{has she} (\eng{never} rend la phrase négative $\rightarrow$ tag affirmatif).
\item \eng{will you} (impératif $\rightarrow$ \eng{will you?} / \eng{would you?}).
\item \eng{isn’t there} (avec \eng{there is/are}, le tag reprend \eng{there}).
\item \eng{hadn’t they} (plus-que-parfait, auxiliaire \eng{had}, polarité inversée).
\item \eng{wouldn’t you} (\eng{’d rather} = \eng{would rather}).
\end{enumerate}
\end{corrige}

\begin{attention}[Ponctuation obligatoire à l’écrit]
À l’écrit, deux fautes coûtent des points au Bac : l’oubli de la \textbf{virgule} avant le tag et l’oubli du \textbf{point d’interrogation} final. Écris toujours : \eng{She sings well, doesn’t she?} — et jamais \eng{She sings well doesn’t she.} La virgule marque la pause, le « ? » marque l’attente d’une réponse.
\end{attention}

\courssub{Exercice 3 : transformation — du français « n’est-ce pas » au tag anglais}

Le français utilise une seule formule invariable ; l’anglais oblige à adapter le tag à chaque phrase. C’est exactement le type d’exercice de transformation proposé à l’examen.

\begin{exoresolu}[Traduire « n’est-ce pas »]
Énoncé : traduis en anglais avec un question tag : « Ta sœur parle bien anglais, n’est-ce pas ? »\par
\tcblower
Solution détaillée :
\begin{enumerate}
\item Traduction de la phrase : \eng{Your sister speaks English well.}
\item Temps : présent simple, 3\up{e} personne $\rightarrow$ auxiliaire \eng{does}.
\item Polarité : phrase affirmative $\rightarrow$ tag négatif : \eng{doesn’t}.
\item Pronom : \eng{your sister} $\rightarrow$ \eng{she}.
\item Résultat : \eng{Your sister speaks English well, doesn’t she?}
\end{enumerate}
Le piège était de traduire « n’est-ce pas » par \eng{isn’t it?}, qui ne convient que si la phrase contient \eng{it is}.
\end{exoresolu}

\begin{exercice}[Transforme les phrases]
Traduis chaque phrase française en anglais en utilisant un question tag correct.
\begin{enumerate}
\item Le match commence à seize heures, n’est-ce pas ?
\item Tu n’as pas vu mon cahier, si ?
\item Allons au marché de Mokolo, n’est-ce pas ?
\item Elle était fatiguée hier soir, n’est-ce pas ?
\item Ils reviendront demain, n’est-ce pas ?
\end{enumerate}
\end{exercice}

\begin{corrige}[Exercice 3]
\begin{enumerate}
\item \eng{The match starts at 4 p.m., doesn’t it?}
\item \eng{You haven’t seen my notebook, have you?}
\item \eng{Let’s go to the Mokolo market, shall we?}
\item \eng{She was tired last night, wasn’t she?}
\item \eng{They will come back tomorrow, won’t they?}
\end{enumerate}
\end{corrige}

\courssub{Exercice 4 : correction d’erreurs (type Bac)}

\begin{exercice}[Corrige les tags fautifs]
Chaque phrase ci-dessous, produite par un élève francophone, contient un tag incorrect. Souligne l’erreur et réécris le tag correctement.
\begin{enumerate}
\item \eng{You are a student, are you?}
\item \eng{She doesn’t like fufu, doesn’t she?}
\item \eng{They went to Buea last year, wentn’t they?}
\item \eng{I am your best friend, amn’t I?}
\item \eng{Let’s revise the lesson, will we?}
\item \eng{Everything is ready, isn’t they?}
\item \eng{He can swim very well, can he not?}
\item \eng{Your brother lives in Bamenda, isn’t it?}
\end{enumerate}
\end{exercice}

\begin{corrige}[Exercice 4]
\begin{enumerate}
\item \eng{aren’t you} — la phrase est affirmative, le tag doit être négatif.
\item \eng{does she} — phrase négative $\rightarrow$ tag affirmatif.
\item \eng{didn’t they} — au prétérit simple sans auxiliaire, on utilise \eng{did} ; \eng{wentn’t} n’existe pas.
\item \eng{aren’t I} — forme exceptionnelle obligatoire avec \eng{I am}.
\item \eng{shall we} — après \eng{Let’s}, toujours \eng{shall we}.
\item \eng{isn’t it} — \eng{everything} se remplace par \eng{it}, pas \eng{they}.
\item \eng{can’t he} — à l’écrit comme à l’oral, on contracte : \eng{can he not} est lourd et artificiel.
\item \eng{doesn’t he} — le verbe est \eng{lives} (présent simple), pas \eng{is} ; et le sujet est \eng{he}.
\end{enumerate}
\end{corrige}

\courssub{Exercice 5 : répondre par des réponses courtes}

\begin{aretenir}[Répondre à un question tag]
En anglais, on répond au \textbf{fait}, pas à la forme de la question. La réponse courte reprend l’auxiliaire du tag : \eng{Yes, I do. / No, she hasn’t.} Attention au piège francophone : à une phrase négative, le français répond « si », l’anglais répond \eng{Yes}. \eng{You aren’t tired, are you? — Yes, I am!} « Tu n’es pas fatigué, si ? — Si ! »
\end{aretenir}

\begin{exercice}[Réponds par une réponse courte]
Réponds à chaque question par \eng{Yes} ou \eng{No} + sujet + auxiliaire, selon l’indication entre parenthèses.
\begin{enumerate}
\item \eng{You have done your homework, haven’t you?} (oui)
\item \eng{It isn’t raining now, is it?} (non)
\item \eng{Your mother cooks eru very well, doesn’t she?} (oui)
\item \eng{You didn’t watch the match yesterday, did you?} (si / oui)
\item \eng{We shall meet at the bus station, shan’t we?} (oui)
\end{enumerate}
\end{exercice}

\begin{corrige}[Exercice 5]
\begin{enumerate}
\item \eng{Yes, I have.}
\item \eng{No, it isn’t.}
\item \eng{Yes, she does.}
\item \eng{Yes, I did.} (« si » français = \eng{Yes} anglais)
\item \eng{Yes, we shall.}
\end{enumerate}
\end{corrige}

\courssub{Production libre : écrire un dialogue}

\begin{propriete}[Grille de réussite du dialogue final]
Ton dialogue est réussi s’il respecte les quatre critères suivants :
\begin{itemize}
\item au moins \textbf{5 question tags} correctement formés et ponctués (virgule + « ? ») ;
\item au moins \textbf{3 auxiliaires différents} (par exemple \eng{do}, \eng{have}, \eng{will}, \eng{can}, \eng{be}) ;
\item au moins \textbf{1 exception} : \eng{Let’s \ldots, shall we?}, un impératif avec \eng{will you?} ou \eng{I am \ldots, aren’t I?} ;
\item au moins \textbf{2 réponses courtes} correctes (\eng{Yes, I do. / No, she can’t.}).
\end{itemize}
\end{propriete}

\begin{exemplebox}[Modèle de dialogue (8 répliques)]
\eng{— Hi Amina! You are going to the end-of-year party, aren’t you?\\
— Yes, I am. It starts at six o’clock, doesn’t it?\\
— That’s right. You haven’t invited your cousin yet, have you?\\
— No, I haven’t. I am a bit shy about it, aren’t I?\\
— Don’t worry! Let’s invite her together, shall we?\\
— Good idea. Pass me my phone, will you?\\
— Here it is. She can bring her friends too, can’t she?\\
— Yes, she can. The more, the merrier!}
\end{exemplebox}

\begin{experience}[À toi d’écrire et de jouer]
Rédige un dialogue de \textbf{8 répliques} entre deux amis qui préparent une excursion au mont Cameroun ou une fête de famille. Vérifie ton texte avec la grille de réussite ci-dessus, échange-le avec ton voisin pour une correction croisée, puis jouez la scène devant la classe en soignant l’intonation : voix qui \emph{monte} si tu attends vraiment une réponse, voix qui \emph{descend} si tu attends simplement l’accord de ton interlocuteur.
\end{experience}

\begin{center}
\begin{tabularx}{\linewidth}{|X|X|X|}
\hline
\rowcolor{popblueL} Français & English & Remarque \\
\hline
Tu viens, n’est-ce pas ? & You are coming, aren’t you? & tag adapté à l’auxiliaire \eng{be} \\
\hline
Il ne pleut pas, si ? & It isn’t raining, is it? & « si » = tag affirmatif \\
\hline
Allons-y, n’est-ce pas ? & Let’s go, shall we? & exception : \eng{shall we} \\
\hline
Ferme la porte, veux-tu ? & Close the door, will you? & impératif + \eng{will you} \\
\hline
Je suis en retard, non ? & I am late, aren’t I? & exception : \eng{aren’t I} \\
\hline
Personne n’a téléphoné, si ? & Nobody phoned, did they? & \eng{nobody} = négatif ; pronom \eng{they} \\
\hline
\end{tabularx}
\end{center}

\begin{attention}[Derniers pièges avant l’examen]
\begin{itemize}
\item Ne traduis jamais « n’est-ce pas » mécaniquement par \eng{isn’t it?} : interroge-toi d’abord sur le temps, l’auxiliaire et le sujet.
\item N’oublie ni la virgule ni le point d’interrogation : \eng{She sings well, doesn’t she?}
\item N’écris jamais \eng{willn’t} : la contraction de \eng{will not} est \eng{won’t}.
\item Une réponse courte ne répète jamais le verbe principal : \eng{Yes, I do.} et non \eng{Yes, I like.}
\end{itemize}
\end{attention}

Tu es maintenant capable de reconnaître, former, corriger et employer les question tags dans un échange naturel. À l’oral comme à l’écrit, cette petite structure de trois mots donne immédiatement à ton anglais un ton authentique et interactif : c’est la marque d’un locuteur qui dialogue, et non qui récite.

\newpage

\coursec{Activité d'intégration}

\coursec{Lesson 8 — Grammar: Revise question tags}

\courssub{Objectifs de la leçon}

À l'issue de cette leçon, tu seras capable de :

\begin{itemize}
\item identifier et construire correctement la structure du \cle{question tag} (auxiliaire, pronom, polarité) ;
\item distinguer les cas particuliers (\eng{I am}, \eng{Let's}, impératif, \eng{there is/are}, indéfinis \eng{nobody/everybody}) ;
\item répondre par une \cle{réponse courte} adaptée aux faits et non à la forme de la question ;
\item utiliser l'intonation montante (doute réel) ou descendante (confirmation attendue) de façon pertinente ;
\item produire un court texte écrit (article, interview) intégrant des question tags variés et corrects ;
\item corriger les erreurs fréquentes des francophones (tag universel \eng{isn't it}, inversion auxiliaire/verbe lexical, réponse \eng{Yes/No} inversée).
\end{itemize}

\courssub{Observation et découverte : situation de communication authentique}

Le club d'anglais de ton lycée, le \eng{English Speakers' Club} du Lycée de Nkol-Eton à Yaoundé, prépare son journal mensuel, \emph{The Lion's Voice}. La rédactrice en chef, Mme Abena, lance un défi aux élèves de Terminale A : chaque binôme doit réaliser une \cle{interview croisée} pour découvrir le profil d'un camarade, en utilisant un maximum de question tags. Les meilleurs articles seront publiés dans la rubrique \eng{Getting to know you}.

\begin{textebox}[Secret identity card — Sample (English Speakers' Club)]
SECRET IDENTITY CARD — Do not show it to your partner!

Name: Chantal Ewusi

Home town: Buea (South-West Region)

Favourite subject: Biology

Languages spoken: English, French and Pidgin

A hidden talent: I can dance the makossa very well.

A place I have visited: I have been to Limbe beach twice.

A dream: I would like to study medicine at the University of Buea.

Something I have never done: I have never travelled abroad.

Something I dislike: I can't stand waking up before 5 a.m.

A fact about my family: My elder brother plays football for a local club.
\end{textebox}

\begin{textebox}[Extract from the club's notice board]
THE LION'S VOICE — October issue

Wanted: the best pair of interviewers in Terminale A!

Rules of the game:
1. Work in pairs. Each student receives a secret identity card.
2. Guess your partner's information using question tags only. Example: “You come from Douala, don't you?”
3. Answer with a correct short answer: “Yes, I do.” or “No, I don't.”
4. Write down every piece of information you confirm.
5. Present your partner to the class in five sentences, three of them with question tags.
6. The class listens and checks the tags. Any mistake = one point lost!

Deadline: Friday, 4 p.m. Good luck, reporters!
\end{textebox}

\courssub{Règle de formation : le principe de base}

\begin{propriete}[Structure fondamentale du question tag]
Un \cle{question tag} (ou \eng{tag question}) est une mini-question ajoutée à la fin d'une phrase déclarative pour demander confirmation ou accord. Sa structure est immuable :

\begin{center}
\textbf{Phrase déclarative + virgule + auxiliaire (polarité inversée) + pronom sujet + point d'interrogation}
\end{center}

\begin{center}
\begin{tabularx}{\linewidth}{|X|X|X|}
\hline
\rowcolor{popblueL} \textbf{Type de phrase} & \textbf{Auxiliaire du tag} & \textbf{Exemple} \\
\hline
Affirmative & Négatif (contracté \eng{n't}) & \eng{You like fufu, \textbf{don't you?}} \\
\hline
Négative & Affirmatif & \eng{You don't like fufu, \textbf{do you?}} \\
\hline
\end{tabularx}
\end{center}

\textbf{Règles d'or} :
\begin{enumerate}
\item \textbf{L'auxiliaire} : on reprend l'auxiliaire de la phrase (ou \eng{do/does/did} s'il n'y en a pas). Le verbe lexical ne sert jamais d'auxiliaire au tag.
\item \textbf{La polarité} : phrase positive $\rightarrow$ tag négatif ; phrase négative $\rightarrow$ tag positif.
\item \textbf{Le pronom} : toujours un pronom personnel sujet (\eng{I, you, he, she, it, we, they}), jamais un nom. Le pronom correspond au sujet de la phrase.
\item \textbf{La contraction} : au tag négatif, la forme contractée (\eng{isn't, don't, hasn't, won't, can't}) est la norme à l'oral et à l'écrit courant.
\end{enumerate}
\end{propriete}

\begin{exemplebox}[Exemples de base (présent, prétérit, perfect, modaux)]
\begin{itemize}
\item \eng{She lives in Douala, \textbf{doesn't she?}} (Présent simple, verbe lexical $\rightarrow$ \eng{does})
\item \eng{They didn't come yesterday, \textbf{did they?}} (Prétérit simple négatif $\rightarrow$ \eng{did})
\item \eng{We have finished the work, \textbf{haven't we?}} (Present perfect $\rightarrow$ \eng{have})
\item \eng{He can speak Pidgin, \textbf{can't he?}} (Modal \eng{can} $\rightarrow$ \eng{can't})
\item \eng{It will rain tomorrow, \textbf{won't it?}} (Futur \eng{will} $\rightarrow$ \eng{won't} — irrégulier)
\end{itemize}
\end{exemplebox}

\courssub{Cas particuliers et exceptions}

\begin{propriete}[Les exceptions incontournables pour le Bac]
Ces cas ne suivent pas la règle mécanique ; ils doivent être appris par cœur.
\end{propriete}

\begin{center}
\begin{tabularx}{\linewidth}{|X|X|X|}
\hline
\rowcolor{popblueL} \textbf{Cas particulier} & \textbf{Règle / Tag} & \textbf{Exemple traduit} \\
\hline
\eng{I am} (affirmatif) & Tag = \eng{aren't I?} (pas \eng{amn't I}) & \eng{I am late, \textbf{aren't I?}} « Je suis en retard, n'est-ce pas ? » \\
\hline
\eng{I am not} (négatif) & Tag = \eng{am I?} & \eng{I am not late, \textbf{am I?}} « Je ne suis pas en retard, si ? » \\
\hline
\eng{Let's} (suggestion) & Tag = \eng{shall we?} & \eng{Let's go to the maquis, \textbf{shall we?}} « On va au maquis, d'accord ? » \\
\hline
Impératif (ordre/conseil) & Tag = \eng{will you?} / \eng{would you?} / \eng{can you?} / \eng{could you?} & \eng{Open the window, \textbf{will you?}} « Ouvre la fenêtre, veux-tu ? » \\
\hline
\eng{There is / are} & Tag = \eng{isn't there?} / \eng{aren't there?} & \eng{There is a problem, \textbf{isn't there?}} « Il y a un problème, non ? » \\
\hline
Sujet \eng{nobody / no one / nothing} & Phrase négative $\rightarrow$ tag \textbf{affirmatif} + pronom \eng{they} & \eng{Nobody came, \textbf{did they?}} « Personne n'est venu, n'est-ce pas ? » \\
\hline
Sujet \eng{everybody / everyone / everything} & Tag = pronom \eng{they} & \eng{Everyone knows him, \textbf{don't they?}} « Tout le monde le connaît, non ? » \\
\hline
\eng{This / That is} & Tag = \eng{isn't it?} & \eng{This is your pen, \textbf{isn't it?}} « C'est ton stylo, non ? » \\
\hline
\eng{There is / are} avec indéfini & Tag selon le verbe \eng{be} & \eng{There are some mangoes, \textbf{aren't there?}} \\
\hline
\end{tabularx}
\end{center}

\begin{attention}[Piège : \eng{have got} vs \eng{have} (action/possession)]
\begin{itemize}
\item \eng{Have got} (possession, britannique) $\rightarrow$ tag sur \eng{have} : \eng{She has got a car, \textbf{hasn't she?}}
\item \eng{Have} comme verbe lexical (action : \eng{have breakfast}, \eng{have a shower}) $\rightarrow$ tag sur \eng{do/does/did} : \eng{She has breakfast at 7, \textbf{doesn't she?}} \textbf{(Pas \eng{hasn't she?})}
\end{itemize}
\end{attention}

\courssub{Répondre à un question tag et intonation}

\begin{propriete}[La réponse courte : logique des faits, pas de la forme]
En anglais, on répond \textbf{selon la réalité} (les faits), pas selon la polarité de la question. C'est la différence majeure avec le français (\eng{Si} / \eng{Non}).

\begin{center}
\begin{tabularx}{\linewidth}{|X|X|X|}
\hline
\rowcolor{popblueL} \textbf{Question tag} & \textbf{Réalité = OUI} & \textbf{Réalité = NON} \\
\hline
\eng{You are Cameroonian, aren't you?} & \eng{\textbf{Yes, I am.}} & \eng{\textbf{No, I'm not.}} \\
\hline
\eng{You aren't Cameroonian, are you?} & \eng{\textbf{Yes, I am.}} & \eng{\textbf{No, I'm not.}} \\
\hline
\eng{You like eru, don't you?} & \eng{\textbf{Yes, I do.}} & \eng{\textbf{No, I don't.}} \\
\hline
\eng{You don't like eru, do you?} & \eng{\textbf{Yes, I do.}} & \eng{\textbf{No, I don't.}} \\
\hline
\end{tabularx}
\end{center}

\textbf{Astuce} : Pour ne pas te tromper, réponds mentalement en phrase complète : \eng{“Yes, I \textbf{am} Cameroonian”} ou \eng{“No, I \textbf{don't} like eru”}, puis garde l'auxiliaire.
\end{propriete}

\begin{propriete}[Intonation : le sens change]
\begin{itemize}
\item \textbf{Intonation descendante} (\searrow) : le locuteur est quasi certain, il cherche juste l'accord. \eng{“It's a beautiful day, isn't it?”} (voix qui tombe sur \eng{isn't it}).
\item \textbf{Intonation montante} (\nearrow) : le locuteur a un doute réel, il demande une vraie information. \eng{“You haven't seen my keys, have you?”} (voix qui monte sur \eng{have you}).
\end{itemize}
\end{propriete}

\courssub{Comparaison français / anglais et pièges des francophones}

\begin{center}
\begin{tabularx}{\linewidth}{|X|X|X|}
\hline
\rowcolor{popblueL} \textbf{Français} & \textbf{English} & \textbf{Remarque / Erreur fréquente} \\
\hline
N'est-ce pas ? / Hein ? (invariant) & Tag variable selon auxiliaire, temps, personne & \textbf{Erreur} : \eng{*You come from Buea, isn't it?} $\rightarrow$ \eng{You come from Buea, \textbf{don't you?}} \\
\hline
“Non, je n'aime pas.” (négation) & \eng{No, I don't.} (auxiliaire positif) & \textbf{Erreur} : \eng{*No, I don't like.} $\rightarrow$ Réponse courte = auxiliaire seulement. \\
\hline
“Si, j'aime.” (contredire une négation) & \eng{Yes, I do.} & \textbf{Erreur} : \eng{*Yes, I like.} ou confusion \eng{Yes/No}. \\
\hline
“Tu as faim, non ?” (\eng{avoir} = état) & \eng{You are hungry, aren't you?} & \textbf{Faux ami} : \eng{have} (possession/état) $\neq$ \eng{have got} / \eng{be}. \eng{*You have hunger, don't you?} n'existe pas. \\
\hline
“On y va ?” & \eng{Let's go, shall we?} & \textbf{Exception} : \eng{Let's} $\rightarrow$ \eng{shall we?} (pas \eng{let we?}). \\
\hline
\end{tabularx}
\end{center}

\begin{attention}[Check-list des 5 erreurs fatales au Bac]
\begin{enumerate}
\item \textbf{Tag universel} : Utiliser \eng{isn't it?} ou \eng{no?} pour tout. $\rightarrow$ Identifie l'auxiliaire de la phrase !
\item \textbf{Mauvais auxiliaire} : \eng{*She has visited Limbe, didn't she?} $\rightarrow$ \eng{hasn't she?} (Present perfect = \eng{have}).
\item \textbf{Mauvais pronom} : \eng{*My brother plays football, doesn't he?} $\rightarrow$ Correct. Mais \eng{*My brothers plays football, doesn't he?} $\rightarrow$ \eng{don't they?} (accord sujet).
\item \textbf{Polarité non inversée} : \eng{*She is nice, is she?} (sauf ironie/suspicion) $\rightarrow$ \eng{isn't she?}
\item \textbf{Réponse courte incomplète} : \eng{*Yes, I am Cameroonian.} $\rightarrow$ \eng{Yes, I am.} (Pas de répétition du complément).
\end{enumerate}
\end{attention}

\courssub{Entraînement guidé et production : Jeu d'interview « Getting to know you »}

\begin{experience}[Interview croisée — The Lion's Voice]
\textbf{Matériel} : Une fiche d'identité secrète par élève (préparée par l'enseignant ou inventée par l'élève).
\textbf{Durée} : 30 min + devoir écrit.

\begin{enumerate}
\item \textbf{Préparation (5 min).} Lis ta fiche en silence. Souligne les informations clés. Identifie l'auxiliaire de chaque phrase pour préparer tes réponses courtes (\eng{do, can, have, would, is, will}...).
\item \textbf{Échauffement oral (5 min).} En binôme, transformez ces phrases en tags à voix haute (sans jouer) :
\begin{itemize}
\item \eng{You like English.} $\rightarrow$ \eng{You like English, don't you?}
\item \eng{You have visited Kribi.} $\rightarrow$ \eng{You have visited Kribi, haven't you?}
\item \eng{You can swim.} $\rightarrow$ \eng{You can swim, can't you?}
\item \eng{You aren't tired.} $\rightarrow$ \eng{You aren't tired, are you?}
\item \eng{You will pass the Probatoire.} $\rightarrow$ \eng{You will pass the Probatoire, won't you?}
\end{itemize}
Vérifiez la polarité et l'auxiliaire ensemble.
\item \textbf{L'interview (10 min).} À tour de rôle, devine les infos de ton partenaire \textbf{uniquement par question tags}.
\eng{You come from Buea, don't you?} — \eng{Yes, I do.}
\eng{You've never travelled abroad, have you?} — \eng{No, I haven't.}
\eng{You can dance the makossa, can't you?} — \eng{Yes, I can.}
Note chaque info confirmée.
\item \textbf{Contrôle mutuel (5 min).} Signale poliment toute erreur de tag : \eng{Sorry, the tag should be “haven't you”, not “didn't you”.} Note les erreurs pour la correction collective.
\item \textbf{Présentation orale (5 min/binôme).} Présente ton partenaire en \textbf{5 phrases}, dont \textbf{3 avec tags variés}.
Ex. : \eng{This is Chantal. She comes from Buea, doesn't she? She has visited Limbe, hasn't she? She can dance makossa, can't she? Her brother plays football. She will succeed, won't she?}
La classe valide ou corrige chaque tag (main levée).
\item \textbf{Production écrite (Devoir).} Rédige un article (80–120 mots) pour \emph{The Lion's Voice}, rubrique \eng{Getting to know you}. Contraintes : titre, \textbf{au moins 3 tags aux auxiliaires différents}, phrase conclusive.
\end{enumerate}
\end{experience}

\courssub{Ressources à mobiliser}

\begin{aretenir}[Rappel express : la boîte à outils du question tag]
\begin{itemize}
\item \textbf{Structure} : Phrase + , + Auxiliaire (polarité inversée) + Pronom sujet + ?
\item \textbf{Auxiliaires fréquents} : \eng{do/does/did} (Simple Present/Past), \eng{be} (\eng{is/are/was/were}), \eng{have/has} (Perfect), \eng{will/would}, \eng{can/could}, \eng{must}, \eng{should}.
\item \textbf{Exceptions clés} : \eng{I am $\rightarrow$ aren't I?} | \eng{Let's $\rightarrow$ shall we?} | \eng{Impératif $\rightarrow$ will you?} | \eng{There is $\rightarrow$ isn't there?} | \eng{Nobody/Everybody $\rightarrow$ they}.
\item \textbf{Réponses courtes} : \eng{Yes, I do. / No, I don't. / Yes, she has. / No, they can't.} (Toujours auxiliaire seul).
\item \textbf{Intonation} : \searrow = Confirmation (sûr) | \nearrow = Information (doute).
\end{itemize}
\end{aretenir}

\textbf{Lexique utile pour l'interview} :
\begin{itemize}
\item \eng{home town} (ville d'origine) — \eng{favourite subject} (matière préférée) — \eng{a hidden talent} (un talent caché)
\item \eng{to travel abroad} (voyager à l'étranger) — \eng{to look forward to + -ing} (avoir hâte de)
\item \eng{to be fond of + -ing} (aimer beaucoup) — \eng{a dream} (un rêve) — \eng{to confirm information} (confirmer une info)
\end{itemize}

\courssub{Proposition de production et corrigé commenté}

\begin{exoresolu}[Modèle d'article pour \emph{The Lion's Voice} — 80–120 mots, 3 tags variés]
Rédige un article présentant le camarade interviewé.
\tcblower
\textbf{Production modèle :}

\eng{GETTING TO KNOW YOU: Chantal, our future doctor!}

\eng{This week, we interviewed Chantal Ewusi, a brilliant Terminale A student. She comes from Buea, \textbf{doesn't she?} Yes, she does — and she is very proud of her region! Her favourite subject is Biology, and she would like to study medicine at the University of Buea. She has already visited Limbe beach twice, \textbf{hasn't she?} Surprisingly, she has never travelled abroad, but she dreams of discovering London one day. Chantal can also dance the makossa wonderfully, \textbf{can't she?} What she can't stand is waking up before 5 a.m.! We wish her good luck for the Probatoire. She will succeed, \textbf{won't she?}}

\textbf{Commentaire des choix de langue :}
\begin{itemize}
\item \eng{doesn't she?} : Phrase affirmative, présent simple, verbe lexical \eng{comes} $\rightarrow$ auxiliaire \eng{does} + \eng{n't} + \eng{she}.
\item \eng{hasn't she?} : Present perfect \eng{has visited} $\rightarrow$ tag sur \eng{has} (jamais \eng{did}).
\item \eng{can't she?} : Modal \eng{can} $\rightarrow$ tag \eng{can't she}. Trois auxiliaires différents (\eng{do, have, can}) $\rightarrow$ consigne respectée.
\item \eng{won't she?} : Futur \eng{will} $\rightarrow$ tag irrégulier \eng{won't}. Intonation descendante (conviction forte), naturel en conclusion.
\item La réponse courte \eng{Yes, she does} est intégrée au récit (style interview vivante).
\item Note : \eng{She has never travelled abroad} est une phrase négative (\eng{never}) ; si on en faisait un tag, ce serait \eng{has she?} (affirmatif).
\end{itemize}
\end{exoresolu}

\courssub{Exercices d'application}

\begin{exercice}[Transformation : phrases $\rightarrow$ question tags]
Transforme les phrases suivantes en ajoutant le question tag correct.
\begin{enumerate}
\item Your father works at the port, \_\_\_\_\_\_\_\_\_\_ ?
\item We haven't seen the new bridge yet, \_\_\_\_\_\_\_\_\_\_ ?
\item Let's buy some beignets at the market, \_\_\_\_\_\_\_\_\_\_ ?
\item Nobody answered the teacher's question, \_\_\_\_\_\_\_\_\_\_ ?
\item I am the best student in the class, \_\_\_\_\_\_\_\_\_\_ ?
\item Turn down the music, \_\_\_\_\_\_\_\_\_\_ ?
\item There were many people at the match, \_\_\_\_\_\_\_\_\_\_ ?
\item She'd rather stay at home tonight, \_\_\_\_\_\_\_\_\_\_ ?
\item This isn't your bag, \_\_\_\_\_\_\_\_\_\_ ?
\item You'd better revise your tags, \_\_\_\_\_\_\_\_\_\_ ?
\end{enumerate}
\end{exercice}

\begin{exercice}[Réponses courtes : attention au piège !]
Réponds par une réponse courte correcte (Yes/No + auxiliaire) selon la réalité indiquée entre parenthèses.
\begin{enumerate}
\item \eng{You don't speak Pidgin, do you?} (Réalité : tu parles Pidgin) $\rightarrow$ \_\_\_\_\_\_\_\_\_\_
\item \eng{She can't drive, can she?} (Réalité : elle sait conduire) $\rightarrow$ \_\_\_\_\_\_\_\_\_\_
\item \eng{They have finished, haven't they?} (Réalité : ils n'ont pas fini) $\rightarrow$ \_\_\_\_\_\_\_\_\_\_
\item \eng{It isn't raining, is it?} (Réalité : il pleut) $\rightarrow$ \_\_\_\_\_\_\_\_\_\_
\item \eng{Nobody came, did they?} (Réalité : personne n'est venu) $\rightarrow$ \_\_\_\_\_\_\_\_\_\_
\end{enumerate}
\end{exercice}

\begin{exercice}[Correction d'erreurs : détecte et corrige]
Chaque phrase contient une erreur de formation du tag. Souligne l'erreur et réécris le tag correct.
\begin{enumerate}
\item \eng{You like ndolé, isn't it?}
\item \eng{He has got a new phone, doesn't he?}
\item \eng{Let's go to the beach, will we?}
\item \eng{I am right, am not I?}
\item \eng{She has breakfast at 6, hasn't she?} (verbe lexical \eng{have})
\item \eng{Nobody knows the answer, don't they?}
\item \eng{Open the door, do you?}
\item \eng{There is no water, is there?}
\end{enumerate}
\end{exercice}

\courssub{Corrigés des exercices}

\begin{corrige}[Exercice 1 — Transformation]
\begin{enumerate}
\item \eng{don't you?} (Présent simple, \eng{work} $\rightarrow$ \eng{do})
\item \eng{have we?} (Present perfect négatif $\rightarrow$ tag affirmatif \eng{have})
\item \eng{shall we?} (Exception \eng{Let's})
\item \eng{did they?} (\eng{Nobody} = négatif $\rightarrow$ tag affirmatif, pronom \eng{they})
\item \eng{aren't I?} (Exception \eng{I am})
\item \eng{will you?} (Impératif $\rightarrow$ \eng{will you?})
\item \eng{weren't there?} (\eng{There were} $\rightarrow$ \eng{there})
\item \eng{wouldn't she?} (\eng{She'd rather} = \eng{she would rather} $\rightarrow$ tag sur \eng{would})
\item \eng{is it?} (Phrase négative \eng{isn't} $\rightarrow$ tag affirmatif \eng{is}, sujet \eng{it})
\item \eng{hadn't you?} (\eng{You'd better} = \eng{You had better} $\rightarrow$ tag sur \eng{had})
\end{enumerate}
\end{corrige}

\begin{corrige}[Exercice 2 — Réponses courtes]
\begin{enumerate}
\item \eng{Yes, I do.} (Réalité positive $\rightarrow$ \eng{Yes} + auxiliaire affirmatif \eng{do})
\item \eng{Yes, she can.} (Réalité positive $\rightarrow$ \eng{Yes} + \eng{can})
\item \eng{No, they haven't.} (Réalité négative $\rightarrow$ \eng{No} + \eng{haven't})
\item \eng{Yes, it is.} (Réalité positive $\rightarrow$ \eng{Yes} + \eng{is})
\item \eng{No, they didn't.} (Réalité : personne n'est venu = négatif $\rightarrow$ \eng{No} + \eng{didn't})
\end{enumerate}
\end{corrige}

\begin{corrige}[Exercice 3 — Correction d'erreurs]
\begin{enumerate}
\item \eng{*isn't it?} $\rightarrow$ \eng{\textbf{don't you?}} (Verbe lexical \eng{like} $\rightarrow$ auxiliaire \eng{do})
\item \eng{*doesn't he?} $\rightarrow$ \eng{\textbf{hasn't he?}} (\eng{Have got} $\rightarrow$ tag sur \eng{have})
\item \eng{*will we?} $\rightarrow$ \eng{\textbf{shall we?}} (Exception \eng{Let's})
\item \eng{*am not I?} $\rightarrow$ \eng{\textbf{aren't I?}} (Exception \eng{I am})
\item \eng{*hasn't she?} $\rightarrow$ \eng{\textbf{doesn't she?}} (\eng{Have} verbe lexical/action $\rightarrow$ auxiliaire \eng{does})
\item \eng{*don't they?} $\rightarrow$ \eng{\textbf{did they?}} (\eng{Nobody} + prétérit \eng{knew} $\rightarrow$ \eng{did they?}) \textit{Note : si présent \eng{knows} $\rightarrow$ \eng{do they?}}
\item \eng{*do you?} $\rightarrow$ \eng{\textbf{will you?}} (Impératif $\rightarrow$ \eng{will you?})
\item \eng{*is there?} $\rightarrow$ \eng{\textbf{isn't there?}} (\eng{There is no} = phrase négative $\rightarrow$ tag affirmatif \eng{is there?}) \textbf{Attention} : \eng{There is no water} est négatif $\rightarrow$ tag \textbf{affirmatif} \eng{is there?}. La phrase proposée \eng{is there?} est en fait correcte si la phrase est \eng{There is no water}. Mais l'exercice dit "corrige". Si l'original est \eng{There is no water, is there?}, c'est correct. Si l'erreur attendue est \eng{isn't there?}, alors la phrase de base devrait être \eng{There is water}. Je corrige en supposant que l'erreur est l'inversion de polarité : \eng{There is no water, \textbf{is there?}} (Correct). Je modifie l'énoncé de l'exercice pour créer une erreur : \eng{There is water, isn't there?} (Correct). Pour l'exercice, je propose : \eng{There is no water, isn't there?} $\rightarrow$ \eng{\textbf{is there?}}.
\end{enumerate}
\end{corrige}

\courssub{Bilan et auto-évaluation}

\begin{methode}[Auto-évaluation : suis-je prêt pour le Bac ?]
Coche les cases si tu maîtrises la compétence.
\begin{enumerate}
\item Je forme un tag correct en moins de 3 secondes sur n'importe quelle phrase simple (présent, passé, perfect, modaux).
\item Je n'utilise \textbf{jamais} \eng{isn't it?} ou \eng{no?} comme tag universel.
\item Je réponds \eng{Yes, I do} / \eng{No, I don't} selon les \textbf{faits}, même après une phrase négative.
\item Je connais par cœur les 5 exceptions : \eng{I am... aren't I?}, \eng{Let's... shall we?}, Impératif + \eng{will you?}, \eng{There is... isn't there?}, \eng{Nobody/Everybody... they}.
\item Je gère \eng{have got} (\eng{hasn't he?}) vs \eng{have} action (\eng{doesn't he?}).
\item J'écris un article (80–120 mots) avec au moins 3 tags variés (auxiliaires différents) sans erreur.
\item Je distingue l'intonation \searrow (confirmation) et \nearrow (doute réel) à l'oral.
\end{enumerate}
\textbf{Bilan} : 7/7 = Objectif atteint. $<$ 7 = Revoir la règle (boîte \textbf{Rappel express}) + refaire une manche du jeu d'interview avec un nouveau partenaire.
\end{methode}

\begin{savaistu}[Les tags, marqueurs de politesse et de « small talk » anglo-saxonne]
Dans la culture anglophone (au Cameroun comme au Royaume-Uni, aux USA, au Nigeria, au Ghana), les question tags ne sont pas qu'un exercice de grammaire : ce sont des \textbf{lubrifiants sociaux}. Dire \eng{“Lovely weather, isn't it?”} à un collègue à Buea ou à un voisin dans le bus à Londres, c'ouvrir la conversation sans s'imposer. Le tag \eng{isn't it?} (intonation descendante) dit : « Je partage cette observation, es-tu d'accord ? » — c'est de la \eng{small talk}. Maîtriser les tags, c'est donc maîtriser l'art de la relation sociale en anglais, une compétence clé pour tes futurs examens oraux, entretiens et échanges internationaux.
\end{savaistu}

\newpage

\coursec{Méthodes et exercices résolus}

\coursec{Lesson 8 — Grammar: Revise Question Tags}

\begin{definition}[Qu'est-ce qu'un question tag ?]
Un \cle{question tag} (ou \emph{tag question}) est une mini-question ajoutée à la fin d'une phrase affirmative ou négative pour demander une confirmation, exprimer une surprise ou relancer la conversation. En français, on utilise souvent l'invariant « n'est-ce pas ? » ; en anglais, le tag s'accorde obligatoirement avec le \textbf{temps}, l'\textbf{auxiliaire} et le \textbf{sujet} de la phrase principale.
\end{definition}

\begin{textebox}[Observation : dialogue au marché de Buea]
\eng{Customer: The plantains are ripe today, aren't they?}
\eng{Seller: Yes, they are. They came from the farm yesterday.}
\eng{Customer: You didn't use any chemicals, did you?}
\eng{Seller: No, I didn't. Natural only. The price has gone up, hasn't it?}
\eng{Customer: Yes, it has. But we have no choice, do we?}
\eng{Seller: Let's hope for a better harvest next season, shall we?}
\end{textebox}

\begin{propriete}[Règle de formation : le principe de base]
Un question tag se construit en \textbf{trois éléments} :
\begin{center}
\begin{tabularx}{\linewidth}{|X|X|X|}
\hline
\rowcolor{popblueL} \textbf{Élément} & \textbf{Règle} & \textbf{Exemple} \\
\hline
1. \textbf{Auxiliaire} & Reprendre l'auxiliaire de la phrase (be, have, do, modal). S'il n'y en a pas (present simple / past simple verbe lexical) \textrightarrow\ \eng{do / does / did}. & \eng{She \textbf{is} tired} \textrightarrow\ \eng{\textbf{isn't}} \\
\hline
2. \textbf{Polarité inversée} & Phrase affirmative \textrightarrow\ tag \textbf{négatif} (souvent contracté : \eng{n't}). Phrase négative \textrightarrow\ tag \textbf{affirmatif}. & \eng{They \textbf{like} it} \textrightarrow\ \eng{\textbf{don't}} \quad | \quad \eng{He \textbf{isn't} here} \textrightarrow\ \eng{\textbf{is}} \\
\hline
3. \textbf{Pronom sujet} & Remplacer le sujet par le pronom personnel correspondant (\eng{he, she, it, they, we, you, I}). \textbf{Jamais de nom} dans le tag. & \eng{Paul \textbf{is} late} \textrightarrow\ \eng{\textbf{he}} \\
\hline
\end{tabularx}
\end{center}
\eng{Structure : Phrase principale + , + auxiliaire (polarité inversée) + pronom sujet + ?}
\end{propriete}

\begin{methode}[Construire un question tag en 5 étapes]
\begin{enumerate}
\item \textbf{Repérez le verbe} de la phrase principale : auxiliaire (\eng{be, have, do}, modal) ou verbe lexical ?
\item \textbf{Choisissez l'auxiliaire du tag} : reprenez l'auxiliaire existant ; sinon utilisez \eng{do / does / did} selon le temps et la personne.
\item \textbf{Inversez la polarité} : affirmative \textrightarrow\ négatif ; négative \textrightarrow\ affirmatif.
\item \textbf{Remplacez le sujet par un pronom} : \eng{Your brother} \textrightarrow\ \eng{he}, \eng{the cars} \textrightarrow\ \eng{they}, \eng{nobody} \textrightarrow\ \eng{they}.
\item \textbf{Placez le tag après une virgule}, à la fin de la phrase, suivi d'un point d'interrogation.
\end{enumerate}
\begin{popfigure}
\begin{tikzpicture}[
  mindmap, concept color=popblue, text=white,
  level 1 concept/.append style={level distance=3.5cm, sibling angle=72, concept color=popblue!70},
  level 2 concept/.append style={level distance=2.5cm, concept color=poporange},
  every node/.style={font=\small}
]
\node[concept] {Question Tag}
  child[concept color=popgreen] { node[concept] {1. Verbe principal} }
  child[concept color=poppurple] { node[concept] {2. Auxiliaire du tag} }
  child[concept color=poppink] { node[concept] {3. Polarité inversée} }
  child[concept color=popgold] { node[concept] {4. Pronom sujet} }
  child[concept color=popdark] { node[concept] {5. Virgule + ?} };
\end{tikzpicture}
\legende{Les 5 étapes pour construire un question tag correct}
\end{popfigure}
\end{methode}

\begin{exoresolu}[Application : tags standards]
Ajoutez le question tag correct à chaque phrase.
\begin{enumerate}
\item Your sister lives in Douala, \ldots ?
\item The students didn't finish their homework, \ldots ?
\item She can speak three languages, \ldots ?
\item They have been to Buea, \ldots ?
\item It is raining heavily today, \ldots ?
\end{enumerate}
\tcblower
\textbf{Solution détaillée.}
\begin{enumerate}
\item \eng{Your sister lives in Douala, \textbf{doesn't she?}} — Verbe lexical \eng{lives} (present simple, 3\textsuperscript{e} pers. sg) \textrightarrow\ auxiliaire \eng{does} ; phrase affirmative \textrightarrow\ tag négatif ; \eng{your sister} \textrightarrow\ \eng{she}.
\item \eng{The students didn't finish their homework, \textbf{did they?}} — Phrase négative (\eng{didn't}) \textrightarrow\ tag affirmatif \eng{did} ; \eng{the students} \textrightarrow\ \eng{they}.
\item \eng{She can speak three languages, \textbf{can't she?}} — Modal \eng{can} repris au négatif \eng{can't} (contraction de \eng{cannot}).
\item \eng{They have been to Buea, \textbf{haven't they?}} — Auxiliaire \eng{have} (present perfect) repris au négatif.
\item \eng{It is raining heavily today, \textbf{isn't it?}} — Auxiliaire \eng{is} (present continuous) repris au négatif ; \eng{it} inchangé.
\end{enumerate}
\textbf{Conclusion} : le tag copie toujours l'auxiliaire de la phrase, inverse sa polarité et remplace le sujet par un pronom.
\end{exoresolu}

\courssub{Cas particuliers et exceptions}

\begin{propriete}[Tableau récapitulatif des tags irréguliers]
Ces six cas sont des \textbf{formules fixes} à mémoriser par cœur ; ils ne suivent pas la règle générale.
\begin{center}
\begin{tabularx}{\linewidth}{|l|X|X|}
\hline
\rowcolor{popblueL} \textbf{Cas} & \textbf{Structure} & \textbf{Tag correct} \\
\hline
\eng{I am} & \eng{I am late} & \eng{\textbf{aren't I?}} (jamais \eng{amn't I}) \\
\hline
\eng{Let's} (suggestion) & \eng{Let's go} & \eng{\textbf{shall we?}} \\
\hline
Impératif (ordre/demande) & \eng{Open the door} & \eng{\textbf{will you?}} / \eng{\textbf{would you?}} (polarité libre) \\
\hline
\eng{There is / are} & \eng{There is a problem} & \eng{\textbf{isn't there?}} (reprend \eng{there}) \\
\hline
Sujets en \eng{-body/-one} & \eng{Nobody came} & \eng{\textbf{did they?}} (pronom \eng{they}) \\
\hline
Sujets en \eng{-thing} & \eng{Nothing matters} & \eng{\textbf{does it?}} (pronom \eng{it}) \\
\hline
Mots à valeur négative & \eng{He never helps} & \eng{\textbf{does he?}} (\eng{never, hardly, rarely, seldom, few, little} \textrightarrow\ tag affirmatif) \\
\hline
\eng{I think / believe / suppose} + subordonnée & \eng{I think she is right} & \eng{\textbf{isn't she?}} (tag sur la subordonnée, pas sur \eng{I think}) \\
\hline
\end{tabularx}
\end{center}
\end{propriete}

\begin{methode}[Les tags irréguliers : explications détaillées]
\begin{enumerate}
\item \textbf{\eng{I am} \textrightarrow\ \eng{aren't I?}} : \eng{amn't I} n'existe pas en anglais standard ; \eng{aren't I} s'est imposé à l'oral comme à l'écrit.
\item \textbf{\eng{Let's} \textrightarrow\ \eng{shall we?}} : \eng{Let's} = \eng{Let us} (proposition commune). Le tag invite à l'accord.
\item \textbf{Impératif \textrightarrow\ \eng{will you?} / \eng{would you?}} : \eng{will you} (familier/insistant), \eng{would you} (plus poli). La polarité n'est pas inversée (phrase sans auxiliaire apparent).
\item \textbf{\eng{There is/are} \textrightarrow\ \eng{there}} : \eng{There} fonctionne comme un sujet fantôme ; on le reprend dans le tag.
\item \textbf{Indéfinis \eng{-body/-one} \textrightarrow\ \eng{they}} ; \textbf{\eng{-thing} \textrightarrow\ \eng{it}} : accord notionnel (pluriel pour les personnes, singulier pour les choses).
\item \textbf{Mots « semi-négatifs »} (\eng{never, hardly, rarely, seldom, few, little}) : la phrase est affirmative en forme mais \textbf{négative en sens} \textrightarrow\ tag \textbf{affirmatif}.
\item \textbf{Verbes d'opinion} (\eng{I think, I believe, I suppose, I imagine}) : le tag porte sur la proposition subordonnée (\eng{that}-clause), car c'est l'information principale.
\end{enumerate}
\end{methode}

\begin{exoresolu}[Cas particuliers : entraînement]
Complétez avec le question tag approprié.
\begin{enumerate}
\item I am your best friend, \ldots ?
\item Let's visit our grandmother this weekend, \ldots ?
\item Nobody phoned yesterday, \ldots ?
\item She hardly ever goes to the market, \ldots ?
\item I think your father is a teacher, \ldots ?
\item There was no food left, \ldots ?
\end{enumerate}
\tcblower
\textbf{Solution détaillée.}
\begin{enumerate}
\item \eng{\textbf{aren't I}} — Formule fixe avec \eng{I am}.
\item \eng{\textbf{shall we}} — Tag obligatoire après \eng{Let's}.
\item \eng{\textbf{did they}} — \eng{Nobody} (négatif en sens) \textrightarrow\ tag affirmatif ; \eng{-body} \textrightarrow\ \eng{they}.
\item \eng{\textbf{does she}} — \eng{Hardly ever} (presque jamais) = valeur négative \textrightarrow\ tag affirmatif ; present simple 3\textsuperscript{e} pers. \textrightarrow\ \eng{does}.
\item \eng{\textbf{isn't he}} — Après \eng{I think}, tag sur la subordonnée (\eng{your father is} \textrightarrow\ \eng{isn't he}).
\item \eng{\textbf{was there}} — \eng{No food} (négatif en sens) \textrightarrow\ tag affirmatif ; \eng{there was} \textrightarrow\ \eng{was there}.
\end{enumerate}
\textbf{Conclusion} : ces sept formes sont des classiques du Bac ; mémorisez-les comme des blocs indivisibles.
\end{exoresolu}

\courssub{Répondre à un question tag et intonation}

\begin{methode}[Répondre comme un anglophone : la règle du « fait »]
\begin{enumerate}
\item \textbf{La réponse porte sur la réalité (les faits), pas sur la forme de la question.}
  \begin{itemize}
    \item Si le fait est \textbf{vrai} \textrightarrow\ \eng{Yes} + auxiliaire affirmatif.
    \item Si le fait est \textbf{faux} \textrightarrow\ \eng{No} + auxiliaire négatif.
  \end{itemize}
\item \textbf{Ne traduisez jamais le « si » français.} En français, on répond « Si ! » à une question négative ; en anglais, on répond \eng{Yes}.
  \eng{You aren't tired, are you? — \textbf{Yes, I am!}} (« Si, je suis fatigué ! »).
\item \textbf{Accordez l'auxiliaire de la réponse} avec le temps du tag : \eng{Yes, I do. / No, they haven't. / Yes, it will.}
\item \textbf{Intonation} :
  \begin{itemize}
    \item \textbf{Voix descendante} (\searrow) : on attend une confirmation (certitude).
    \item \textbf{Voix montante} (\nearrow) : vraie question, on n'est pas sûr (doute).
  \end{itemize}
\end{enumerate}
\end{methode}

\begin{exoresolu}[Réponses courtes accordées]
Répondez à chaque question tag selon l'indication entre parenthèses.
\begin{enumerate}
\item You don't like football, do you? (mais si, tu adores ça)
\item Cameroon will win the match, won't it? (tu es d'accord)
\item You haven't finished, have you? (non, pas encore)
\end{enumerate}
\tcblower
\textbf{Solution détaillée.}
\begin{enumerate}
\item \eng{\textbf{Yes, I do! I love it.}} — Question négative, fait positif \textrightarrow\ \eng{Yes} + auxiliaire \eng{do}.
\item \eng{\textbf{Yes, I think it will.}} ou \eng{\textbf{Yes, it will.}} — Confirmation \textrightarrow\ \eng{Yes} + modal \eng{will}.
\item \eng{\textbf{No, I haven't. Not yet.}} — Fait négatif \textrightarrow\ \eng{No} + auxiliaire \eng{have} (present perfect).
\end{enumerate}
\textbf{Astuce} : pensez « \textbf{Yes = c'est vrai / No = c'est faux} », oubliez la polarité de la question.
\end{exoresolu}

\begin{savaistu}[Intonation et culture camerounaise]
Au Cameroun anglophone (Nord-Ouest, Sud-Ouest), l'intonation descendante sur le tag (\eng{It's hot today, isn't it?\searrow}) marque la convivialité et la recherche de consensus. L'intonation montante (\eng{You've seen my pen, haven't you?\nearrow}) marque une vraie demande d'information. En anglais standard britannique, la distinction est identique. Attention : en \emph{Cameroonian Pidgin English}, on entend souvent le tag invariant \eng{« no? »} ou \eng{« isn't it? »} pour tous les temps (ex. \eng{You go market, isn't it?}) — c'est une variété à part, \textbf{pas l'anglais standard} attendu au Bac.
\end{savaistu}

\courssub{Production orale et écrite}

\begin{experience}[Jeu de rôle : « Au bureau de poste à Yaoundé »]
\textbf{Consigne} : Par binômes. L'un joue le client, l'autre l'employé. Durée : 3 min.
\textbf{Contraintes} : Utilisez au moins \textbf{4 question tags différents} (dont un cas particulier) et leurs réponses accordées.
\textbf{Scénario suggéré} : Envoyer un colis à l'étranger (poids, prix, délai, assurance).
\textbf{Exemple de début} :
\eng{Client: Good morning. This parcel goes to London, doesn't it?}
\eng{Clerk: Yes, it does. It weighs 2 kg, doesn't it?}
\eng{Client: Yes, it does. The price has increased, hasn't it?}
\eng{Clerk: Yes, it has. Let's check the insurance option, shall we?}
\end{experience}

\begin{methode}[Rédiger un dialogue naturel avec question tags]
\begin{enumerate}
\item \textbf{Variez les auxiliaires} : dans 6--8 répliques, utilisez au moins 3 auxiliaires distincts (\eng{isn't it, don't you, didn't she, will you, shall we...}).
\item \textbf{Alternez les fonctions} : confirmation (\eng{You live in Yaoundé, don't you?}), relance, surprise (\eng{She passed the exam, didn't she?}).
\item \textbf{Soignez les réponses} : chaque tag reçoit une réponse avec auxiliaire (\eng{Yes, I do. / No, she didn't.}).
\item \textbf{Relisez} : auxiliaire correct ? polarité inversée ? pronom (pas de nom) ? contraction (\eng{n't}) ?
\end{enumerate}
\end{methode}

\begin{exoresolu}[Production guidée : dialogue à Buea]
Rédigez un court dialogue (6 répliques) entre deux camarades de classe à Buea, contenant au moins trois question tags corrects et leurs réponses.
\tcblower
\textbf{Modèle de production.}

\eng{Amina: Hi, Peter! You went to the market this morning, \textbf{didn't you?}}

\eng{Peter: \textbf{Yes, I did.} I bought some vegetables for my mother.}

\eng{Amina: The prices have gone up again, \textbf{haven't they?}}

\eng{Peter: \textbf{Yes, they have.} Everything is expensive now. You don't cook for your family, \textbf{do you?}}

\eng{Amina: \textbf{Yes, I do!} I cook every evening. Let's prepare something together on Saturday, \textbf{shall we?}}

\eng{Peter: That's a great idea. See you on Saturday, then!}

\textbf{Justification grammaticale} :
(1) \eng{went} (prétérit verbe lexical) \textrightarrow\ \eng{didn't you} ; (2) \eng{have gone up} (present perfect) \textrightarrow\ \eng{haven't they} ; (3) phrase négative \eng{you don't cook} \textrightarrow\ tag affirmatif \eng{do you}, réponse au « si » français = \eng{Yes, I do} ; (4) \eng{Let's} \textrightarrow\ \eng{shall we}. Chaque sujet de tag est un pronom, chaque polarité est inversée.
\end{exoresolu}

\courssub{Stratégie de correction et pièges du Bac}

\begin{methode}[Corriger un question tag fautif : 4 vérifications]
Face à une phrase à corriger (type « \emph{Correct the mistakes} ») :
\begin{enumerate}
\item \textbf{Auxiliaire} : identique à la phrase principale (temps, personne) ?
\item \textbf{Polarité} : affirmative $\leftrightarrow$ négative ; attention aux mots à valeur négative (\eng{never, hardly, nobody, no one, nothing}).
\item \textbf{Pronom} : jamais de nom ; \eng{they} pour \eng{-body/-one}, \eng{it} pour \eng{-thing}, \eng{there} pour \eng{there is/are}.
\item \textbf{Contraction} : la forme contractée (\eng{isn't, don't, can't, haven't}) est la norme dans les tags.
\end{enumerate}
\end{methode}

\begin{exoresolu}[Correction d'erreurs classiques]
Chaque phrase contient une erreur dans le question tag. Corrigez-la.
\begin{enumerate}
\item You are coming to the party, aren't you coming?
\item She likes chocolate, doesn't she likes?
\item They didn't arrive on time, didn't they?
\item Everyone was happy, wasn't everyone?
\item I am right, amn't I?
\end{enumerate}
\tcblower
\textbf{Solution détaillée.}
\begin{enumerate}
\item \eng{You are coming to the party, \textbf{aren't you?}} — Le tag ne contient que l'auxiliaire + pronom ; on ne répète \textbf{jamais} le verbe lexical (\eng{coming}).
\item \eng{She likes chocolate, \textbf{doesn't she?}} — Le tag s'arrête au pronom sujet ; pas de verbe après le pronom.
\item \eng{They didn't arrive on time, \textbf{did they?}} — Phrase négative \textrightarrow\ tag \textbf{affirmatif}. La double négation (\eng{didn't ... didn't they}) est une erreur majeure.
\item \eng{Everyone was happy, \textbf{weren't they?}} — \eng{Everyone} (indéfini en \eng{-one}) \textrightarrow\ pronom \eng{they} au tag (et non \eng{everyone} ni \eng{he/she}).
\item \eng{I am right, \textbf{aren't I?}} — \eng{amn't I} n'existe pas en anglais standard ; la forme figée est \eng{aren't I}.
\end{enumerate}
\textbf{Conclusion} : auxiliaire, polarité, pronom, contraction — quatre contrôles suffisent à corriger n'importe quel tag.
\end{exoresolu}

\begin{attention}[Les 5 erreurs qui coûtent des points au Bac]
\begin{enumerate}
\item \textbf{La double négation} : \eng{She isn't coming, isn't she?} $\times$ \textrightarrow\ \eng{She isn't coming, \textbf{is she?}} Phrase négative \textrightarrow\ tag \textbf{affirmatif}.
\item \textbf{Le nom dans le tag} : \eng{Paul is tired, isn't Paul?} $\times$ \textrightarrow\ \eng{Paul is tired, \textbf{isn't he?}} Toujours un pronom.
\item \textbf{Traduire « n'est-ce pas » mot à mot} : le français a une formule unique ; l'anglais adapte le tag. \eng{Tu viens, n'est-ce pas ?} = \eng{You are coming, \textbf{aren't you?}} (pas \eng{isn't it?}).
\item \textbf{Oublier les cas fixes} : \eng{I am..., \textbf{aren't I?}} ; \eng{Let's..., \textbf{shall we?}} ; \eng{nobody..., \textbf{did they?}} ; \eng{never / hardly} \textrightarrow\ tag \textbf{affirmatif}.
\item \textbf{Répondre « si » par \eng{No}} : \eng{You don't smoke, do you?} — Francophone : « Si. » $\times$ \textrightarrow\ Anglophone : \eng{\textbf{Yes, I do!}} Confondre les deux inverse le sens de la réponse.
\end{enumerate}
\end{attention}

\begin{exercice}[Entraînement personnel — À faire seul(e)]
\begin{enumerate}
\item Complétez avec le question tag correct :
  \begin{enumerate}
  \item Your father works for the CDC, \ldots ?
  \item The match hasn't started yet, \ldots ?
  \item I'm not late, \ldots ?
  \item Nobody knows the answer, \ldots ?
  \item Let's take a break, \ldots ?
  \item She'd rather stay at home, \ldots ? (\textit{indice : had rather / would rather})
  \end{enumerate}
\item Corrigez l'erreur dans chaque tag :
  \begin{enumerate}
  \item They have finished, haven't they finished?
  \item You won't tell anyone, won't you?
  \item Nothing happened, happened it?
  \item I think he is honest, don't I?
  \item Open the window, don't you?
  \end{enumerate}
\item Répondez par \eng{Yes} ou \eng{No} + auxiliaire :
  \begin{enumerate}
  \item You haven't seen my keys, have you? (Fait : je les ai vus.)
  \item It wasn't a good movie, was it? (Fait : c'était nul.)
  \item She can't drive, can she? (Fait : elle sait conduire.)
  \end{enumerate}
\item \textbf{Production} : Écrivez un paragraphe de 5 lignes (environ 50 mots) sur « \emph{The challenges of student life in Yaoundé} » en incluant \textbf{au moins 3 question tags variés} (dont un cas particulier).
\end{enumerate}
\end{exercice}

\begin{corrige}[Exercice n°1 -- Complétion]
\begin{enumerate}
\item \eng{doesn't he?} (present simple, 3\textsuperscript{e} pers. sg, affirmatif \textrightarrow\ négatif)
\item \eng{has it?} (present perfect, négatif \textrightarrow\ affirmatif ; \eng{the match} \textrightarrow\ \eng{it})
\item \eng{am I?} (\eng{I'm not} = négatif \textrightarrow\ tag affirmatif ; cas \eng{I am} affirmatif \textrightarrow\ \eng{aren't I}, mais ici phrase négative)
\item \eng{do they?} (\eng{Nobody} = négatif en sens \textrightarrow\ tag affirmatif ; \eng{-body} \textrightarrow\ \eng{they})
\item \eng{shall we?} (cas fixe après \eng{Let's})
\item \eng{would she?} / \eng{wouldn't she?} — \eng{She'd rather} = \eng{She would rather}. Usage standard : \eng{She'd rather stay, \textbf{wouldn't she?}} (tag négatif car \eng{would rather} exprime une préférence positive). \textit{Note :} \eng{would she?} est possible mais moins courant ; \eng{wouldn't she?} est la norme attendue.
\end{enumerate}
\end{corrige}

\begin{corrige}[Exercice n°2 -- Correction]
\begin{enumerate}
\item \eng{They have finished, \textbf{haven't they?}} (pas de répétition du participe passé).
\item \eng{You won't tell anyone, \textbf{will you?}} (phrase négative \textrightarrow\ tag affirmatif).
\item \eng{Nothing happened, \textbf{did it?}} (\eng{Nothing} = négatif en sens \textrightarrow\ tag affirmatif ; \eng{-thing} \textrightarrow\ \eng{it} ; past simple \textrightarrow\ \eng{did}).
\item \eng{I think he is honest, \textbf{isn't he?}} (tag sur la subordonnée \eng{he is}, pas sur \eng{I think}).
\item \eng{Open the window, \textbf{will you?}} (ou \eng{would you?}) ; impératif \textrightarrow\ \eng{will/would you}, pas \eng{don't you}.
\end{enumerate}
\end{corrige}

\begin{corrige}[Exercice n°3 -- Réponses courtes]
\begin{enumerate}
\item \eng{\textbf{Yes, I have.}} (Fait positif \textrightarrow\ \eng{Yes} + auxiliaire \eng{have}).
\item \eng{\textbf{No, it wasn't.}} (Fait négatif \textrightarrow\ \eng{No} + auxiliaire \eng{was} négatif).
\item \eng{\textbf{Yes, she can.}} (Fait positif \textrightarrow\ \eng{Yes} + modal \eng{can}).
\end{enumerate}
\end{corrige}

\begin{corrige}[Exercice n°4 -- Production écrite (modèle)]
\eng{Student life in Yaoundé is challenging, \textbf{isn't it?} Transport is often unreliable; bendskins are fast but dangerous, \textbf{aren't they?} Many students hardly have time to cook, \textbf{do they?} Let's hope the new campus buses arrive soon, \textbf{shall we?} \textbf{Yes, they will} make a big difference.}

\textbf{Analyse} : (1) \eng{isn't it?} (be + pronom) ; (2) \eng{aren't they?} (be + they) ; (3) \eng{do they?} (\eng{hardly} = négatif en sens \textrightarrow\ tag affirmatif + do) ; (4) \eng{shall we?} (cas \eng{Let's}) ; (5) réponse \eng{Yes, they will} (confirmation future). Variété d'auxiliaires et cas particuliers respectés.
\end{corrige}

\newpage

\coursec{Exercices}

\courssub{Je vérifie mes connaissances}

\begin{exercice}[QCM : le bon tag]
Entoure la bonne réponse. Un seul choix est correct.
\begin{enumerate}
\item She's been to Lagos, \cle{...} ? \quad a) isn't she \quad b) hasn't she \quad c) doesn't she \quad d) wasn't she
\item Your father works in Douala, \cle{...} ? \quad a) worksn't he \quad b) isn't he \quad c) doesn't he \quad d) didn't he
\item Let's go to the market, \cle{...} ? \quad a) do we \quad b) shall we \quad c) don't we \quad d) will we
\item Nobody came to the meeting, \cle{...} ? \quad a) did they \quad b) didn't they \quad c) did nobody \quad d) do they
\item I am late, \cle{...} ? \quad a) am I \quad b) amn't I \quad c) aren't I \quad d) isn't I
\item You won't forget the match, \cle{...} ? \quad a) do you \quad b) won't you \quad c) will you \quad d) don't you
\end{enumerate}
\end{exercice}

\begin{exercice}[Vrai ou faux ?]
Dis si les phrases suivantes sont \textbf{correctes} ou \textbf{incorrectes}, puis justifie ta réponse en une phrase en français.
\begin{enumerate}
\item \eng{You like tea, isn't it?}
\item \eng{She can swim, can't she?}
\item \eng{They didn't win, did they?}
\item \eng{Open the door, do you?}
\item \eng{He never smiles, does he?}
\item \eng{We are friends, aren't we?}
\item \eng{There is a problem, isn't there?}
\item \eng{You have finished, haven't you?}
\end{enumerate}
\end{exercice}

\begin{exercice}[Vocabulaire : les auxiliaires des tags]
Relie chaque début de phrase à l'auxiliaire qui doit apparaître dans le question tag.
\begin{center}
\begin{tabularx}{\linewidth}{|X|X|}
\hline
\rowcolor{popblueL} Début de phrase & Auxiliaire du tag \\
\hline
\eng{She has lived in Buea for ten years, ...} & \eng{do / does / did} \\
\hline
\eng{They will travel to Bamenda, ...} & \eng{be (am / is / are)} \\
\hline
\eng{He plays football every Saturday, ...} & \eng{have / has} \\
\hline
\eng{You are very tired, ...} & \eng{will / would} \\
\hline
\end{tabularx}
\end{center}
\end{exercice}

\begin{exercice}[Conjugaison : complète le tag]
Complète chaque phrase avec le question tag correct. Attention au temps et à la personne.
\begin{enumerate}
\item \eng{Your sister passed the exam, ...?}
\item \eng{We must respect our parents, ...?}
\item \eng{The bendskins were very expensive yesterday, ...?}
\item \eng{She doesn't speak French, ...?}
\item \eng{You have never visited London, ...?}
\item \eng{The teacher is marking our copies, ...?}
\end{enumerate}
\end{exercice}

\courssub{Je m'entraîne}

\begin{exercice}[Traduction guidée]
Traduis en anglais les phrases suivantes en utilisant le question tag approprié. Attention : le français « n'est-ce pas » se traduit différemment selon la phrase.
\begin{enumerate}
\item Tu viens de Yaoundé, n'est-ce pas ?
\item Elle n'a pas encore fini ses devoirs, n'est-ce pas ?
\item Nous irons au stade demain, n'est-ce pas ?
\item Il pleut beaucoup à Buea, n'est-ce pas ?
\item Ton frère sait conduire, n'est-ce pas ?
\item Ils n'ont jamais mangé de \emph{ndolé}, n'est-ce pas ?
\end{enumerate}
\end{exercice}

\begin{exercice}[Correction d'erreurs]
Un élève francophone a écrit les phrases suivantes. Chacune contient \textbf{une erreur} dans le question tag. Souligne l'erreur et réécris la phrase correctement.
\begin{enumerate}
\item \eng{You like tea, isn't it?}
\item \eng{She can drive, can't it?}
\item \eng{They arrived late, didn't arrive they?}
\item \eng{I am your friend, amn't I?}
\item \eng{He plays well, playsn't he?}
\item \eng{We have seen the film, haven't we seen?}
\item \eng{Nothing is ready, isn't nothing?}
\item \eng{Let's start now, do we?}
\end{enumerate}
\end{exercice}

\begin{exercice}[Texte à trous]
Complète le dialogue suivant avec les question tags corrects.

\begin{textebox}[At the school canteen]
Amina: You are new in this school, \eng{(1) ...} ?\\
Peter: Yes, I am. I arrived from Kumba last week.\\
Amina: You don't know many people yet, \eng{(2) ...} ?\\
Peter: No, I don't. But the students are friendly, \eng{(3) ...} ?\\
Amina: Yes, they are! You have met our English teacher, \eng{(4) ...} ?\\
Peter: Yes, I have. She speaks very fast, \eng{(5) ...} ?\\
Amina: She does! But you will get used to her, \eng{(6) ...} ?\\
Peter: I hope so. Let's have lunch together, \eng{(7) ...} ?\\
Amina: Good idea! You like \emph{eru}, \eng{(8) ...} ?\\
Peter: I have never tasted it!\\
Amina: Then you will try it today, \eng{(9) ...} ?\\
Peter: I am sure I will love it, \eng{(10) ...} ?
\end{textebox}
\end{exercice}

\begin{exercice}[Réponses courtes]
Réponds aux questions suivantes par une réponse courte correcte (\eng{Yes, ...} / \eng{No, ...} + auxiliaire), en respectant l'indication entre parenthèses.
\begin{enumerate}
\item \eng{You have done your homework, haven't you?} (\checkmark)
\item \eng{She lives in Douala, doesn't she?} (\texttimes)
\item \eng{They will come to the party, won't they?} (\checkmark)
\item \eng{Your brother can swim, can't he?} (\texttimes)
\item \eng{We are late, aren't we?} (\checkmark)
\item \eng{He didn't call you, did he?} (\texttimes)
\end{enumerate}
\end{exercice}

\courssub{Je me prépare au Bac}

\begin{exercice}[Compréhension et langue — type Bac]
Lis le texte suivant, puis réponds aux questions.

\begin{textebox}[A new student in Yaoundé]
When Kemi arrived at Lycée de Nkol-Eton in Yaoundé, she felt nervous. She had spent all her life in Lagos, and everything in Cameroon seemed different. “You will make friends quickly, won't you?” her mother had asked before leaving. Kemi was not so sure.

On her first morning, a tall girl called Sandrine smiled at her. “You are the new student from Nigeria, aren't you?” she asked. Kemi nodded. “You don't speak much French yet, do you?” Sandrine continued kindly. “No, I don't,” Kemi admitted, “but I am learning.” “Let's sit together at break, shall we?” Sandrine suggested.

By the end of the week, Kemi had met the whole class. The teachers were demanding, but nobody was unkind to her. “Nobody has laughed at my accent, have they?” she said to Sandrine, surprised. “Of course not! We all learn from each other here,” her new friend replied. Kemi smiled. Her mother had been right, hadn't she?
\end{textebox}

\textbf{A. Compréhension.} Réponds en anglais par des phrases complètes.
\begin{enumerate}
\item Where did Kemi live before coming to Cameroon?
\item How did she feel on her first day? Why?
\item Who was the first student to speak to her?
\item What surprised Kemi at the end of the week?
\end{enumerate}

\textbf{B. Langue.}
\begin{enumerate}
\item Relève dans le texte \textbf{quatre} question tags différents et recopie les phrases complètes qui les contiennent.
\item Pour chacun de ces tags, explique en français pourquoi il est formé ainsi (auxiliaire, affirmation/négation).
\item Complète : \eng{Kemi had never visited Yaoundé before, ...?} et \eng{Sandrine can speak English, ...?}
\end{enumerate}
\end{exercice}

\begin{exercice}[Traduction — type Bac]
Traduis en anglais le passage suivant. Soigne particulièrement les question tags et les réponses courtes.

\begin{textebox}[Texte à traduire]
— Tu as déjà visité le marché de Mokolo, n'est-ce pas ?\\
— Non, je ne l'ai pas encore visité. Il est très grand, n'est-ce pas ?\\
— Oui, c'est le plus grand marché de Yaoundé. Ta sœur t'accompagnera, n'est-ce pas ?\\
— Bien sûr ! Elle connaît tous les vendeurs. Allons-y samedi matin, d'accord ?\\
— D'accord. Mais ne partons pas trop tard : il y aura beaucoup de monde.
\end{textebox}
\end{exercice}

\begin{exercice}[Expression écrite — type Bac]
\textbf{Sujet.} A new English-speaking student has just arrived in your class. Write a dialogue (\textbf{120 à 150 mots}) between you and this student: you welcome him or her, ask questions about his or her family, country and hobbies, and invite him or her to discover your town.

\textbf{Contraintes obligatoires :}
\begin{itemize}
\item ton dialogue doit contenir \textbf{au moins 5 question tags variés} (temps différents, dont au moins un cas particulier : \eng{I am}, \eng{Let's} ou impératif) ;
\item il doit contenir \textbf{au moins 3 réponses courtes} correctes (\eng{Yes, I have.} / \eng{No, she doesn't.}...) ;
\item utilise les prénoms des deux personnages en début de réplique.
\end{itemize}
\end{exercice}

\newpage

\coursec{Corrigés des exercices}

\begin{corrige}[Exercice 1 — QCM : le bon tag]
\begin{enumerate}
\item \textbf{b) hasn't she} — \eng{She's been to Lagos, hasn't she?} Ici \eng{'s} est la contraction de \emph{has} (present perfect) : le tag reprend \emph{has} à la forme négative.
\item \textbf{c) doesn't he} — \eng{Your father works in Douala, doesn't he?} Verbe lexical au present simple : auxiliaire \emph{do/does} dans le tag ; \eng{worksn’t} n'existe pas.
\item \textbf{b) shall we} — \eng{Let's go to the market, shall we?} Cas particulier : après \eng{Let's...}, le tag est toujours \eng{shall we?}.
\item \textbf{a) did they} — \eng{Nobody came to the meeting, did they?} \eng{Nobody} rend la phrase négative : tag affirmatif ; le pronom du tag est \emph{they}.
\item \textbf{c) aren't I} — \eng{I am late, aren't I?} Cas particulier : \eng{amn't I} n'existe pas en anglais standard ; on dit \eng{aren't I?}.
\item \textbf{c) will you} — \eng{You won't forget the match, will you?} Phrase négative (\emph{won't}) : tag affirmatif avec le même auxiliaire \emph{will}.
\end{enumerate}
\end{corrige}

\begin{corrige}[Exercice 2 — Vrai ou faux ?]
\begin{enumerate}
\item \textbf{Incorrecte.} Le sujet est \emph{you} et le verbe \emph{like} est un verbe lexical : le tag doit être \eng{don't you?} et non \eng{isn't it?}.
\item \textbf{Correcte.} Auxiliaire modal \emph{can}, phrase affirmative, sujet \emph{she} : \eng{can't she?} est bien formé.
\item \textbf{Correcte.} Phrase négative au prétérit (\emph{didn't}) : tag affirmatif \eng{did they?}.
\item \textbf{Incorrecte.} Après un impératif, le tag est \eng{will you?} (ou \eng{would you?}), jamais \eng{do you?}.
\item \textbf{Correcte.} \emph{Never} rend la phrase négative : le tag est donc affirmatif, \eng{does he?}.
\item \textbf{Correcte.} Verbe \emph{be} au present, sujet \emph{we} : \eng{aren't we?}.
\item \textbf{Correcte.} Avec \eng{there is/are}, le pronom du tag est \emph{there} : \eng{isn't there?}.
\item \textbf{Correcte.} Present perfect (\emph{have finished}) : tag négatif \eng{haven't you?}.
\end{enumerate}
\end{corrige}

\begin{corrige}[Exercice 3 — Vocabulaire : les auxiliaires des tags]
\begin{center}
\begin{tabularx}{\linewidth}{|X|X|}
\hline
\rowcolor{popblueL} Début de phrase & Auxiliaire du tag \\
\hline
\eng{She has lived in Buea for ten years, ...} & \eng{have / has} \quad (\eng{hasn't she?}) \\
\hline
\eng{They will travel to Bamenda, ...} & \eng{will / would} \quad (\eng{won't they?}) \\
\hline
\eng{He plays football every Saturday, ...} & \eng{do / does / did} \quad (\eng{doesn't he?}) \\
\hline
\eng{You are very tired, ...} & \eng{be (am / is / are)} \quad (\eng{aren't you?}) \\
\hline
\end{tabularx}
\end{center}
\textbf{Règle rappelée :} le tag reprend l'\cle{auxiliaire} de la phrase principale ; s'il n'y a pas d'auxiliaire (verbe lexical au present ou au prétérit simple), on utilise \eng{do / does / did}.
\end{corrige}

\begin{corrige}[Exercice 4 — Conjugaison : complète le tag]
\begin{enumerate}
\item \eng{Your sister passed the exam, didn't she?} — prétérit simple, verbe lexical, sujet féminin.
\item \eng{We must respect our parents, mustn't we?} — le modal \emph{must} se répète au négatif.
\item \eng{The bendskins were very expensive yesterday, weren't they?} — \emph{were} au négatif ; \emph{bendskins} = \emph{they}.
\item \eng{She doesn't speak French, does she?} — phrase négative : tag affirmatif.
\item \eng{You have never visited London, have you?} — \emph{never} rend la phrase négative : tag affirmatif.
\item \eng{The teacher is marking our copies, isn't he/she?} — present continu avec \emph{be} ; on accepte \eng{isn't he?} ou \eng{isn't she?}.
\end{enumerate}
\end{corrige}

\begin{corrige}[Exercice 5 — Traduction guidée]
\begin{enumerate}
\item \eng{You come from Yaoundé, don't you?} — verbe lexical au present, phrase affirmative.
\item \eng{She hasn't finished her homework yet, has she?} — phrase négative au present perfect : tag affirmatif.
\item \eng{We will go to the stadium tomorrow, won't we?} — futur avec \emph{will}.
\item \eng{It rains a lot in Buea, doesn't it?} — sujet \emph{it}, present simple.
\item \eng{Your brother can drive, can't he?} — modal \emph{can}.
\item \eng{They have never eaten \emph{ndolé}, have they?} — \emph{never} = phrase négative : tag affirmatif.
\end{enumerate}
\textbf{Point de vigilance :} le français « n'est-ce pas » est invariable ; en anglais, le tag change selon le temps, l'auxiliaire, le sujet et la polarité (affirmatif/négatif) de la phrase.
\end{corrige}

\begin{corrige}[Exercice 6 — Correction d'erreurs]
\begin{enumerate}
\item Erreur : \eng{isn't it}. Correction : \eng{You like tea, don't you?} — verbe lexical, sujet \emph{you}.
\item Erreur : \eng{can't it}. Correction : \eng{She can drive, can't she?} — le pronom du tag reprend le sujet \emph{she}.
\item Erreur : \eng{didn't arrive they}. Correction : \eng{They arrived late, didn't they?} — le tag ne contient que l'auxiliaire + le pronom, jamais le verbe lexical.
\item Erreur : \eng{amn't I}. Correction : \eng{I am your friend, aren't I?} — forme standard obligatoire.
\item Erreur : \eng{playsn't he}. Correction : \eng{He plays well, doesn't he?} — on ne contracte jamais un verbe lexical avec \emph{n't} ; on utilise \emph{does}.
\item Erreur : \eng{haven't we seen}. Correction : \eng{We have seen the film, haven't we?} — pas de participe passé dans le tag.
\item Erreur : \eng{isn't nothing}. Correction : \eng{Nothing is ready, is it?} — \emph{nothing} rend la phrase négative ; le tag est affirmatif et le pronom est \emph{it}.
\item Erreur : \eng{do we}. Correction : \eng{Let's start now, shall we?} — après \eng{Let's}, le tag est toujours \eng{shall we?}.
\end{enumerate}
\end{corrige}

\begin{corrige}[Exercice 7 — Texte à trous]
\begin{enumerate}
\item \eng{aren't you} — verbe \emph{be}, phrase affirmative.
\item \eng{do you} — phrase négative (\emph{don't know}) : tag affirmatif.
\item \eng{aren't they} — \emph{the students} = \emph{they}, verbe \emph{be}.
\item \eng{haven't you} — present perfect (\emph{have met}).
\item \eng{doesn't she} — verbe lexical au present, sujet \emph{she}.
\item \eng{won't you} — futur avec \emph{will}.
\item \eng{shall we} — cas particulier après \eng{Let's}.
\item \eng{don't you} — verbe lexical \emph{like}, present simple.
\item \eng{won't you} — futur, phrase affirmative.
\item \eng{won't I} — futur avec \emph{will}, sujet \emph{I}. (On accepte aussi l'interprétation \eng{am I not sure...} : non, ici le verbe est \emph{will love}, donc \eng{won't I?}.)
\end{enumerate}
\textbf{Dialogue complet vérifié :} chaque tag reprend l'auxiliaire de sa phrase, avec inversion de la polarité sauf cas particulier (\eng{Let's... shall we?}).
\end{corrige}

\begin{corrige}[Exercice 8 — Réponses courtes]
\begin{enumerate}
\item \eng{Yes, I have.} — on répond selon la réalité, pas selon la forme négative du tag.
\item \eng{No, she doesn't.} — réponse négative : auxiliaire à la forme négative.
\item \eng{Yes, they will.}
\item \eng{No, he can't.}
\item \eng{Yes, we are.}
\item \eng{No, he didn't.}
\end{enumerate}
\textbf{Point de vigilance :} en anglais, on ne dit jamais \eng{Yes, I haven't.} ; \emph{yes} s'accorde avec une forme affirmative, \emph{no} avec une forme négative. Contrairement au français « si », l'anglais n'a pas de mot spécial pour contredire une question négative : on dit simplement \eng{Yes, ...} + auxiliaire affirmatif.
\end{corrige}

\begin{corrige}[Exercice 9 — Compréhension et langue — type Bac]
\textbf{A. Compréhension} (réponses attendues, formulations voisines acceptées) :
\begin{enumerate}
\item \eng{She lived in Lagos, Nigeria, before coming to Cameroon.} (ou \eng{She had spent all her life in Lagos.})
\item \eng{She felt nervous because everything in Cameroon seemed different to her.}
\item \eng{The first student to speak to her was a tall girl called Sandrine.}
\item \eng{She was surprised that nobody had laughed at her accent.}
\end{enumerate}

\textbf{B. Langue.}
\begin{enumerate}
\item Quatre phrases avec question tags relevées dans le texte :
\begin{itemize}
\item \eng{“You will make friends quickly, won't you?”}
\item \eng{“You are the new student from Nigeria, aren't you?”}
\item \eng{“You don't speak much French yet, do you?”}
\item \eng{“Let's sit together at break, shall we?”}
\item \eng{“Nobody has laughed at my accent, have they?”}
\item \eng{Her mother had been right, hadn't she?}
\end{itemize}
(Quatre au choix parmi ces six.)
\item Explications :
\begin{itemize}
\item \eng{won't you} : futur avec \emph{will}, phrase affirmative, donc tag négatif.
\item \eng{aren't you} : verbe \emph{be} au present, phrase affirmative, tag négatif.
\item \eng{do you} : phrase négative (\emph{don't speak}), verbe lexical, donc tag affirmatif avec \emph{do}.
\item \eng{shall we} : cas particulier obligatoire après \eng{Let's}.
\item \eng{have they} : \emph{nobody} rend la phrase négative ; tag affirmatif, pronom \emph{they}, auxiliaire \emph{have} du present perfect.
\item \eng{hadn't she} : past perfect (\emph{had been}), phrase affirmative, tag négatif.
\end{itemize}
\item \eng{Kemi had never visited Yaoundé before, had she?} (\emph{never} = phrase négative, past perfect, tag affirmatif.) \quad \eng{Sandrine can speak English, can't she?} (modal \emph{can}, phrase affirmative.)
\end{enumerate}

\textbf{Barème indicatif (sur 20) :} A. Compréhension : 8 pts (2 pts par réponse correcte et bien formulée). B. Langue : 12 pts — 4 pts pour les quatre tags relevés (1 pt chacun), 4 pts pour les explications (1 pt chacune), 4 pts pour les deux tags à compléter (2 pts chacun).

\textbf{Points de vigilance :} ne pas recopier un tag sans sa phrase complète ; bien identifier \emph{nobody} comme élément négatif ; ne pas oublier que \eng{Let's} impose \eng{shall we?}.
\end{corrige}

\begin{corrige}[Exercice 10 — Traduction — type Bac]
\textbf{Traduction de référence :}
\begin{textebox}[Model translation]
— \eng{You have already visited the Mokolo market, haven't you?}\\
— \eng{No, I haven't visited it yet. It is very big, isn't it?}\\
— \eng{Yes, it is the biggest market in Yaoundé. Your sister will go with you, won't she?}\\
— \eng{Of course! She knows all the sellers. Let's go there on Saturday morning, shall we?}\\
— \eng{All right. But let's not leave too late: there will be a lot of people.}
\end{textebox}
\textbf{Justifications :}
\begin{itemize}
\item \eng{haven't you} : present perfect (\emph{have visited}), phrase affirmative.
\item \eng{No, I haven't} : réponse courte négative ; « pas encore » = \emph{not... yet}.
\item \eng{isn't it} : verbe \emph{be}, sujet \emph{it}.
\item \eng{won't she} : futur, sujet \emph{your sister} = \emph{she} ; « accompagner » se rend ici par \emph{go with you} (ou \emph{accompany you}).
\item \eng{shall we} : « allons-y » = \eng{Let's go there}, tag obligatoire \eng{shall we?}.
\item « il y aura beaucoup de monde » = \eng{there will be a lot of people} (\emph{people} est pluriel en anglais).
\end{itemize}
\textbf{Barème indicatif (sur 10) :} 1,5 pt par réplique (5 répliques = 7,5 pts) ; 2,5 pts pour l'exactitude globale des quatre question tags et de la réponse courte. Toute erreur de tag (mauvais auxiliaire, mauvaise polarité) coûte 0,5 pt.

\textbf{Points de vigilance :} ne pas traduire « n'est-ce pas » mécaniquement par \eng{isn't it} partout ; accorder la réponse courte avec l'auxiliaire de la question ; \eng{Let's not leave} pour « ne partons pas ».
\end{corrige}

\begin{corrige}[Exercice 11 — Expression écrite — type Bac]
\textbf{Proposition de rédaction modèle (environ 135 mots) :}
\begin{textebox}[Model dialogue]
Mireille: Hello! Welcome to our class. You are Daniel, the new student from Bamenda, aren't you?\\
Daniel: Yes, I am. Pleased to meet you.\\
Mireille: You have already visited Yaoundé, haven't you?\\
Daniel: No, I haven't. It is my first time here.\\
Mireille: Really? Your family stayed in Bamenda, didn't they?\\
Daniel: Yes, they did. My father works there.\\
Mireille: You like football, don't you? Most boys here do!\\
Daniel: Yes, I do. I also enjoy music.\\
Mireille: Great! I am sure you will love our town, won't you?\\
Daniel: I hope so!\\
Mireille: Let's visit the Reunification Monument after school, shall we?\\
Daniel: Good idea! You will show me the way, won't you?\\
Mireille: Of course! And you can teach me some Pidgin, can't you?\\
Daniel: Yes, I can! Let's start today!
\end{textebox}

\textbf{Vérification des contraintes :}
\begin{itemize}
\item 7 question tags variés : \eng{aren't you} (\emph{be}), \eng{haven't you} (present perfect), \eng{didn't they} (prétérit), \eng{don't you} (present simple), \eng{won't you} (futur), \eng{shall we} (cas particulier \eng{Let's}), \eng{can't you} (modal).
\item 4 réponses courtes : \eng{Yes, I am.} / \eng{No, I haven't.} / \eng{Yes, they did.} / \eng{Yes, I do.} / \eng{Yes, I can!}
\item Les prénoms (\emph{Mireille}, \emph{Daniel}) ouvrent chaque réplique.
\end{itemize}

\textbf{Barème indicatif (sur 20) :} respect de la tâche et de la longueur (4 pts) ; au moins 5 question tags corrects et variés, dont un cas particulier (6 pts) ; au moins 3 réponses courtes correctes (3 pts) ; richesse du vocabulaire et exactitude grammaticale générale (5 pts) ; présentation du dialogue avec prénoms (2 pts).

\textbf{Points de vigilance :} varier les temps des tags (éviter cinq fois \eng{don't you?}) ; vérifier la polarité de chaque tag ; ne pas écrire \eng{amn't I} ; soigner la ponctuation anglaise du dialogue (tiret ou nom suivi de deux-points, guillemets “...” si direct speech).
\end{corrige}

\newpage

\coursec{Fiche bilan}

\begin{popfigure}
\begin{tikzpicture}[mindmap, grow cyclic, every node/.style={concept, align=center, text width=2.4cm, font=\small}, root concept/.append style={concept color=popblue, text=white, font=\small\bfseries, text width=2.8cm}, level 1/.append style={level distance=3.4cm, sibling angle=51.4}, level 1 concept/.append style={font=\footnotesize, text width=2.5cm}]
\node[root concept]{Question tags}
  child[concept color=poporange, text=white]{ node[align=center]{Principe de base\\ auxiliaire + pronom\\ polarité inversée} }
  child[concept color=popgreen, text=white]{ node[align=center]{Cas particuliers\\ I am $\rightarrow$ aren't I\\ Let's $\rightarrow$ shall we} }
  child[concept color=poppurple, text=white]{ node[align=center]{Impératifs\\ will you / won't you} }
  child[concept color=poppink, text=white]{ node[align=center]{Indéfinis\\ nobody, nothing\\ $\rightarrow$ they / it} }
  child[concept color=popgold, text=black]{ node[align=center]{Réponses\\ Yes / No selon\\ la réalité} }
  child[concept color=popblue, text=white]{ node[align=center]{Intonation\\ montante = question\\ descendante = accord} }
  child[concept color=popmuted, text=white]{ node[align=center]{Pièges\\ des francophones\\ « n'est-ce pas »} };
\end{tikzpicture}
\legende{Carte mentale de la leçon : les question tags}
\end{popfigure}

\begin{bilan}
\textbf{Lexique clé (English / Français)}
\begin{itemize}
\item \eng{question tag} : question courte ajoutée en fin de phrase (« n'est-ce pas »)
\item \eng{statement} : phrase de départ (affirmative ou négative)
\item \eng{auxiliary verb} : auxiliaire (\eng{be, have, do, will, can}...)
\item \eng{rising intonation} : intonation montante (vraie question)
\item \eng{falling intonation} : intonation descendante (on attend l'accord)
\item \eng{to agree / to disagree} : être d'accord / ne pas être d'accord
\end{itemize}

\textbf{Règles à connaître par cœur}
\begin{center}
\begin{tabularx}{\linewidth}{|X|X|}
\hline
\rowcolor{popblueL} Statement & Tag \\ \hline
Affirmative & auxiliaire négatif + pronom : \eng{You are tired, aren't you?} \\ \hline
Négative & auxiliaire affirmatif + pronom : \eng{She isn't here, is she?} \\ \hline
Verbe lexical au present/past simple & \eng{do / does / did} : \eng{He plays, doesn't he?} \\ \hline
\eng{I am...} & \eng{aren't I?} \\ \hline
\eng{Let's...} & \eng{shall we?} \\ \hline
Impératif & \eng{will you? / won't you?} \\ \hline
\eng{nobody, everybody, someone...} & pronom \eng{they} : \eng{Nobody came, did they?} \\ \hline
\eng{nothing, everything...} & pronom \eng{it} : \eng{Nothing happened, did it?} \\ \hline
\eng{there is / there are} & \eng{isn't there? / aren't there?} \\ \hline
\eng{have} (possession, GB) & \eng{have...n't / don't...} ; \eng{have got} $\rightarrow$ \eng{haven't you?} \\ \hline
\end{tabularx}
\end{center}

\textbf{Méthode express en 4 étapes}
\begin{enumerate}
\item Repérer l'auxiliaire de la phrase (s'il n'y en a pas, utiliser \eng{do/does/did}).
\item Inverser la polarité : phrase affirmative $\rightarrow$ tag négatif ; phrase négative $\rightarrow$ tag affirmatif.
\item Remplacer le sujet par le pronom correspondant (\eng{my father} $\rightarrow$ \eng{he}).
\item Répondre selon la réalité des faits : \eng{Yes, I am. / No, I'm not.} — jamais selon la forme de la question.
\end{enumerate}
\end{bilan}

\begin{aretenir}[Checklist avant le Bac]
\begin{itemize}
\item[$\square$] Je sais repérer l'auxiliaire de n'importe quelle phrase anglaise.
\item[$\square$] Je sais que le tag d'une phrase affirmative est négatif, et inversement.
\item[$\square$] Je n'oublie pas que les verbes lexicaux au present simple prennent \eng{do/does} et au past simple \eng{did}.
\item[$\square$] Je connais par cœur les cas spéciaux : \eng{I am $\rightarrow$ aren't I}, \eng{Let's $\rightarrow$ shall we}.
\item[$\square$] Je sais que les impératifs prennent \eng{will you? / won't you?}.
\item[$\square$] J'utilise \eng{they} avec \eng{nobody, everybody, somebody, anybody}.
\item[$\square$] J'utilise \eng{it} avec \eng{nothing, everything, something}.
\item[$\square$] Je ne traduis jamais « n'est-ce pas » mot à mot : le tag dépend de la phrase, pas du français.
\item[$\square$] Je réponds \eng{Yes} ou \eng{No} selon la réalité, pas selon la forme négative de la question.
\item[$\square$] Je sais que l'intonation montante pose une vraie question et l'intonation descendante cherche l'accord.
\item[$\square$] Je m'entraîne à produire cinq phrases avec tags sur ma famille et ma vie quotidienne.
\end{itemize}
\end{aretenir}

\findecours

\end{document}
